\documentclass[12pt,leqno]{amsart}
\usepackage{graphicx}
\baselineskip=16pt

\usepackage{showlabels}
\usepackage{indentfirst,csquotes}

% \topmargin= .5cm
% \textheight= 20cm
% \textwidth= 32cc
% \baselineskip=16pt

% \evensidemargin= .9cm
% \oddsidemargin= .9cm
\usepackage[margin=1in]{geometry}

\usepackage{amssymb,amsthm,amsmath}
\usepackage{xcolor,paralist,hyperref,fancyhdr,etoolbox}
\newtheorem{theorem}{Theorem}
\newtheorem{definition}{Definition}
\newtheorem{remark}{Remark}
\newtheorem{example}{Example}
\newtheorem{assumption}{Assumption}
\newtheorem{lemma}{Lemma}
\newtheorem{proposition}{Proposition}
\newtheorem{corollary}{Corollary}
\newtheorem{conjecture}{Conjecture}
%\newtheorem{example}{Example}
% our packages:
\usepackage{enumitem}
\usepackage{soul}
\usepackage{color-edits}
\usepackage{amstext,amsfonts,amssymb}%amsthm
\usepackage{booktabs}
\usepackage{graphicx,color}
\usepackage{bbm}
\usepackage[normalem]{ulem}
\usepackage{algorithm}
\usepackage{algpseudocode}
\usepackage{tikz}
\usepackage{comment}
\usepackage{graphicx}
\usepackage{subcaption}
\usepackage{mathrsfs}
\usepackage{multirow}
% end our packages

% Natbib setup for author-number style
\usepackage{natbib}
 \bibpunct[, ]{(}{)}{,}{a}{}{,}%
 \def\bibfont{\small}%
 \def\bibsep{\smallskipamount}%
 \def\bibhang{24pt}%
 \def\newblock{\ }%
 \def\BIBand{and}%
% end Natbib

% counter set-ups
\newcounter{experiment}
\renewcommand{\theexperiment}{\arabic{experiment}}
% end of counter set-ups

\hypersetup{ colorlinks=true, linkcolor=black, filecolor=black, urlcolor=black }
\def\proof{\noindent {\it Proof. $\, $}}
\def\endproof{\hfill $\Box$ \vskip 5 pt }

\usepackage{lipsum}

\newcommand{\harsha}[1]{\textcolor{red}{#1}}
\newcommand{\zihe}[1]{\textcolor{blue}{#1}}

\newenvironment{ack}{%
  \renewcommand{\abstractname}{Acknowledgement}% Rename Abstract to Acknowledgements
  \begin{abstract}
}{%
  \end{abstract}
}

\begin{document}

\title[SQP for Optimal Transport]{Beyond Brenier: Sequential Quadratic Programming for General-Cost Monge Optimal Transport}
%Sequential Quadratic Programming for the Monge Optimal Transport Problem
\author{Zihe Zhou \and Harsha Honnappa \and Raghu Pasupathy}
\date{March 2026}
\address{315 N Grant St, West Lafayette, IN 47907-2023, United States}
\email{\{zhou408,honnappa,pasupath\}@purdue.edu}
\maketitle
\let\thefootnote\relax
\footnotetext{MSC2020: Primary 90C15, Secondary 93E20, 60H30.} %%%%%%%%%%

\begin{abstract}
  The Monge optimal transport (OT) problem seeks to optimize transportation cost from a source probability measure to a target probability measure. The optimization is over a space of transport maps and the transportation cost is defined by a cost functional of the maps. As recent interest arises in OT, many works focus on tackling the OT problem computationally using finite-dimensional approximations. In this work, we present the infinite-dimensional formulation of the OT problem over a Banach space. We provide explicit expressions of the first and second-order variation of the objective functional, and of the functional constraint with respect to the transport map. We propose a Sequential Quadratic Programming (SQP) framework for this infinite-dimensional problem. We show that subject to reasonable regularity assumptions, our framework satisfies known conditions for local convergence. Moreover, we demonstrate that a merit functional with sufficient constraint penalization serves as an effective step-size monitor within SQP iterations, leading to global convergence towards critical points. To the best of our knowledge, this is the first attempt at a globally convergent SQP operator recursion over infinite-dimensional spaces.
\end{abstract} %%%%%%%%%

\bigskip


\section{Introduction and Perspective}

Optimal transport (OT) provides a mathematical framework for transforming probability measures in a cost-optimal manner. Given probability measures $\mu$ on $\mathcal{X}$ and $\nu$ on $\mathcal{Y}$, and a measurable cost function $c:\mathcal{X}\times\mathcal{Y}\to\mathbb R\cup\{+\infty\}$, the modern formulation due to Kantorovich considers couplings $\pi$ on $\mathcal{X}\times\mathcal{Y}$ with prescribed marginals. The Kantorovich problem is
\begin{equation}
\inf_{\pi\in\Pi(\mu,\nu)} \int_{\mathcal{X}\times\mathcal{Y}} c(x,y)\,d\pi(x,y),
\end{equation}
where $\Pi(\mu,\nu)$ denotes the set of all probability measures on $\mathcal{X}\times\mathcal{Y}$ with marginals $\mu$ and $\nu$. This problem is a linear program over a convex feasible set. Under mild regularity conditions on the cost (e.g., lower semicontinuity and integrability), the feasible set is weakly compact and strong duality holds, yielding a well-understood dual formulation and guaranteed existence of optimal solutions \citep{Villani2003,Villani2009}. These structural properties have enabled the development of reliable and scalable algorithms, including linear and interior-point methods for discrete OT, first-order primal--dual schemes, and, after entropic regularization, highly efficient Sinkhorn-type algorithms \citep{Cuturi2013,PeyreCuturi2019}.

\subsection{Monge OT Often Induces Nonconvexity}
The classical Monge formulation predates Kantorovich and seeks a deterministic transport map $s:\mathcal{X}\to\mathcal{Y}$ satisfying the pushforward constraint $s_\#\mu=\nu$, minimizing
\begin{equation}\label{monge}
\int_{\mathcal{X}} c(x,s(x))\,d\mu(x).
\end{equation}
While conceptually appealing and often required in applications, the Monge problem is inherently nonconvex because the pushforward constraint $s_\#\mu=\nu$ is nonlinear, and the set of feasible maps is not convex. As a result, existence, uniqueness, and algorithmic tractability are substantially more delicate than in the Kantorovich setting.

A notable exception arises in the Brenier setting, where $\mathcal{X}=\mathcal{Y}=\mathbb R^d$, the cost is quadratic, and $\mu$ is absolutely continuous with respect to Lebesgue measure. In this case, the optimal Kantorovich solution is induced by a unique Monge map that can be expressed as the gradient of a convex potential \citep{Brenier1991,McCann1995}. This representation restores convex structure at the level of the potential and underlies a substantial literature on convex optimization and Newton-type solvers for the associated Monge--Amp\`ere equation \citep{FroeseOberman2011,Benamou2014}. However, this convexity is highly fragile and does not persist beyond this specific setting.

Outside the Brenier regime, convexity typically breaks down even when Monge maps exist. For nonquadratic costs, e.g., $c(x,y)=\|x-y\|^p$ with $p\neq 2$ or more general costs satisfying a twist condition, existence and uniqueness of Monge solutions may still hold under suitable regularity assumptions, but the resulting formulations rely on $c$-convex geometry and no longer correspond to convex optimization problems \citep{GangboMcCann1996,Villani2009}. Similarly, when the source measure $\mu$ is singular, as in discrete or empirical measures, Brenier's theorem fails even for quadratic cost. Kantorovich solutions typically split mass, whereas Monge solutions lead to genuinely nonconvex optimization problems.

Despite these difficulties, there are many settings in which a deterministic Monge map is explicitly required. In economics and matching theory, agents must be matched deterministically and cost or surplus functions are rarely quadratic \citep{Galichon2016}. In data-driven applications such as domain adaptation, generative modeling, and simulation, the source measure is often empirical, yet the goal is to learn a deterministic transformation rather than a probabilistic coupling \citep{Arjovsky2017,PeyreCuturi2019}. In control, particle transport, and dynamical systems, one frequently seeks an explicit flow map rather than a relaxed transport plan. In these contexts, the Kantorovich formulation provides a useful convex relaxation, but its solutions are not directly deployable, and common post-processing heuristics (e.g., barycentric projection) only partially address the modeling requirements.

\subsection{Why SQP Methods?} There appears to be a clear methodological gap between the modeling demands that motivate Monge optimal transport and the algorithmic tools currently available for solving the resulting problems outside the convex Brenier regime. While the Kantorovich formulation admits powerful convex relaxations and enjoys global strong duality, these guarantees rely on optimization over couplings and generally fail to transfer to the Monge formulation once deterministic maps are imposed. In such cases, the restriction to Monge maps introduces genuine nonconvexity, and global strong duality with respect to the Kantorovich dual typically breaks down, reflecting the possibility of mass splitting in the relaxed problem. This failure is structural rather than pathological and is analogous to the loss of strong duality commonly encountered in nonconvex constrained optimization when convexity assumptions are violated \citep{Luenberger1969,Bertsekas1999}.

Importantly, the absence of global strong duality does not preclude the existence of meaningful dual structure at the level relevant for algorithmic design. In nonlinear programming, it is well understood that even for nonconvex problems, local saddle-point structure and strong duality in a \emph{local} sense persist under appropriate constraint qualifications and second-order sufficient conditions \citep{Luenberger1969,NocedalWright2006,Bertsekas1999}. When such conditions hold, Lagrange multipliers exist, the Karush--Kuhn--Tucker (KKT) conditions characterize local optima, and the associated Lagrangian provides a valid local dual model. This notion of local duality, rather than global convex duality, is precisely the foundation on which sequential quadratic programming (SQP) methods are built.

Viewed through this lens, Monge optimal transport fits naturally within the scope of constrained nonconvex optimization. In an SQP framework, the transport map (or an associated potential or parameterization) is treated as the primary decision variable, the nonlinear pushforward constraint is enforced explicitly and linearized at each iterate, and the resulting subproblems take the form of structured quadratic programs that approximate the Lagrangian to second order while preserving the linearized constraints \citep{NocedalWright2006,Bertsekas1999}. Standard globalization strategies, e.g., merit functions, line search, or trust regions, can then be employed to ensure convergence from remote starting points, even in the absence of global convexity \citep{Luenberger1969}. While global strong duality may fail for Monge optimal transport, the persistence of local duality and second-order structure provides a rigorous theoretical foundation for SQP-type methods, suggesting a principled pathway for developing algorithms with provable local convergence guarantees in regimes where deterministic transport maps are indispensable and convexity is absent.

A key feature of our framework is that it does not assume the twist condition on the cost function. In one dimension, when the cost satisfies the twist condition, specifically, when the mixed derivative $D^2_{xy}c$ has a definite sign, the structure of the optimal Monge map is determined a priori see \citep[Theorem 3.2]{guillen2010five}. If $D^2_{xy}c<0$, the optimal map is monotone increasing and given by the quantile coupling $s^*(x)=G^{-1}\circ F(x)$; if $D^2_{xy}c>0$, it is monotone decreasing and given by $s^*(x)=G^{-1}\circ(1-F(x))$. 
The regime our framework targets where such structural guarantees are absent. When the mixed derivative $D^2_{xy}c$ does not have a uniform sign, the monotonicity of the optimal map is unknown and neither inverse-CDF formula applies. In this paper, the twist condition is not a standing assumption, but enters only in Section~\ref{sec:cp}, where we invoke classical existence and uniqueness results to ensure that the Monge problem admits a well-defined optimizer $s^*$ against which our computational approximation can be benchmarked. The SQP framework itself requires smoothness of the cost and admissibility of the transport map (monotonicity and nondegeneracy, as formalized through the problem domain $\Omega$), but does not require the twist condition or any a priori knowledge of which monotonicity class the optimizer belongs to.

\subsection{Summary of Insight and Contribution}
Our goal is to study and solve the general cost Monge OT problem with with (i) a theoretical analysis posed directly in the infinite-dimensional space of transport maps, and
(ii) a computational scheme that is implementable and whose objective values provably approach the true objective as the approximation and discretization is refined. 

We assume that $\mathcal{X}, \mathcal{Y} \subset \mathbb R$. Working directly over the space of transport maps, we impose measure preservation through the pushforward constraint:

% \begin{align*}
%     &\min_{s\in\mathcal{S}} J(s)=\int_{\mathcal{X}} c(x, s(x))\,f(x)dx\\
%     &\text{s.t.}\qquad Pf(s)\equiv g,
% \end{align*}
\begin{align*}
    &\min_{s\in\Omega} J(s)=\int_{\mathcal{X}} c(x,s(x))\,d\mu(x)\\
    &\text{s.t.}\quad s_\#\mu=\nu,
\end{align*}
Throughout, we assume $\mu$ and $\nu$ are absolutely continuous with respect to Lebesgue measure on $\mathcal{X}$,
so we may write $d\mu(x)=f(x)dx$ and $d\nu(y)=g(y)\,dy$ for densities $f$ and $g$.
Under absolute continuity, the pushforward constraint $s_{\#}\mu=\nu$ can be expressed equivalently at the level of
densities. To this end, we introduce the Perron--Frobenius (pushforward) operator $Pf(s)$, which maps an input density
$d\mu=f$ to the density of the pushforward measure $s_{\#}\mu$. With this notation,
\[
s_{\#}\mu=\nu
\quad \Longleftrightarrow \quad
Pf(s) = g \ \ \text{(a.e. on } \mathcal{X}\text{)}.
\]
Hence, the measure-preserving constraint can be written as the infinite-dimensional equality constraint
$Pf(s) - g = 0$.
Most classical results that characterize the solution of Monge OT require ``nice'' cost structures, most notably transportation costs of the form $c(x,y)=\tilde c(x-y)$ with $\tilde c$ quadratic or strictly convex
(e.g., \cite{brenier1991polar,GangboMcCann1996}).
In contrast, we do not require strict convexity of the cost on the translation $|x-y|$, nor do we require $c(x, y)$ to take the form of $c(x,y)=\tilde c(x-y)$.
Instead, we work under milder conditions on $c$, such as smoothness and injectivity properties, which allow us to treat a broader class of costs while remaining in the Monge setting. The key contributions of this paper are as follows:

%A key observation in this work is on the local geometry of the feasible set: depending on the choice of the approximation space for $\mathcal{S}$, feasible points can be isolated, and feasibility cannot be preserved by small local perturbation. These two features motivate an approach that can move through infeasible region and handle nonconvexity. We therefore adopt a sequential quadratic programming (SQP) approach, which solves a sequence of sub-problems at each iteration of the method, and gradually approaches the true solution.

\begin{enumerate}[leftmargin=*, label=\arabic*.]

\item (Local Geometry of the Feasible Sets) 
A central structural fact in our analysis is that the first-variation operator of the pushforward constraint is injective, and that this injectivity has different geometric implications in finite- versus infinite-dimensional spaces.
When the solution space is finite-dimensional, injectivity of the firs-variation operator $Pf_s$ of $Pf(\cdot)$ implies that feasible points are isolated. Intuitively, if the Jacobian has full column rank, then the tangent
space is trivial:
\[
Pf_s(h)=0 \quad \Longrightarrow \quad h=0, \quad \mu\text{--a.e}
\]
so there is no nonzero first-order direction that preserves the constraint. Geometrically, one can think of the feasible set as a collection of isolated points in parameter space and moving in any direction immediately breaks feasibility.
In an infinite-dimensional function space, however, the same conclusion no longer follows from injectivity alone.
The finite-dimensional isolation argument relies on a key fact of the injective operator, namely a positive minimum modulus:
\[
\inf_{\|h\|=1}\|Pf_s(h)\| \;>\; 0,
\]
which provides a uniform lower bound and forces the null space to be robustly trivial. In infinite dimensions, injectivity does not imply a positive minimum modulus as one can construct a sequence
$\{h_n\}\subset \mathcal{S}$ with $\|h_n\|=1$ such that
\[
\|Pf_s(h_n)\| \to 0 \quad \text{as } n\to\infty,
\]
so the linearized constraint can be made arbitrarily small along nontrivial directions. As a consequence, the sufficient condition behind the finite-dimensional isolation result breaks down, and feasible points are not necessarily isolated in the infinite-dimensional setting.

\item (Consistency of the Computational Objective)
The original optimization problem is posed directly over transport maps and includes a pushforward constraint. While this formulation is conceptually well-posed, it is computationally intractable in its native form: the decision variable lives in an infinite-dimensional function space, and the pushforward constraint is a functional equality constraint. To make the problem solvable, we introduce a finite-dimensional computational objective that approximates the original one, and we prove that this approximation becomes exact in the limit as the approximation and discretization are refined. 

We start by recording standard existence and uniqueness results for the continuous Monge problem. Since our consistency analysis compares the computational approximation to the true objective, we first ensure that the continuous problem admits a unique optimizer $s^*$ (unique $\mu$--a.e.). Two complementary regimes are relevant in our setting. For strictly convex transportation costs of the form
$c(x,y)=\tilde c(x-y)$, strict convexity of $\tilde c$ yields uniqueness of the optimal transport map and an almost-everywhere inverse map via \cite{GangboMcCann1996}. To cover objectives beyond strict convexity and beyond the transportation form, we also rely on the existence and uniqueness theorem of \cite{philippis2013regularity}. In that regime we do not assume convexity; instead we impose regularity and injectivity conditions on $c$, which are
sufficient to guarantee a unique optimal transport map. Under additional non-degeneracy and H\"older regularity of the cost, the optimal map is an almost-everywhere $C^{1,\alpha}_{\mathrm{loc}}$ bijection between full-measure subsets of $\mathrm{supp}(f)$ and $\mathrm{supp}(g)$. This bijection property is a key input for discretization: it legitimizes a pointwise change of variables
representation of the pushforward constraint involving $s^{-1}$ and its derivative, which can then be enforced on a grid over the support of $g$.

We bridge the theoretical and the computational objectives by introducing two intermediate objectives. First, we restrict the transport maps to a finite-dimensional family $\mathcal{S}_{n\infty}$ (e.g., $n$ degrees piecewise linear maps), isolating the approximation error from representing an infinite-dimensional map by finitely many parameters. Second, we replace the functional pushforward constraint
by finitely many pointwise constraints enforced on a grid, isolating the approximation error from enforcing the constraint only at a
discrete set of locations. These two approximations naturally define intermediate problems that isolate each source of error, and the
final computational problem combines both.

The convergence proof relies on the error decomposition brought by the two intermediate problems. We start by only enforcing the functional
pushforward constraint with a finite set of pointwise equality constraints over a grid, and then additionally restrict the transport map to a finite-dimensional subspace. Denote the corresponding optimal values by
$J(s_{\infty m}^*)$ (finite-dimensional constraint) and $J(s_{nm}^*)$ (finite-dimensional solution space). Two ingredients drive the argument. First, under our regularity assumptions, $Pf(s)$ admits a continuous representative; therefore,
if $Pf(s)=g$ holds on a refining sequence of grids whose union is dense in $\mathrm{supp}(g)$, then $Pf(s)=g$ holds everywhere. This yields
the pointwise limit
\[
J(s_{\infty m}^*)\;\longrightarrow\;J(s^*)\qquad (\text{as }m\to\infty),
\]
i.e., the constraint gridding-only optimal value converges to the true optimal value. Second, Lipschitz continuity of $J$ with respect to $\|\cdot\|_{\infty}$ implies that a uniform approximation error in the map yields a proportional error in the objective value; specifically,
\[
|J(s)-J(\hat s)| \le K\,\|s-\hat s\|_{\infty} .
\] Using this and monotonicity of feasible sets in $n$,
we prove that for each fixed $m$,
\[
J(s_{nm}^*)\;\longrightarrow\;J(s_{\infty m}^*)\qquad (\text{as }n\to\infty),
\]
and moreover this convergence is \emph{uniform in $m$}. With these two facts in hand, the Moore--Osgood theorem allows us to exchange limits and conclude
\[
\lim_{n\to\infty}\lim_{m\to\infty} J(s_{nm}^*)
\;=\;
\lim_{m\to\infty}\lim_{n\to\infty} J(s_{nm}^*)
\;=\;
J(s^*),
\]
i.e., the optimal objective value of the computational problem converges to the optimal value of the original infinite-dimensional problem
as $n,m\to\infty$.


Finally, we quantify the approximation with explicit rates. For fixed grid size $m$, we show
\[
\big|J(s_{nm}^*)-J(s_{\infty m}^*)\big|\;=\;O\!\left(\frac{1}{n}\right),
\]
by constructing a feasible piecewise linear approximation of $s_{\infty m}^*$ with sup-norm error $O(1/n)$ and invoking the
Lipschitz continuity of $J$. For the pushforward constraint approximation, we show
\[
\big|J(s_{\infty m}^*)-J(s^*)\big|\;=\;O\!\left(\frac{1}{m}\right),
\]
by refining the grid and making localized corrections that enforce each newly added constraint (as we increase the number of grid points) while preserving the earlier ones. Each correction is controlled in $\|\cdot\|_{\infty}$, and Lipschitz continuity of $J$ makes the resulting objective changes summable, yielding the
$O(1/m)$ rate.

\item (SQP Framework for General Cost Monge OT)
We develop a SQP framework for the general cost Monge OT, and we establish convergence guarantees in both the functional and implementable finite-dimensional settings. 
In the infinite-dimensional setting, we formulate a functional SQP step. Using the first- and second-order Fr\'echet derivative formulas, we verify the local convergence conditions of \cite{alt1990lagrange} at a local solution, which yields a well-defined SQP step and a local convergence result for the resulting iterates. To obtain globalization, we introduce a merit functional that balances objective decrease and constraint violation, and we show that for a sufficiently large penalty parameter, each SQP step produces a descent direction for the merit functional. Coupled with an Armijo line search, this yields global convergence of the SQP iterates to a critical point of the merit functional from an arbitrary initialization, this is to be differentiated from convergence to the global optimizer.
In the finite-dimensional setting, we derive the discrete SQP iteration for the computational problem by writing the quadratic subproblem and its KKT linear system. We then characterize standard regularity assumptions under which the discrete SQP method can converge locally with appropriate grid choice for the constraint: when the constraint Jacobian satisfies LICQ and a second-order sufficient condition holds at a solution pair, the iterates converge quadratically with the exact Lagrangian Hessian, following the classical SQP theory in \cite{wright1999numerical}. Together, these results provide an optimization framework for nonconvex Monge OT that is grounded in explicit variational calculus and comes with both theoretical and computational convergence guarantees.

\item (First and Second Variations) 
At each SQP iteration, we solve a local quadratic program obtained by (i) a second-order approximation of the objective and (ii) a first-order approximation of the pushforward constraint. Concretely, given the current iterate $s\in\mathcal{S}$ and an
increment direction $h\in\mathcal{S}$, SQP uses the local expansions:
\[
J(s+h)\;\approx\; J(s) + J_s(h) + \tfrac12\,H_s(h,h),
\qquad
Pf(s+h)\;\approx\; Pf(s) + Pf_s(h),
\]
where $J_s\in\mathcal{S}^*$ and $H_s\in\mathcal{S}^{**}$ are the first and second Fr\'echet derivatives of $J$, and $Pf_s$ is the Fr\'echet derivative of the pushforward operator. This leads to a Newton-Lagrange system (equivalently, the KKT conditions of the quadratic subproblem) that couples the quadratic objective with linearized constraint. As a foundation for this construction, we derive explicit first- and second-order variations (which are Fr\'echet) for both the Monge objective functional and the pushforward constraint operator. These formulas serve as the basic components for the SQP subproblems that support the convergence analysis, and also entail the local geometry of feasibility, as described above.
\end{enumerate}


\section{Example Settings}

In this section, we outline several example applications where the Monge OT in~\eqref{monge} becomes the natural formulation. As noted, nonconvexities in~\eqref{monge} can be caused either due to the cost $c(\cdot,\cdot)$ not being quadratic, or due to the source measure $\mu$ being singular. The examples below are chosen to exemplify these different sources of nonconvexity. 

\subsection{Deterministic Matching}
\label{sec:econ_monge_example}

As a first example, consider a classical matching problem arising in economics, where the objective is to deterministically assign agents from one population to agents in another so as to maximize aggregate surplus. Let worker types be represented by $x \in \mathcal{X} \subset \mathbb R^d$, distributed according to a probability measure $\mu$, and job (or firm) types by $y \in \mathcal{Y} \subset \mathbb R^d$, distributed according to $\nu$. A feasible assignment is given by a measurable map
$s : \mathcal{X} \to \mathcal{Y}, s_\# \mu = \nu,$
so that each worker type is matched to exactly one job type and the induced distribution of jobs matches $\nu$.

A widely used specification for the match surplus (or equivalently, negative cost) takes the form
\begin{equation}
\label{eq:econ_cost}
c(x,y)
=
\phi(x)
+
\psi(y)
-
x^\top A y
+
r(x,y),
\end{equation}
where each term has a clear economic interpretation: $\phi:\mathcal{X} \to \mathbb R$ represents worker-specific productivity or outside options that are independent of the match; $\psi:\mathcal{Y} \to \mathbb R$ represents job-specific productivity or reservation utilities; $A \in \mathbb R^{d\times d}$ encodes bilinear complementarities between worker and job characteristics; and $r(x,y)$ captures nonlinear interaction effects, frictions, or institutional features that are not well represented by a purely bilinear model. (Since the term $\phi(x)$ does not depend on $s$, and since $\int_{\mathcal{X}} \psi(s(x))\, d\mu(x) = \int_{\mathcal{Y}} \psi(y)d\nu(y)$, the terms $\phi$ and $\psi$ in~\eqref{eq:econ_cost} do not affect the optimizer, but they are included for modeling completeness and economic interpretation.)

In many such economic settings, couplings that split mass are inadmissible, as they correspond to random or mixed assignments of worker types to job types. Consequently, the Monge formulation in~\eqref{monge} is the more meaningful model, whereas the Kantorovich formulation serves only as a relaxation. Under suitable regularity conditions on $\mu$ and $\nu$ (e.g., absolute continuity) and mild assumptions on the cost, such as a \emph{twist condition} $y \mapsto \nabla_x c(x,y)$ being injective for each $x,$ it is known that a unique optimal Monge map exists for problem \eqref{monge}. The bilinear term $-x^\top A y$ typically enforces this condition, and small or smooth nonlinear perturbations $r(x,y)$ preserve it. However, the problem (while well-defined) becomes nonconvex if $r(x,y)$ is nonquadratic, and the classical Brenier representation of the optimal map as the gradient of a convex function fails.

\subsection{Transportation and Routing}
\label{sec:transport_monge_example}

We next consider a class of transportation and routing problems in which individual agents, vehicles, or flow units must be assigned deterministically to routes, schedules, or trajectories. Let $x \in \mathcal{X} \subset \mathbb R^d$ denote an agent descriptor such as an origin-destination pair, departure time, or state vector, distributed according to a probability measure $\mu$. Let $y \in \mathcal{Y} \subset \mathbb R^m$ represent a transportation outcome, such as a route choice, arrival time, or trajectory summary, with target distribution $\nu$. A feasible transportation policy is given by a measurable map $s : \mathcal{X} \to \mathcal{Y}, s_\# \mu = \nu,$ so that each agent is assigned a unique transportation outcome and the aggregate distribution of outcomes matches prescribed system-level requirements.

A representative cost functional takes the form
\begin{equation}
\label{eq:transport_cost}
c(x,y)
=
\ell(x,y)
+
g\bigl(y,\,R(s)\bigr),
\end{equation}
where $\ell(x,y)$ represents an individual travel cost (e.g., travel time, distance, or delay), and $g(y,R(s))$ captures congestion effects, capacity penalties, or other system-level interactions that depend on the aggregate flow induced by the assignment $s$. Here, $J(s)$ denotes a suitable flow or load functional determined by the pushforward measure $s_\#\mu$, that is, $R(s) \equiv J(s_\#\mu)$ and $R$ is often a linear functional.

As in the first example, since randomized or mass-splitting assignments are not operationally meaningful here, the Monge formulation in~\eqref{monge} provides the appropriate modeling framework. Under suitable regularity conditions, e.g., smooth $\ell$ and monotonicity properties of $g$, optimal assignments may exist and be unique, particularly when the induced marginal $\nu$ lies in the interior of the feasible flow polytope. However, the problem is often nonconvex because the presence of congestion or capacity effects introduces nonlinear dependence of the cost on the induced pushforward measure. Even when $\ell(x,y)$ is convex in $y$, the aggregate interaction term $g(y,R(s))$ typically destroys any underlying convex structure.

\subsection{Deterministic Map Learning}
\label{sec:ml_monge_example}

We finally consider a class of problems arising in machine learning and data-driven modeling, where the objective is to learn a deterministic transformation between probability distributions. Let $\mu$ denote a source distribution supported on $\mathcal{X} \subset \mathbb R^d$, and let $\nu$ denote a target distribution supported on $\mathcal{Y} \subset \mathbb R^m$. In many applications, $\mu$ is an empirical measure of the form 
$\mu = \frac{1}{n}\sum_{i=1}^n \delta_{x_i},$
corresponding to a finite dataset, while $\nu$ may be continuous or empirical. A feasible model is given by a deterministic map $s : \mathcal{X} \to \mathcal{Y}, s_\# \mu \approx \nu,$
which transforms each data point $x_i$ into an output $s(x_i)$ so that the induced distribution matches the target.

In contrast to the Kantorovich formulation, which permits mass splitting and randomized couplings, many machine learning applications require an explicit deterministic transformation. Examples include normalizing flows, generative models, and domain adaptation, where a map $s$ must be evaluated pointwise to generate samples, transfer representations, or deploy a learned model. In such settings, transport plans are not directly implementable, and the Monge formulation in~\eqref{monge} provides the appropriate modeling framework.

A representative objective takes the form
\begin{equation}
\label{eq:ml_cost}
\min_{s \in \mathcal{S}}
\left\{
\int_{\mathcal{X}} c\bigl(x,s(x)\bigr)\,d\mu(x)
\;+\;
\lambda\,\mathsf D\!\left(s_\#\mu,\,\nu\right)
\right\},
\end{equation}
where $c(x,y)$ is a pointwise loss function, $\mathsf D(\cdot,\cdot)$ is a discrepancy or divergence between probability measures (e.g., an integral probability metric or moment-matching penalty), and $\mathcal{S}$ denotes a parameterized class of maps, such as neural networks or other function approximators.

Even when the loss $c(x,y)$ is convex in $y$ and the divergence $\mathsf D$ is convex in its arguments, problem~\eqref{eq:ml_cost} is generally nonconvex. First, the source measure $\mu$ is singular, so the classical Brenier theory does not apply even for quadratic cost, and the optimal Kantorovich solution typically splits mass. Imposing a deterministic map therefore introduces a nonconvex restriction on the feasible set. Second, the parameterization of $s$ through a finite-dimensional family $\mathcal{S}$ leads to additional nonconvexity in the decision variables. As a result, the optimization problem lies outside the convex regimes traditionally associated with optimal transport, despite being motivated by distributional matching objectives.



\subsection{Physics-Inspired Transport Costs}
\label{sec:physics_monge_example}

Our final class of examples arises from physical systems in which the transport cost $c(x,y)$ is determined by an underlying physical law rather than prescribed by modeling convenience. We present three example settings, in optics, classical mechanics, and quantum mechanics, and, in each case, the cost takes the problem outside the quadratic Brenier regime, and the Monge formulation is the physically meaningful one. Specifically, the optimizer seeks a deterministic assignment of source configurations to target configurations, not a probabilistic coupling.

\paragraph{\it Optical path costs.}
In freeform illumination design~\citep{yadav2019monge,romijn2019freeform,prins2015least}, a lens or reflector must redirect a source intensity $f(x)$ into a prescribed target intensity $g(y)$. Each ray entering at position $x$ and arriving at $y$ incurs a cost equal to its optical path length. For a cylindrical lens with screen distance $D$, this yields
\begin{equation}\label{eq:lens-cost-sec2}
    c_{\textup{lens}}(x,y) = \sqrt{(x-y)^2 + D^2},
\end{equation}
which interpolates between the quadratic cost ($|x-y| \ll D$) and the $L^1$ cost ($|x-y| \gg D$). The far-field reflector antenna problem~\citep{wang2004design,glimm2003optical} gives rise to the logarithmic cost
\begin{equation}\label{eq:reflector-cost-sec2}
    c_{\textup{refl}}(x,y) = -\log|x - y|,
\end{equation}
posed naturally on the sphere $S^{n-1}$, where source and target occupy disjoint hemispheres and the cost is bounded. In both cases, the Monge formulation is the physically correct one: each ray must be assigned a unique destination, and mass-splitting (i.e., a single ray arriving at multiple locations) is not physically realizable.

Both costs satisfy the twist condition ($D^2_{xy}c < 0$ away from the diagonal), so the Gangbo--McCann theory~\citep{GangboMcCann1996} applies and optimal maps exist and are unique. However, the costs are genuinely non-quadratic, and the optimal map is not the gradient of a convex function in any simple sense. The lens cost has a twist strength $|D^2_{xy}c_{\textup{lens}}|$ that decays as $|x-y|^{-3}$ for large displacement, while the reflector cost has a singularity at $x = y$. The regularity theory for the reflector cost is well-developed: Wang~\citep{wang2004design} established existence via a Monge--Amp\`ere formulation, and Loeper~\citep{loeper2011regularity} showed that the Ma--Trudinger--Wang condition is satisfied on the sphere, yielding smooth optimal maps for smooth data~\citep{ma2005regularity,loeper2009regularity}.

\paragraph{\it Lagrangian obstacle costs.}
A more challenging setting arises when the cost is not given in closed form but is itself defined through a variational problem. Consider transporting mass from $x$ to $y$ along a path $\gamma:[0,1] \to \mathbb{R}$ through a potential landscape $V(z)$. The cost is the minimum-action path:
\begin{equation}\label{eq:lagrangian-cost-sec2}
    c(x,y) = \inf_{\substack{\gamma \in C^1([0,1]) \\ \gamma(0)=x,\, \gamma(1)=y}}
    \int_0^1 \bigl[\lambda\,|\dot\gamma(t)|^2 + V(\gamma(t))\bigr]\, dt,
\end{equation}
where $\lambda > 0$ controls the kinetic energy penalty and $V \geq 0$ is a potential barrier. This is the action functional of classical mechanics; the optimal path satisfies the Euler--Lagrange equation, and $c(x,y)$ is the corresponding Hamilton--Jacobi value function. The Monge problem is then a \emph{bilevel optimization}: the outer problem chooses the transport map $s$, while the inner problem~\eqref{eq:lagrangian-cost-sec2} computes the cost of each assignment.

When $V = 0$, the inner problem yields straight-line paths and $c(x,y) = \lambda(x-y)^2$, recovering the quadratic cost. A nonzero barrier $V$ deforms the cost in two ways that take it outside the Gangbo--McCann framework. First, the cost function is a value function of a variational problem; by the theory of Hamilton--Jacobi equations, such value functions are Lipschitz but not necessarily $C^1$ at points where multiple optimal paths exist, violating the smoothness hypotheses of the twist condition. Second, even where $c$ is smooth, the twist condition can fail locally: in regions where the barrier forces all paths through a common bottleneck, the mapping $y \mapsto D_x c(x,y)$ loses injectivity. Since one cannot verify a priori whether the optimal transport plan avoids the twist-failure region, the standard sufficient conditions for existence and uniqueness are not applicable.

This class of problems has been studied in the recent machine learning literature under the name ``optimal transport with Lagrangian costs''~\citep{pooladian2024neural,asadulaev2024neural,fan2021scalable}, where both the transport map and the inner paths are parameterized by neural networks. As in the preceding examples, the Monge formulation is the natural one: the physical system requires a deterministic assignment of source positions to destinations, not a probabilistic coupling.

\paragraph{\it Coulomb repulsive costs.}
The most extreme departure from the Brenier regime arises in the strong-interaction limit of Density Functional Theory (DFT), where the ground-state energy of $N$ electrons with one-body density $\rho$ is determined by a multi-marginal optimal transport problem with Coulomb cost~\citep{Seidl1999,ButCDP2012}. For $N = 2$ electrons in one dimension, this reduces to finding a map $f$ pushing $\rho/2$ to itself that minimizes the Coulomb repulsion:
\begin{equation}\label{eq:sce-cost-sec2}
    \min_{f:\; f_\#(\rho/2) = \rho/2} \int \frac{1}{|x - f(x)|}\,\frac{\rho(x)}{2}\,dx.
\end{equation}
The optimal map $f$ is the \emph{co-motion function}~\citep{Seidl1999}: $f(x) = F^{-1}(1 - F(x))$, where $F$ is the CDF of $\rho$. This is the anti-quantile coupling, a monotone \emph{decreasing} map that pairs each electron with its ``opposite'' in the density distribution.

The Coulomb cost $c(x,y) = 1/|x-y|$ violates the hypotheses of the Gangbo--McCann theory in essentially every way. The cost is singular on the diagonal ($c \to +\infty$ as $x \to y$); it is convex and decreasing in $|x-y|$ (repulsive), the opposite regime from the attractive costs for which classical theory is developed; the mixed partial $D^2_{xy}c = 2/|x-y|^3 > 0$ has the wrong sign for the standard twist condition; and the self-transport constraint $f_\# \mu = \mu$ is degenerate for classical uniqueness results. Nonetheless, the problem is well-posed for bimodal densities with sufficient separation: the co-motion function avoids the singularity by mapping each atomic mode to the opposite mode, and the mathematical structure has been analyzed by Colombo, De Pascale, and Di Marino~\citep{ColomboDeP2015}.

The Monge formulation is again the physically correct one. In the SCE framework, the co-motion function encodes the deterministic relationship between electron positions in the strongly correlated ground state; couplings that split mass correspond to partially correlated states that do not arise in the strong-interaction limit. Beyond the two-electron case, the $N$-electron problem gives rise to a multi-marginal optimal transport problem with pairwise Coulomb cost $\sum_{i < j} 1/|x_i - x_j|$, where computational cost scales exponentially in $N$ and the existence of optimal maps in general three-dimensional settings remains open~\citep{ButCDP2012}.

\paragraph{\it Common structure.}
The three settings above share a common structure: the cost is determined by a physical law, the Monge formulation is the natural one, and the problem lies outside the regime where Brenier's theorem or the Gangbo--McCann theory provides a complete answer. For the optical costs, the twist condition holds but the cost is non-quadratic, so the optimal map is not the gradient of a convex function. For the Lagrangian cost, the twist condition fails locally, so existence and uniqueness cannot be guaranteed by standard theory. For the Coulomb cost, the twist has the wrong sign, the cost is singular, and the transport is self-referential. In each case, a direct optimization approach to the Monge problem --- one that does not rely on structural properties of the cost such as the twist condition or $c$-convexity --- provides the appropriate computational framework.


\section{Mathematical Preliminaries}\label{sec:pre}
\subsection{Notation}\label{subsec:notation}
Let $\mathcal{X}\subseteq\mathbb R$ be compact. Let $\mu$ and $\nu$ be two probability measures on $\mathcal{X}$, and assume they are absolutely continuous
with respect to Lebesgue measure with density functions $f,g\in \mathcal{P}(\mathcal{X})$, respectively; that is,
\[
d\mu(x)=f(x)dx, \qquad d\nu(x)=g(x)dx.
\]

Let $\mathcal{Z} := L^1(\mathcal{X})$ and define the mean-zero subspace $\mathcal{Z}_0:= \left\{w \in L^1(\mathcal{X}):\int_{\mathcal{X}} w(x)\,dx = 0\right\}$. 
Since $\mathcal{Z} = L^1(\mathcal{X})$, its Banach dual is $\mathcal{Z}^* := \mathcal{B}(\mathcal{Z},\mathbb R)$, the space of bounded linear functionals on $\mathcal{Z}$.
Moreover, by the standard $L^1$--$L^\infty$ duality, $\mathcal{Z}^*$ is isometrically isomorphic to $L^\infty(\mathcal{X})$ (denoted by $\mathcal{Z}^* \cong L^\infty(\mathcal{X})$),
in the sense that each $\lambda\in L^\infty(\mathcal{X})$ induces a functional $z^* \in \mathcal{Z}^*$ via the pairing
\[
z^*(\eta) := \int_{\mathcal{X}} \lambda(x)\,\eta(x)dx, \qquad \forall\,\eta \in \mathcal{Z},
\]
and every $z^* \in \mathcal{Z}^*$ admits such a representation for some $\lambda\in L^\infty(\mathcal{X})$, unique up to equality almost everywhere~\citep[Theorem~6.15]{folland1999real}.
With this identification, we define the dual pairing
\[
\langle z^*,\eta\rangle_{\mathcal{Z}^*,\mathcal{Z}} := z^*(\eta) = \int_{\mathcal{X}} \lambda(x)\,\eta(x)dx, \qquad z^*\in \mathcal{Z}^*,\ \eta\in \mathcal{Z}.
\]
with the associated operator norm
\[
\|z^*\|_{\mathcal{Z}^*} := \sup\Big\{ \langle z^*,\eta\rangle_{\mathcal{Z}^*,\mathcal{Z}} : \|\eta\|_{\mathcal{Z}}\le 1\Big\} = \|\lambda\|_{L^\infty(\mathcal{X})}.
\]

The dual space of $\mathcal{Z}_0$ is $\mathcal{Z}_0^* \cong L^\infty(\mathcal{X})/\mathbb{R}$.
This quotient structure arises because every $w \in \mathcal{Z}_0$ integrates to zero, so adding a constant to a representative $\lambda \in L^\infty(\mathcal{X})$ does not change the induced functional on $\mathcal{Z}_0$:
\[
\int_{\mathcal{X}} \big(\lambda(x)+c\big)w(x)\,dx
=
\int_{\mathcal{X}} \lambda(x)w(x)\,dx.
\]
Accordingly, elements of $\mathcal{Z}_0^*$ are represented by equivalence classes $[\lambda] \in L^\infty(\mathcal{X})/\mathbb{R}$.

We consider the space of transport maps to be a subspace $\mathcal{S}\subseteq L^\infty(\mathcal{X})$ to ensure they map the compact set $\mathcal{X}$ into itself. Because our analysis relies on inner products for variational calculus and uniform bounds for approximations, we utilize a two-norm framework. Throughout this paper, the default norm $\| \cdot \|$ on any element $h\in\mathcal{S}$ denotes the $L^2(\mu)$ norm:
\[ 
\|h\| := \left(\int_{\mathcal{X}}|h(x)|^2 f(x)dx\right)^{1/2}, 
\]
where $|\cdot|$ denotes the absolute value on $\mathbb{R}$. The associated inner product is
\[
\langle a,b\rangle_{L^2(\mu)} := \int_{\mathcal{X}} a(x)b(x)f(x)dx. 
\]
When uniform bounds are required (such as for establishing the Lipschitz continuity of the objective functional), we explicitly invoke the supremum norm, denoted by $\| \cdot \|_\infty$.

We define the mean-zero subspace of $L^2(\mu)$ as $L^2_0(\mu) := \left\{ \psi \in L^2(\mu) : \mathbb{E}[\psi(Y)] = 0 \right\}$, and the mean-zero elements of the transport map space as $\mathcal{S}_0 := \{h \in \mathcal{S} : \mathbb{E}[h(Y)] = 0\}$.

For any vector $v\in\mathbb{R}^d$, the same notation $\|v\|$ denotes the Euclidean norm,
\[
\|v\| := \|v\|_2 := \Big(\sum_{i=1}^d v_i^2\Big)^{1/2},
\]
and $\langle u,v\rangle := u^\top v$ denotes the standard Euclidean inner product.
For densities (elements of $\mathcal{Z}$) we define the $L^1$ norm with respect to Lebesgue measure.
\[
\|z\| _{\mathcal{Z}}:=\int_{\mathcal{X}}|z(x)|dx,\qquad z\in \mathcal{Z}.
\]
Denote by $\mathcal{B}(U,V)$ the space of bounded linear maps from $U$ to $V$.

For $L\in \mathcal{S}^*:=\mathcal{B}(\mathcal{S},\mathbb R)$, define
\[
\|L\|_{\mathcal{S}^*} :=\sup_{\|h\|=1}|L(h)|.
\]
For $I\in \mathcal{S}^{**}:=\mathcal{B}(\mathcal{S},\mathcal{S}^*)$, define
\[
\|I\|_{\mathcal{S}^{**}} :=\sup_{\|h\|=1}\|I(h)\|_{\mathcal{S}^*}.
\]

For $A\in\mathcal{B}(\mathcal{S},\mathcal{Z})$, we define the induced operator norm by
\[
\|A\|_{\mathcal{B}(\mathcal{S},\mathcal{Z})} := \sup_{\|h\|=1}\|A(h)\|_{\mathcal{Z}}.
\]
Likewise, for $A\in\mathcal{B}(\mathcal{S},\mathcal{Z}_0)$, we define
\[
\|A\|_{\mathcal{B}(\mathcal{S},\mathcal{Z}_0)} := \sup_{\|h\|=1}\|A(h)\|_{\mathcal{Z}_0}.
\]

Let \(C^{k,\alpha}_{\mathrm{loc}}\) denote the space of functions whose derivatives up to order \(k\) exist and whose \(k\)th derivatives are locally \(\alpha\)-H\"older continuous. We write \(W^{1,\infty}(\mathcal X)\) for the Sobolev space of essentially bounded functions with essentially bounded weak derivatives, equipped with the norm
\[
\|s\|_{W^{1,\infty}}
:=
\|s\|_{L^\infty}+\|s'\|_{L^\infty}.
\]
Since \(\mathcal X\) is compact, one has the continuous embeddings
\[
W^{1,\infty}(\mathcal X)\hookrightarrow L^\infty(\mathcal X)\hookrightarrow L^2(\mu).
\]
Hence every \(s\in W^{1,\infty}(\mathcal X)\) belongs in particular to \(L^\infty(\mathcal X)\cap L^2(\mu)\). Therefore, any estimate in the sequel that only requires \(L^\infty\)- or \(L^2(\mu)\)-control remains applicable to functions in \(W^{1,\infty}(\mathcal X)\). Finally, \(ACLR(c,d)\) denotes the class of absolutely continuous functions on \((c,d)\) vanishing at both endpoints \(c\) and \(d\).

Lastly, $AC_{LR}(c,d)$ denotes the class of absolutely continuous functions on $(c,d)$ vanishing at both endpoints $c$ and $d$.


\section{Problem Formulation}\label{sec:formulation}
We now introduce the precise Monge optimal transport formulation that will be used throughout the paper.
Let $\mathcal{X}\subseteq \mathbb R$, and let $c(\cdot,\cdot):\mathcal{X}\times\mathcal{X}\rightarrow\mathbb R$ be a cost function. Let $\mathcal{Z} := L^1(\mathcal{X})$ and $\mathcal{P}(\mathcal{X})\subset \mathcal{Z} $ be the set of probability density functions on $\mathbb R$.

Let $\mu$ and $\nu$ be two probability measures on $\mathcal{X}$, and assume they are absolutely continuous
with respect to Lebesgue measure with density functions $f,g\in \mathcal{P}(\mathcal{X})$, respectively; that is,
\[
d\mu(x)=f(x)dx, \qquad d\nu(x)=g(x)dx.
\]
Let $\mathcal{S}\subseteq L^2(\mu)$ be a normed space of transport maps $s:\mathcal{X}\to\mathcal{X}$. Consider the objective functional $J:\mathcal{S}\rightarrow \mathbb R$ defined as 
\[
J(s) :=\int_{\mathcal{X}} c(x, s(x))f(x)dx.
\]
We denote its first and second variation operators by
\[
J':\mathcal{S}\to \mathcal{S}^*, \qquad J'':\mathcal{S}\to \mathcal{S}^{**},
\]
and for $s\in\mathcal{S}$ write
\[
J'(s)=:J_s\in\mathcal{S}^*, \qquad J''(s)=:H_s\in\mathcal{S}^{**}.
\]
We say that $J$ is Fr\'echet differentiable at $s \in\mathcal{S}$ if for all $h\in\mathcal{S}$, 
\[
\big|J(s+h)-J(s)-J_s(h)\big|=o(\|h\|),
\]
and twice Fr\'echet differentiable at $s$ if, in addition,
\[
\|J'(s+h)-J'(s)-H_s(h)\|_{\mathcal{S}^*}=o(\|h\|).
\]

Recall the Perron--Frobenius (PF) operator $Pf:\mathcal{S}\rightarrow \mathcal{P}(\mathcal{X})$ from \cite{boyarsky2012laws},
\[
Pf(s) := \left\{\frac{d}{dx}\int_{s^{-1}(A)} f(y)dy , x\in \mathbb R \right\}
\]
where $A =(-\infty, x]$.  The function $Pf(s)$ is the density of the pushforward measure $s_{\#}\mu$. 
Since $Pf(s)$ is a probability density, we regard $Pf$ as a map into $\mathcal{Z}$ (which includes $\mathcal{P}(\mathcal{X})$), while its first variation takes values in $\mathcal{Z}_0$.
Its first and second variation operators are
\[
Pf':\mathcal{S}\to \mathcal{B}(\mathcal{S},\mathcal{Z}_0), \qquad Pf'':\mathcal{S}\to \mathcal{B}\!\big(\mathcal{S},\mathcal{B}(\mathcal{S},\mathcal{Z}_0)\big),
\]
and for $s\in\mathcal{S}$ we write
\[
Pf'(s)=:Pf_s\in \mathcal{B}(\mathcal{S},\mathcal{Z}_0), \qquad Pf''(s)=:Qf_s\in \mathcal{B}\!\big(\mathcal{S},\mathcal{B}(\mathcal{S},\mathcal{Z}_0)\big).
\]
We say that $Pf$ is Fr\'echet differentiable at $s$ if for all $h\in\mathcal{S}$,
\[
\|Pf(s+h)-Pf(s)-Pf_s(h)\|_{\mathcal{Z}}=o(\|h\|),
\]
and twice Fr\'echet differentiable at $s$ if, in addition,
\[
\|Pf'(s+h)-Pf'(s)-Qf_s(h)\|_{\mathcal{B}(\mathcal{S},\mathcal{Z})} =o(\|h\|).
\]

We call a transport map $s\in\mathcal{S}$ measure-preserving (with respect to $\mu$ and $\nu$) if $s_{\#}\mu=\nu$.
When $d\mu(x)=f(x)dx$ and $d\nu(y)=g(y)\,dy$, this is equivalent to the following identity: for every Lebesgue measurable set $A\subseteq \mathcal{X}$,
\[
\int_{s^{-1}(A)} f(x)dx = \int_A g(y)\,dy,
\]
i.e.,
\[
\mu\big(s^{-1}(A)\big)=\nu(A).
\]
In density form, the same condition can be expressed using the Perron--Frobenius (pushforward) operator:
\[
Pf(s) = g \qquad \text{(Lebesgue a.e. on $\mathcal{X}$)}.
\]
Thus, the statements $s_{\#}\mu=\nu$, the set-wise identity above, and $Pf(s) = g$ are equivalent ways of expressing the same measure-preserving constraint. The pushforward density $Pf(s)$ is well-defined and continuous when $s$ is monotone and nondegenerate  (see Lemma~\ref{lem:pf-continuous-version}). For the local convergence analysis of Section~\ref{subsec:loc_conv}, we further require these structural properties to persist on an entire neighborhood of the minimizer $s_0$, so that $Pf(s)$ and its Fr\'echet derivatives are defined for all $s$ near $s_0$. We therefore introduce the problem domain
\[
\Omega
:=
\bigl\{\, s \in W^{1,\infty}(\mathcal{X}) : s \text{ is monotone and } |s'(x)| \ge c_s > 0 \text{ for some } c_s>0, \text{ a.e.\ } x\in\mathcal{X} \,\bigr\}.
\]
Equivalently, $\Omega = \Omega_+\cup\Omega_-$ where $\Omega_\pm := \{s\in W^{1,\infty}(\mathcal{X}):\pm s'(x)\ge c_s>0 \text{ for some }c_s>0,\text{ a.e.}\}$. Since the $W^{1,\infty}$-norm controls the derivative uniformly, $\Omega$ is open in $W^{1,\infty}(\mathcal{X})$: if $s_0\in \Omega_+$ with $s_0'(x)\ge c>0$ a.e., then for any $\tilde s$ with $\|\tilde s-s_0\|_{W^{1,\infty}}<c/2$, we have $\tilde s'(x)\ge c/2>0$ a.e., so $\tilde s\in \Omega_+$.  This openness guarantees that sufficiently small perturbations of any admissible transport map remain admissible.

All optimization problems in this paper are posed over $\Omega$ or its finite-dimensional restriction $\Omega_n:=\mathcal{S}_n\cap\Omega$ (defined in Section~\ref{subsec:approx}). The role of $\Omega$ is purely one of \emph{admissibility}: it ensures that $Pf(s)$ and its derivatives are well-defined throughout the feasible region. All norms, inner products, and convergence statements continue to use $L^2(\mu)$ and $L^\infty$ as defined above.

We can thus reformulate the measure-preserving constraint into an infinite-dimensional equality constraint and our objective now becomes:

\begin{comment}
    \begin{align*}
    &\min_{s\in\mathcal{S}} J(s)=\int_{\mathcal{X}} c(x, s(x))f(x)dx \tag{P}\\
    &\text{s.t.} \\
    & Pf(s)-g =0_{\mathcal{Z}},
\end{align*}
where $Pf(s)-g$ is thought of as a map from $\mathcal{S}$ to $\mathcal{Z}$. $\mathcal{Z}$ is the range of $Pf(s)-g$. $0_{\mathcal{Z}}$ is the zero element in $\mathcal{Z}$. Note that $\mathcal{Z}$ is a subset of the space of real-valued, integrable functions that integrate to zero.

Let $Z := L^1(\mathbb R)$ and define the closed subspace
\[
Z_0 := \Big\{ \eta \in Z \;:\; \int_{\mathbb R} \eta(x)dx = 0 \Big\}.
\]
Since $Pf(s),g\in \mathcal{P}(\mathcal{X})\subset L^1(\mathbb R)$, we have $Pf(s)-g\in Z_0$ for all $s\in\mathcal{S}$.
Accordingly, define the equality constraint mapping $\mathcal{H}:\mathcal{S}\to Z_0$ by
\[
\mathcal{H}(s) := Pf(s)-g,
\]
and rewrite (P) as $\mathcal{H}(s)=0_{Z_0}$.
\end{comment}


\begin{align*}
    &\min_{s\in\Omega} J(s)=\int_{\mathcal{X}} c(x, s(x))f(x)dx \tag{P}\\
    &\text{s.t.} \\
    & \mathcal{H}(s)=0_{\mathcal{Z}},
\end{align*}

where $0_{\mathcal{Z}}$ is the zero element in $\mathcal{Z}$, and the equality constraint mapping $\mathcal{H}:\mathcal{S}\to \mathcal{Z}$ is defined by
\[
\mathcal{H}(s) := Pf(s)-g \in \mathcal{Z}.
\]
Accordingly, the Lagrangian of (P) is the mapping $\mathcal{L}:\mathcal{S}\times \mathcal{Z}^*\to \mathbb R$ defined by \citep[cf.][{\S}9.3]{luenberger1969optimization}
\[
\mathcal{L}(s,z^*) := J(s) + \langle z^*,\mathcal{H}(s)\rangle_{\mathcal{Z}^*,\mathcal{Z}}, \qquad z^*\in \mathcal{Z}^*.
\]
When $z^*$ is represented by a function $\lambda\in L^\infty(\mathcal{X})$ (i.e., $z^*\equiv \lambda$ under $\mathcal{Z}^*\cong L^\infty(\mathcal{X})$), the Lagrangian of (P) can be written in the integral form
\[
\mathcal{L}(s,\lambda) = J(s) + \int_{\mathcal{X}} \lambda(x)\,\mathcal{H}(s)(x)dx.
\]

\section{Key Assumptions}\label{subsec:assume}
In this section, we list some assumptions. 

\begin{assumption}\label{ass:domain-support}
Let $\mathcal{X}\subset \mathbb R$ be compact. The supports $\mathrm{supp}(f)$ and $\mathrm{supp}(g)$ are nonempty open intervals in $\mathcal{X}$.
\end{assumption}

\begin{assumption}\label{ass:density-cont}
The functions $f$ and $g$ are continuous probability densities on $\mathcal{X}$.
\end{assumption}

\begin{assumption}\label{ass:cost-lipschitz}
The cost function $c(\cdot,\cdot):\mathcal{X} \times\mathcal{X}\rightarrow\mathbb R$ is $K$-Lipschitz in its second argument.
\end{assumption} 

\begin{assumption}\label{ass:cost-second}
The cost function $c(\cdot,\cdot):\mathcal{X}\times\mathcal{X}\to\mathbb R$ is twice continuously differentiable in its second argument.
There exist constants $L,M>0$ such that for all $x\in\mathcal{X}$ and all $y_1,y_2\in\mathcal{X}$,
\[
|D_y c(x,y_1)-D_y c(x,y_2)|\le L\,|y_1-y_2|, \qquad |D^2_{yy}c(x,y_1)-D^2_{yy}c(x,y_2)|\le M\,|y_1-y_2|.
\]
\end{assumption}


\begin{assumption}\label{ass:philippis-C2}
The cost function  $c(\cdot,\cdot):\mathcal{X} \times\mathcal{X}\rightarrow\mathbb R$ is twice continuously differentiable on
$\mathrm{supp}(f)\times \mathrm{supp}(g)$, and
\[
c \in C^2\!\big(\mathrm{supp}(f)\times \mathrm{supp}(g)\big), \qquad \|c\|_{C^2(\mathrm{supp}(f)\times \mathrm{supp}(g))}<\infty,
\]
where, for $X,Y\subseteq \mathcal{X}$,
\[
\|c\|_{C^2(X\times Y)} :=\sup_{(x,y)\in X\times Y}\Bigl( |c(x,y)| +\|\nabla c(x,y)\| +\|D^2 c(x,y)\|_{\mathrm{op}} \Bigr).
\]
Here $\nabla c(x,y)$ denotes the gradient of $c$ with respect to $(x,y)$, and
$D^2 c(x,y)$ denotes the corresponding Hessian matrix. The matrix norm $\|\cdot\|_{\mathrm{op}}$ is the operator norm induced by the Euclidean norm $\|\cdot\|$.
\end{assumption}

\begin{assumption}\label{ass:mono}
$\Omega$ consists of monotone functions $s:\mathrm{supp}(f)\to\mathrm{supp}(g)$ such that $|s'(x)|\ge c_s>0$ for some $c_s>0$, Lebesgue a.e.\ $x\in\mathcal{X}$.
\end{assumption}

\section{Fr\'echet Derivatives of the Objective and Constraint}\label{sec:deriv}
In this section, we derive the first and second Fr\'echet derivatives for both the Monge objective functional $J(s)$ and the Perron--Frobenius pushforward operator $Pf(s)$. These explicit operator representations serve as the computational and theoretical foundation for the local quadratic approximations utilized in the remainder of the paper.
\begin{theorem}\label{thm: J1deriv} 
    Under Assumption~\ref{ass:cost-lipschitz}, the map $J_s(\cdot): \mathcal{S}\rightarrow\mathbb{R}$ defined by
    \begin{align*}
        J_s(h) := \int_{\mathcal{X}}h(x)\, D_y c\big(x,s(x)\big)\, f(x)dx
    \end{align*}
    is the Fr\'echet derivative of $J(\cdot)$ at $s$. 
\end{theorem}

\begin{theorem}\label{thm: J2deriv} 
    Under Assumption~\ref{ass:cost-second}, the bilinear form $H_s(\cdot)(\cdot): \mathcal{S}\times\mathcal{S}\rightarrow\mathbb{R}$ defined by
    \begin{align*}
        H_s(h)(\Tilde{h}) :=\int_{\mathcal{X}}\Tilde{h}(x)\,h(x)\, D^2_{yy}c\big(x,s(x)\big)\, f(x)dx    
    \end{align*} 
    is the Fr\'echet derivative of $J_s(h)$ at $s$. Equivalently, $H_s(\cdot)(\cdot)$ defines an operator
    $H_s(\cdot)(\cdot):\mathcal{S}\to\mathcal{S}^*$ via $h\mapsto (\Tilde h\mapsto H_s(h)(\Tilde h))$.
\end{theorem}
Note that $H_s(h)(\cdot)\in\mathcal{S}^*$ is a real-valued functional on $\mathcal{S}$ and $H_s(\cdot)(\cdot)\in\mathcal{S}^{**}$ can be viewed as an operator on $\mathcal{S}^*$.

\begin{theorem}\label{thm: PF1deriv}
Let $Y$ be a random variable with density $f$ on the compact set $\mathcal{X}\subset\mathbb R$.
Fix $s\in\mathcal{S}$. Assume:

\begin{enumerate}[label=(A\arabic*),itemsep=0.4em]
\item
The characteristic function of the random variable $s(Y)-\mathbb E[s(Y)]$ is integrable, and there exists $K>0$ such that for every
$0<|\epsilon|<K$ the characteristic function of $(s+\epsilon h)(Y)-\mathbb E[(s+\epsilon h)(Y)]$ is integrable.

\item
The random variable $s(Y)$ admits a once continuously differentiable density on $\mathcal{X}$, denoted
$g_s(x):=Pf(s)(x)$, with $g_s'\in L^1(\mathcal{X})$.

\item
There exists a function $M_1\in L^1(\mathbb R)$ such that for all $h\in\mathcal{S}$ and all $0<|\epsilon|<K$,
\begin{align}\label{eq:M1_dom}
|t|\,
\left|
\mathbb E\!\left[
e^{it\left((s+\epsilon h)(Y)-\mathbb E[(s+\epsilon h)(Y)]\right)}\,
\big(h(Y)-\mathbb E[h(Y)]\big)
\right]
\right|
\;\le\; M_1(t)\,\|h\|.
\end{align}
\end{enumerate}
Then $Pf$ is Fr\'echet differentiable at $s$ as a map $Pf:\mathcal{S}\to\mathcal{Z}$, meaning there exists a bounded linear operator $Pf_s \in \mathcal{B}(\mathcal{S},\mathcal{Z}_0)$ such that:
\[
\|Pf(s+h) - Pf(s) - Pf_s(h)\|_{\mathcal{Z}} = o(\|h\|)
\]
as $\|h\| \to 0$. This Fr\'echet derivative $Pf'(s) := Pf_s$ is evaluated at $h\in\mathcal{S}$ is given by
\begin{align}~\label{eq:first-var}
Pf_s(h)(x)
=&\;\frac{1}{2\pi}\int_{\mathbb R} -e^{itx}\,e^{-it\,\mathbb E[s(Y)]}\,
\mathbb E\!\left[
it\,e^{it(s(Y)-\mathbb E[s(Y)])}\big(h(Y)-\mathbb E[h(Y)]\big)
\right]\,dt\\
&\;-\;g_s'(x)\,\mathbb E[h(Y)] ,
\qquad x\in\mathcal{X}.
\end{align}
\end{theorem}

\begin{comment}
Theorem~\ref{thm: PF1deriv} establishes the general form of the Fr\'echet derivative of the Perron--Frobenius operator,
expressed in terms of expectations involving the characteristic function of $s(Y)$.
We now restrict attention to the subspace 
of strictly increasing and differentiable transport maps. 
Within this setting, the derivative admits a closed-form expression that depends only on the 
density $f$ and the inverse map $s^{-1}$.
\begin{corollary}\label{cor: PF1deriv}  
Assume $\bar{\mathcal{S}}\subset\mathcal{S}$ is a subspace comprised of strictly increasing and differentiable functions. Then
\[
g_s(x) = Pf_s(h)(x)=\frac{d}{dz}\frac{f((s+zh)^{-1}(x))}{(s+zh)'((s+zh)^{-1}(x))}\Big|_{z=0}.
\]
\end{corollary}
\end{comment}

\begin{theorem}\label{thm: PF2deriv} 
Under the conditions of Theorem~\ref{thm: PF1deriv}, further assume that the random variable $s(Y)$ has a twice
continuously differentiable density function $g_s(x):=Pf(s)(x)$ with $g_s''\in L^1(\mathcal{X})$.
Assume in addition that 
\begin{enumerate}[label=(A\arabic*),itemsep=0.4em]
\setcounter{enumi}{3}
\item
There exists a function $M_2\in L^1(\mathbb R)$ and a constant $K>0$ such that for all
$h,\Tilde h\in\mathcal{S}$ and all $0<|\epsilon|<K$,
\begin{align}\label{eq:M2_dom}
|t|^2\,
\Big|\mathbb E\!\Big[
e^{it\left((s+\epsilon\Tilde h)(Y)-\mathbb E[(s+\epsilon\Tilde h)(Y)]\right)}
\big(h(Y)-\mathbb E[h(Y)]\big)\big(\Tilde h(Y)-\mathbb E[\Tilde h(Y)]\big)
\Big]\Big|
\;\le\; M_2(t)\,\|h\|\,\|\Tilde h\|.
\end{align}
\end{enumerate}
Then the derivative map $Pf':\mathcal{S}\to\mathcal{B}(\mathcal{S},\mathcal{Z}_0)$ is Fr\'echet differentiable at $s$, i.e., there exists a bounded linear operator $Qf_s \in \mathcal{B}(\mathcal{S}, \mathcal{B}(\mathcal{S},\mathcal{Z}_0))$ such that:
\[\|Pf'(s+\Tilde{h}) - Pf'(s) - Qf_s(\Tilde{h})\|_{\mathcal{B}(\mathcal{S},\mathcal{Z})} = o(\|\Tilde{h}\|)
\]
as $\|\Tilde{h}\| \to 0$. The second Fr\'echet derivative $Pf''(s) := Qf_s$ evaluated at $h,\Tilde{h}\in\mathcal{S}$ is given by:
\begin{align*}
&Qf_s(h)(\Tilde{h})(x) \\
& \quad =\frac{1}{2\pi}\int_{\mathbb R} e^{itx} e^{-it\,\mathbb E [s(Y)] }\mathbb E\Bigg[ e^{it(s(Y)-\mathbb E [s(Y)])}
\big(h(Y)-\mathbb E[h(Y)]\big)\bigg(t^2\big(\Tilde{h}(Y)-\mathbb E[\Tilde{h}(Y)]\big)+it\,\mathbb E[\Tilde{h}(Y)]\bigg)\Bigg]dt\\
&\qquad+g_s''(x)\,\mathbb E [h(Y)] \mathbb E [\Tilde{h}(Y)].
\end{align*}
\end{theorem}

\section{Local Geometry of the Constraints}\label{sec:local-geometry}
In the previous section, we derived explicit first and second variations for the objective and for the pushforward constraint through the Perron–Frobenius operator. These derivatives determine the local linearization of feasibility and, consequently, the local geometry of the feasible set. Let $\bar{s} \in \mathcal{S}$ be any feasible point of (P). Recall that $\bar{s}$ is feasible if the constraint mapping $\mathcal{H}(s) := Pf(s) - g$ satisfies $\mathcal{H}(\bar{s}) = 0_{\mathcal{Z}_0}$, where $\mathcal{Z}_0 := \left\{w\in L^1(\mathcal{X}) : \int_{\mathcal{X}} w(x)dx = 0 \right\}$ is the mean-zero subspace of $L^1(\mathcal{X})$.

In this section, we show that if the solution space $\mathcal{S}$ is finite-dimensional, every feasible point $\bar{s}$ is isolated as a direct consequence of the injectivity of $Pf_{\bar{s}}(\cdot)$, the first variation operator of $Pf$ evaluated at $\bar{s}$. In contrast, when $\mathcal{S}$ is infinite-dimensional, injectivity alone no longer guarantees a uniform lower bound, meaning the minimum modulus may vanish and the isolation of feasible points cannot be concluded in general. Nevertheless, under appropriate conditions, we can establish the closedness of the range of $Pf_{\bar{s}}(\cdot)$, which allows us to recover full surjectivity.

We start by establishing that the first variation operator at a feasible point is unconditionally injective. We start by recalling the following classical consequence of the Stone-Weierstrass theorem regarding the completeness of complex exponentials.
\begin{lemma}\label{lem:completeness}
Let $\mu$ be a finite Borel measure on $\mathbb{R}$ with compact support $\mathcal{X} \subset \mathbb{R}$. Then the family of functions $\{e^{itx} : t \in \mathbb{R}\}$ is complete in $L^2(\mu)$. That is, if for some $u\in L^2(\mu)$ we have
\[ \int_{\mathcal{X}} u(x) e^{itx} \, d\mu(x) = 0 \quad \text{for all } t \in \mathbb{R}, \]
then $u(x) = 0$ $\mu$--almost everywhere.
\end{lemma}
\begin{proposition}\label{prop:injective}
Let $\bar{s} \in \mathcal{S}$ be a feasible transport map for (P) such that $Pf(\bar{s})=g$. Suppose that Assumptions~\ref{ass:domain-support}, \ref{ass:density-cont}, and \ref{ass:mono} hold, and assume further that the regularity conditions \emph{(A1)--(A3)} of Theorem~\ref{thm: PF1deriv} are satisfied at $\bar{s}$. Then, the first variation operator $Pf_{\bar{s}}: \mathcal{S} \to \mathcal{Z}_0$ is injective. In particular, if $Pf_{\bar{s}}(h) = 0_{\mathcal{Z}_0}$, then $h = 0$ $\mu$--a.e.
\end{proposition}
\begin{proof}
Define the centered random variable $\bar{h}:=h(Y)-\mathbb E[h(Y)]$. Using the first-variation formula in~\eqref{eq:first-var}, the assumption $Pf_{\bar{s}}(h)=0_{\mathcal{Z}}$ yields:
\begin{equation}\label{eq:F_inv_equals}
\mathscr{F}^{-1}\!\big(e^{-it\,\mathbb E[\bar{s}(Y)]}\phi_h\big)(x) 
= g'(x)\,\mathbb E[h(Y)] 
\quad\text{for a.e. }x\in\mathcal{X},
\end{equation}
where the function $\phi_h(t):=it\,\mathbb E\!\left[\bar{h}\,e^{it(\bar{s}(Y)-\mathbb E[\bar{s}(Y)])}\right]$. By Theorem~\ref{thm: PF1deriv}\,(A3) evaluated at $\epsilon=0$, $\phi_h(t)$ is bounded by an $L^1(\mathbb{R})$ dominating envelope. Thus, applying the Fourier transform $\mathscr{F}$ to both sides of \eqref{eq:F_inv_equals} in the classical $L^1$ sense, we obtain the transform-domain identity:
\begin{equation}\label{eq:Fourier_identity}
e^{-it\,\mathbb E[\bar{s}(Y)]}\phi_h(t)=\mathscr{F}(g')(t)\,\mathbb E[h(Y)] \qquad \text{for a.e. }t\in\mathbb R.
\end{equation}
For the right-hand side, since
\[
\mathscr F(g')(t) = \int_{\mathcal{X}} e^{-itx} g'(x)dx,
\]
we have
\[
\mathscr F(g')'(0)= -i\int_{\mathcal{X}} x\,g'(x)dx.
\]
Since $\mathrm{supp}(g)=(a,b)\subset \mathcal{X}$ is an open interval by Assumption~\ref{ass:domain-support}, adn $g$ vanishes outside $(a,b)$, we have $g(a)=g(b)=0$. Therefore integration by
parts on $(a,b)$ yields
\[
\int_{\mathcal{X}} xg'(x)dx = \int_a^b xg'(x)dx
= \Big[xg(x)\Big]_a^b - \int_a^b g(x)dx = -1,
\]
since $g$ is a probability density.
Consequently,
\[
\mathscr F(g')'(0) = -i(-1)= i \neq 0.
\]
Equating the derivatives at $t=0$ now yields
\[
0=\mathscr F(g')'(0)\,\mathbb E[h(Y)]=i\mathbb E[h(Y)],
\]
which forces $\mathbb E[h(Y)]=0$.

With $\mathbb E[h(Y)]=0$, \eqref{eq:Fourier_identity} immediately reduces to $e^{-it\,\mathbb E[\bar{s}(Y)]}\phi_h(t) \equiv 0$, which implies $\phi_h(t) \equiv 0$ for all $t \in \mathbb{R}$. Substituting the definition of $\phi_h$ and dividing by $it$ for $t \neq 0$, we have
\begin{equation}\label{eq:completeness_ortho}
\mathbb E\!\left[h(Y)\,e^{it(\bar{s}(Y)-\mathbb E[\bar{s}(Y)])}\right] = 0, \qquad \forall t \in \mathbb{R}.
\end{equation}

Now, let $\mu_{\bar{s}}$ denote the pushforward measure of $\bar{s}(Y)$. By the tower property of expectation, this implies
\[ 
\int_{\mathcal{X}} \mathbb E[h(Y)\mid \bar{s}(Y)=x]\,e^{it(x-\mathbb E[\bar{s}(Y)])}\,d\mu_{\bar{s}}(x) = 0 \qquad \forall t\in\mathbb R. 
\]
Define $u(x):=\mathbb E[h(Y)\mid \bar s(Y)=x]$, so that
\[
\mathbb E[h(Y)e^{it(\bar s(Y)-\mathbb E[\bar s(Y)])}]
= \int_{\mathcal{X}} u(x)\,e^{it(x-\mathbb E[\bar s(Y)])}\,d\mu_{\bar s}(x),
\qquad \forall t\in\mathbb R,
\]
where $\mu_{\bar s}$ denotes the pushforward measure of $\bar s(Y)$.
Hence,
\[
\int_{\mathcal{X}} u(x)\,e^{itx}\,d\mu_{\bar s}(x)=0,\qquad \forall t\in\mathbb R,
\]
which means that $u$ is orthogonal in $L^2(\mu_{\bar s})$ to the collection of
complex exponentials $\{e^{it(\cdot)}:t\in\mathbb R\}$.
Since $\mathcal{X}$ is compact, this collection is complete in $L^2(\mu_{\bar s})$
(by Lemma~\ref{lem:completeness}), and therefore $u(x)=0$ $\mu_{\bar s}$-a.e. Equivalently,
\[
\mathbb E[h(Y)\mid \bar s(Y)] \;=\; u(\bar s(Y)) \;=\; 0 \qquad \mu\text{-a.s.}
\]

Moreover, since $\bar s$ is a bijection by Assumption~\ref{ass:mono}, it follows that $Y$ is a
measurable function of $\bar s(Y)$. This implies that $h(Y)$ is $\sigma(\bar s(Y))$-measurable.
Thus $\mathbb E[h(Y)\mid \bar s(Y)] = h(Y)$ $\mu$-a.s., and we conclude that
$h(Y)=0$ $\mu$-a.s. Since $Y\sim \mu$, this implies $h=0$ $\mu$-a.e, completing the proof.
\end{proof}

\subsection{Isolation of Feasible Points in Finite Dimensions}
When the solution space $\mathcal{S}$ is restricted to a finite-dimensional approximation space $\mathcal{S}_n$ (as is necessary for any computational implementation), the injectivity of $Pf_{\bar{s}}$ dictates a highly rigid local geometry.
\begin{proposition}[Isolation of feasible points in finite-dimensional spaces]
\label{prop:isolated-feasible}
Suppose that $\mathcal{S}_n$ is a finite-dimensional Banach space of transport maps and that $Pf:\mathcal{S}_n\to \mathcal{Z}_0$ is continuously Fr\'echet differentiable. Let $\bar{s} \in \mathcal{S}_n$ be a feasible point. Then, $\bar{s}$ is an isolated feasible point in $\mathcal{S}_n$.
\end{proposition}
\begin{proof}
Let $A:=P f_{\bar{s}}\in \mathcal{B}(\mathcal{S}_n,\mathcal{Z}_0)$. Since $\mathcal{S}_n$ is finite-dimensional, the unit sphere
\[
\mathbb S:=\{h\in\mathcal{S}_n:\|h\|=1\}
\]
is compact. The map $h\mapsto \|A(h)\|_{\mathcal{Z}}$ is continuous on $\mathbb S$, and $A$ is injective, hence
\[
\sigma_{\min}(A):=\min_{h\in\mathbb S}\|A(h)\|_{\mathcal{Z}_0}>0,
\]
we have
\begin{align}\label{eqn:sigma-min}
\|A(h)\|_{\mathcal{Z}_0}\ge \sigma_{\min}(A)\|h\|\ \ \forall h\in\mathcal{S}_n.     
\end{align}
Next, by continuity of the derivative/first variation map $s\mapsto P f_s\in\mathcal{B}(\mathcal{S}_n,\mathcal{Z}_0)$ at $\bar{s}$
(implied by the differentiability of $P f_s$),
there exists a neighborhood $W$ of $\bar{s}$ such that
\begin{align}\label{eqn:neigbourhood}
\|P f_s-A\|_{\mathcal{B}(\mathcal{S}_n,\mathcal{Z})}\le \frac{\sigma_{\min}(A)}{2},
\qquad \forall s\in W,   
\end{align}
Now suppose we have another feasible $s\in W$ and set $h:=s-\bar{s}$.
By the integral mean value theorem for Fr\'echet differentiable maps
\[
0=P f(s)-P f(\bar{s})=\int_0^1 P f_{\bar{s}+t h}(h)\,dt.
\]
Add and subtract $A(h)$ inside the integral:
\[
0=\int_0^1\bigl(P f_{\bar{s}+t h}(h)-A(h)\bigr)\,dt+\int_0^1 A(h)\,dt,
\]
hence
\[
A(h)=-\int_0^1\bigl(P f_{\bar{s}+t h}-A\bigr)(h)\,dt.
\]
Taking norms and using (\ref{eqn:neigbourhood}),
\[
\|A(h)\|_{\mathcal{Z}_0}
\le \int_0^1 \|P f_{\bar{s}+t h}-A\|_{\mathcal{B}(\mathcal{S}_n,\mathcal{Z})}\,\|h\|\,dt
\le \frac{\sigma_{\min}(A)}{2}\|h\|.
\]
Combining with (\ref{eqn:sigma-min}) yields
\[
\sigma_{\min}(A)\|h\|\le \|A(h)\|_{\mathcal{Z}_0}\le \frac{\sigma_{\min}(A)}{2}\|h\|,
\]
which is only possible if $\|h\|=0$, and therefore if $s=\bar{s}$. Thus $\bar{s}$ is an isolated feasible point.
\end{proof}

We stress that the proof of Proposition~\ref{prop:isolated-feasible} essentially depends on the finite-dimensionality of $\mathcal{S}_n$.
When $\dim(\mathcal{S}_n)<\infty$, the Rank--Nullity Theorem~\citep{hoffmann1971linear} implies that the range $R(Pf_{\bar{s}})$ is finite-dimensional, and hence closed by \cite{Rudin1974FunctionalAnalysis}[Theorem~1.21] (every finite-dimensional subspace of a normed space is closed).
Therefore, since $Pf_{\bar{s}}$ is also injective, \cite{hildebrand2012methods}[Theorem~8.6] yields that $Pf_{\bar{s}}$ has a strictly positive minimum modulus.
Equivalently, there exists $\sigma>0$ such that
\[
\|Pf_{\bar{s}}(h)\|_{\mathcal{Z}_0}\ge \sigma\|h\|\qquad \forall\,h\in\mathcal{S}_n.
\]
This lower bound is the key ingredient of the isolation of feasible points.

\subsection{Non-Isolation and Surjectivity in Infinite Dimensions}
In contrast, the isolated feasible points implication may fail in infinite-dimensional spaces where injectivity of $Pf_{\bar{s}}$ does not ensure that it is bounded below or that its range is closed.
Example~\ref{ex:infinite-counterexample} in Appendix~\ref{proof:loc} constructs a sequence $\{h_n\}\subset \mathcal{S}$ with $\|h_n\|=1$ but
$\|Pf_{\bar{s}}(h_n)\|_{\mathcal{Z}_0}\to 0$, showing that the minimum modulus can be zero.
Consequently, in the infinite-dimensional setting one cannot infer isolation of feasible points from injectivity alone. 

When $\mathcal{S}$ is infinite dimensional, injectivity of $Pf_{\bar{s}}$ at a feasible $\bar{s}$ alone does not imply that the operator is bounded below. Consequently, the range of $Pf_{\bar{s}}$ need not be closed. However, by imposing natural structural and integrability conditions on the optimal map, we can recover not just a closed range, but full surjectivity onto $\mathcal{Z}_0$. This requires establishing two properties: adjoint injectivity (yielding a dense image) and a weighted Hardy inequality (yielding a closed range).

\begin{lemma}\label{lem:density-mean-zero}
Suppose Assumptions~\ref{ass:domain-support}, \ref{ass:density-cont}, and \ref{ass:mono} hold, and that $\mathcal{S}$ contains functions from $\mathrm{supp}(f)$ to $\mathrm{supp}(g)$. Then the mean-zero elements of $\mathcal{S}$, 
\[ \mathcal{S}_0 := \{h \in \mathcal{S} : \mathbb{E}[h(Y)] = 0\}, \]
are dense in $L^2_0(\mu) := \{\psi \in L^2(\mu) : \mathbb{E}[\psi(Y)] = 0\}$.
\end{lemma}
\begin{proof}
Let $\psi \in L^2_0(\mu)$ and $\epsilon > 0$ be given. Under Assumptions~\ref{ass:domain-support} and \ref{ass:density-cont}, $\mu$ has a continuous and bounded density $f$ on the compact domain $\mathcal{X}$. Consequently, the norms $\|\cdot\|$ (the $L^2(\mu)$ norm) and $\|\cdot\|_{L^2(\mathcal{X})}$ are equivalent on bounded subsets. By the standard density of smooth functions in $L^2(\mathcal{X})$, there exists $\tilde{\psi} \in C^\infty(\mathcal{X})$ such that $\|\tilde{\psi} - \psi\| < \epsilon/2$.

Defining $\tilde{\psi}_0 := \tilde{\psi} - \mathbb{E}[\tilde{\psi}(Y)]$, we guarantee that $\tilde{\psi}_0 \in C^\infty(\mathcal{X})$ and $\mathbb{E}[\tilde{\psi}_0(Y)] = 0$, which places $\tilde{\psi}_0$ in $L^2_0(\mu)$. The approximation error is controlled by the triangle inequality:
\[
\|\tilde{\psi}_0 - \psi\| \le \|\tilde{\psi} - \psi\| + |\mathbb{E}[\tilde{\psi}(Y) - \psi(Y)]| \le 2\|\tilde{\psi} - \psi\| < \epsilon.
\]
To verify that $\tilde{\psi}_0$ is a valid transport direction at any reference point $s\in\mathcal{S}$, we consider the perturbed map $s_\delta := s + \delta \tilde{\psi}_0$ for some $\delta > 0$. By Assumption~\ref{ass:mono}, $s'$ is bounded away from zero almost everywhere. Because $\tilde{\psi}_0 \in C^\infty$, its derivative is bounded. Therefore, for sufficiently small $\delta$, we have $|s_\delta'(y)| \ge c_0 - \delta\|\tilde{\psi}_0'\|_{L^\infty} > 0$ a.s., ensuring $s_\delta$ remains strictly monotone and satisfies Assumption~\ref{ass:mono}. Because $\mathbb{E}[\tilde{\psi}_0(Y)]=0$, we conclude $\tilde{\psi}_0 \in \mathcal{S}_0$. This establishes that $\mathcal{S}_0$ is dense in $L^2_0(\mu)$.
\end{proof}

\begin{lemma}[Adjoint Injectivity]\label{lem:adjoint-injectivity}
Suppose the Lagrange multiplier $[\lambda] \in L^\infty(\mathcal{X})/\mathbb{R}$ satisfies $\lambda \in W^{1,\infty}(\mathcal{X})$. Under the conditions of Lemma \ref{lem:density-mean-zero} and the conditions of Theorem~\ref{thm: PF1deriv}, the adjoint operator $Pf_{\bar{s}}^* : \mathcal{Z}_0^* \to \mathcal{S}^*$ is injective at feasible point $\bar{s}$. 
\end{lemma}
\begin{proof}
Suppose $\lambda\in L^\infty(\mathcal{X})$ (equivalently, $[\lambda] \in L^\infty(\mathcal{X})/\mathbb{R}$) satisfies $Pf_{\bar{s}}^*(\lambda) = 0$, which requires $\int_{\mathcal{X}} \lambda(x) Pf_{\bar{s}}(h)(x)dx = 0$ for all $h \in \mathcal{S}$. Restricting this identity to $h \in \mathcal{S}_0$, substituting the expression for $Pf_{\bar{s}}(h)$ from Theorem~\ref{thm: PF1deriv}, and applying Fubini's theorem yields:
\[
\mathbb{E}\left[h(Y) \cdot \frac{1}{2\pi} \int_{\mathbb{R}} it \hat{\lambda}(t) e^{it\bar{s}(Y)}dt \right] = 0,
\]
where $\hat{\lambda}(t) = \int_{\mathcal{X}} \lambda(x)e^{-itx}dx$. Because $\lambda \in W^{1,\infty}(\mathcal{X})$, $it\hat{\lambda}(t)$ acts as the Fourier transform of $\lambda'$ in the $L^2$ sense. The inner integral evaluates to $\lambda'(\bar{s}(Y))$, simplifying the condition to $\mathbb{E}[h(Y)\lambda'(\bar{s}(Y))] = 0$ for all $h \in \mathcal{S}_0$. 

By Lemma~\ref{lem:density-mean-zero}, $\mathcal{S}_0$ is dense in $L^2_0(\mu)$, implying that $\lambda'(\bar{s}(Y))$ is orthogonal to $L^2_0(\mu)$. This orthogonality forces $\lambda'(\bar{s}(Y)) = c_0$ almost surely for some constant $c_0 \in \mathbb{R}$. As $\bar{s}$ is surjective onto $\mathrm{supp}(g)$, $\lambda' = c_0$ almost everywhere on $\mathrm{supp}(g)$, integrating to $\lambda(x) = c_0 x + c_1$. 

To determine $c_0$, we extend the evaluation to any general $h \in \mathcal{S}$ with a non-zero mean. The variation condition simplifies to $\int_{\mathcal{X}} \lambda(x) g'(x)dx = 0$. Substituting $\lambda(x) = c_0 x + c_1$ and integrating by parts on $\mathrm{supp}(g) = (a,b)$ yields:
\[
c_0 \left( [xg(x)]_a^b - \int_a^b g(x)dx \right) = c_0(0 - 1) = -c_0 = 0.
\]
Thus $c_0 = 0$, making $\lambda$ constant almost everywhere. Consequently, $[\lambda] = [0]$ in $\in L^\infty(\mathcal{X})/\mathbb{R}$, proving the injectivity of the adjoint operator.
\end{proof}

Having established adjoint injectivity, the remaining ingredient required for full surjectivity is ensuring the operator's range is closed. This corresponds to proving a strict lower bound, which we establish next via a weighted Hardy inequality reduction.

\begin{proposition}[Hardy inequality reduction]\label{prop:hardy-reduction}
Suppose Assumptions~\ref{ass:domain-support}, \ref{ass:density-cont}, and \ref{ass:mono} hold. Let $\bar{s} \in \mathcal{S}$ be a reference transport map such that the ratio function $v(y) := f(y)/\bar{s}'(y)$ vanishes at the boundaries $\partial \mathrm{supp}(f)$, and the weight integrability condition $f/(\bar{s}')^2 \in L^1(\mathrm{supp}(f))$ is satisfied. Then the lower bound $\|Pf_{\bar{s}}(h)\|_{\mathcal{Z}_0} \ge \sigma\|h\|$ for all $h\in\mathcal{S}_0$ is equivalent to the weighted Hardy inequality:
\[
\|(v h)'\|_{L^1(\mathrm{supp}(f))} \ge \sigma\|h\|.
\]
\end{proposition}
\begin{proof}
Under Assumption~\ref{ass:mono}, $\bar{s}$ is a strictly monotone bijection. Differentiating the change-of-variables formula $Pf(\bar{s})(x) = f(\bar{s}^{-1}(x))(\bar{s}^{-1})'(x)$ yields the first variation $Pf_{\bar{s}}(h)(\bar{s}(y)) = -\frac{1}{\bar{s}'(y)}\frac{d}{dy}[v(y)h(y)]$, where $v(y) := f(y)/\bar{s}'(y)$. Computing the $L^1$ norm via the substitution $x = \bar{s}(y)$ with $dx = \bar{s}'(y)dy$ provides:
\[
\|Pf_{\bar{s}}(h)\|_{L^1(\mathcal{X})} = \int_{\mathrm{supp}(g)} |Pf_{\bar{s}}(h)(x)|dx = \int_{\mathrm{supp}(f)} \left|\frac{(v h)'(y)}{\bar{s}'(y)}\right|\bar{s}'(y)dy = \|(v h)'\|_{L^1(\mathrm{supp}(f))}.
\]
By assumption, the function $v(y)$ vanishes at $\partial \mathrm{supp}(f)$. Because $h$ is bounded in our domain, the product $v h$ therefore strictly vanishes at the boundaries. This ensures that $vh$ belongs to the proper absolutely continuous space with zero boundary conditions, making the direct substitution valid and establishing the equivalence.
\end{proof}

We rely on the bilateral Hardy inequality of Opic and Kufner for functions vanishing at both endpoints, $AC_{LR}(c,d)$.
\begin{theorem}[Opic-Kufner \cite{opic1990hardy}]\label{thm:opic-kufner}
Let $(c,d) \subset \mathbb{R}$ be a bounded open interval and let $w$ be a positive measurable weight function. The inequality
\[
\left(\int_c^d |u(y)|^2 w(y) dy\right)^{1/2} \le C \int_c^d |u'(y)| dy \qquad \forall u \in AC_{LR}(c,d)
\]
holds with a finite constant $C>0$ if and only if the value
\[
\mathcal{B} := \sup_{c < \xi < \eta < d} \left(\int_\xi^\eta w(t)dt\right)^{1/2} \min(\xi-c, d-\eta)
\]
is finite.
\end{theorem}

\begin{assumption}\label{ass:hardy}
For any feasible $s\in\mathcal{S}$, we have $\int_{\mathrm{supp}(f)} \frac{s'(y)^2}{f(y)}\,dy < \infty$.
\end{assumption}

\begin{lemma}[Reduction to $L^1$ integrability]\label{lem:L1-reduction}
With weight $w(y) = \bar{s}'(y)^2/f(y)$ and $\mathrm{supp}(f) = (c,d)$, Assumption~\ref{ass:hardy} implies the condition $\mathcal{B} < \infty$.
\end{lemma}
\begin{proof}
Let $I := \int_c^d w(y)dy = \int_c^d \frac{\bar{s}'(y)^2}{f(y)}dy$. Because $\bar{s} \in \mathcal{S}$ is a feasible transport map, Assumption~\ref{ass:hardy} guarantees that $I < \infty$. For any $c < \xi < \eta < d$, the integral over the subinterval is bounded by the total integral: $\int_\xi^\eta w(t)dt \le I$. Furthermore, the distance to the boundaries is bounded by half the total interval length: $\min(\xi-c, d-\eta) \le \frac{d-c}{2}$. 
Taking the supremum over all such $\xi$ and $\eta$ yields:
\[
\mathcal{B} \le I^{1/2} \frac{d-c}{2} < \infty.
\]
Thus, the condition $\mathcal{B} < \infty$ is satisfied.
\end{proof}

Combining these components yields the complete geometric characterization of the functional constraint mapping.

\begin{theorem}\label{thm:pf-surjectivity}
Let $\bar{s} \in \mathcal{S}$ be a feasible transport map for (P) such that $Pf(\bar{s})=g$. Under the assumptions of Lemma~\ref{lem:adjoint-injectivity} and Assumption~\ref{ass:hardy}, the operator $Pf_{\bar{s}} : \mathcal{S} \to \mathcal{Z}_0$ is surjective. 
\end{theorem}
\begin{proof}
By Lemma~\ref{lem:adjoint-injectivity}, the adjoint operator $Pf_{\bar{s}}^* : \mathcal{Z}_0^* \to \mathcal{S}^*$ is injective. The Hahn-Banach theorem~\citep[Theorem 3.2]{Rudin1974FunctionalAnalysis} dictates that a bounded linear operator with an injective adjoint inherently possesses a dense image, establishing that $Pf_{\bar{s}}(\mathcal{S})$ is dense in $\mathcal{Z}_0$. 

Furthermore, by Proposition~\ref{prop:hardy-reduction}, the closed range lower bound is equivalent to the Hardy inequality on the space $AC_{LR}$. Under Assumption~\ref{ass:hardy}, Lemma~\ref{lem:L1-reduction} guarantees that the relevant Opic-Kufner constant satisfies $\mathcal{B} < \infty$. Applying Theorem~\ref{thm:opic-kufner} confirms the Hardy inequality holds. This provides the strict lower bound $\|Pf_{\bar{s}}(h)\|_{\mathcal{Z}_0} \ge \sigma\|h\|$, confirming that $Pf_{\bar{s}}$ has a closed range in $\mathcal{Z}_0$. The combination of a dense image and a closed range proves that $Pf_{\bar{s}}$ is surjective, meaning $Pf_{\bar{s}}(\mathcal{S}) = \mathcal{Z}_0$. 
\end{proof}

\section{Computational Formulation and Consistency}\label{sec:cp}
%Since the infinite-dimensional problem (P) cannot be solved directly, 
Any implementable method must restrict the search to a finite-dimensional model of $\mathcal{S}_n$ and replace the functional constraint by finitely many constraints. In this section, we construct a computable approximate problem (CP) by (i) restricting $s$ to a finite-dimensional parameterization and (ii) discretizing the pushforward constraint by enforcing it on a finite grid.
This section specifies the problem formulation and states in which sense does (CP) approximate (P) as the discretization is refined.

\subsection{Approximation and Discretization}\label{subsec:approx}
Recall our original objective 
\begin{align*}
    &\min_{s\in\Omega} J(s)=\int_{\mathcal{X}} c(x, s(x))f(x)dx \tag{P}\\
    &\text{s.t.} \\
    & Pf(s)-g =0_{\mathcal{Z}}.
\end{align*}
The above infinite-dimensional problem cannot be solved directly, in general. 

\subsubsection{Finite-Dimensional Function Space Approximation}
To make the problem computationally tractable, we first approximate the continuous transport map $s$ using finite-dimensional parameterizations. Let $\mathcal{S}_n \subset \mathcal{S}$ be an $n$-dimensional approximation space for the transport maps, where each map $s_n \in \mathcal{S}_n$ is uniquely determined by a parameter vector $a = (a_0, a_1, \dots, a_{n-1}) \in \mathbb{R}^n$. We define $\Omega_n:=\mathcal{S}_n\cap\Omega$, the set of admissible maps in the finite-dimensional approximation space. Two common and computationally convenient choices are described below.

\begin{example}{Polynomial approximation.}
Consider a transport map as a degree $n-1$ polynomial of the form
\begin{align*}
    s_n(x) = a_0 + a_1x + \cdots+ a_{n-1}x^{n-1},
\end{align*}
with parameter vector $a:=(a_0, a_1, \cdots, a_{n-1})\in\mathbb R^n$.
\end{example}

\begin{example}{Piecewise linear approximation.}
Consider a transport map as an $n$-piece piecewise linear function on $\mathcal{X}$. For a set of grid points $(x_0, \dots, x_n)$, let $a:= (a_0, \dots, a_{n-1})\in\mathbb R^n$ denote the map values at these locations. The approximated piecewise linear transport map is represented by:
\begin{align*}
    s_n(x) = \frac{a_{i+1}-a_i}{x_{i+1}-x_i}(x-x_i)+a_i,\quad x\in [x_i, x_{i+1}].
\end{align*}
\end{example}

\subsubsection{Constraint Discretization}
We next approximate the measure-preserving constraint by imposing it only at a finite collection of $m$ grid points: $(d_1, \dots, d_m)\in\mathbb{R}^m$. We evaluate the Perron--Frobenius operator associated with $s_n \in \mathcal{S}_n$ on each grid point. Hence, our constraints become:
\begin{align*}
    \tau_1^{(m)}(s_n) &= Pf(s_n)(d_1)-g(d_1) = 0,\\
    &\vdots\\
    \tau_m^{(m)}(s_n) &= Pf(s_n)(d_m)-g(d_m) = 0.
\end{align*}
The computational objective (CP) is:
\begin{align*}
    &\min_{s_n\in\Omega_n} J(s_n)=\int_{\mathcal{X}} c(x, s_n(x)) f(x)dx \tag{CP}\\
    &\text{s.t.} \\
    & \tau_1^{(m)}(s_n) = \tau_2^{(m)}(s_n) =\cdots=\tau_m^{(m)}(s_n) = 0.
\end{align*}

For simplicity, denote $\tau^{(m)}(s_n) :=(\tau_1^{(m)}(s_n),\dots,\tau_m^{(m)}(s_n))$. The Jacobian of the constraints at $s_n$ is the matrix $\nabla\tau^{(m)}(s_n)\in\mathbb{R}^{m\times n}$ with:
\[
\bigl(\nabla\tau^{(m)}(s_n)\bigr)_{i,j}=\frac{\partial \tau_i^{(m)}}{\partial a_j}(s_n).
\]

Let $s_n^*$ be a feasible point of (CP). 
If $\nabla\tau^{(m)}(s_n^*)$ has full row rank $m$ (i.e., LICQ holds at $s_n^*$, as shown in Proposition~\ref{prop:discrete-LICQ}), then the implicit function theorem implies that the feasible set
\[
\mathcal{F}_{nm}:=\{s_n\in\mathcal{S}_{n}:\ \tau^{(m)}(s_n)=0\}
\]
is locally an $(n-m)$-dimensional manifold through $s_n^*$.
In particular, $s_n^*$ can be isolated only if $n-m=0$, i.e., only in the square case $m=n$ where $\nabla\tau^{(m)}(s_n^*)$ is nonsingular.
When $m<n$, the Jacobian $\nabla\tau^{(m)}(s_n^*)\in\mathbb{R}^{m\times n}$ has fewer rows than columns and therefore cannot be injective as a linear map from $\mathbb{R}^n$ to $\mathbb{R}^m$. 
In particular,
\[
\dim\ker(\nabla\tau^{(m)}(s_n^*)) \;=\; n-\operatorname{rank}(\nabla\tau^{(m)}(s_n^*)) \;\ge\; n-m \;>\; 0,
\]
so there exist nonzero directions $p$ such that $\nabla\tau^{(m)}(s_n^*)p=0$. 
These vectors are the first-order feasible directions for the discretized constraints, i.e., directions along which the linearly approximated constraints do not change.
Therefore, under LICQ and $m<n$, no feasible point of (CP) can be isolated.

This contrasts with the discussion in Section~\ref{sec:local-geometry}, where the first variation of the functional constraint in \textnormal{(P)} is injective at feasible points and, in finite-dimensional solution spaces, such injectivity leads to a feasible set that consists of isolated points.
For (CP), unless the discretization dimensions satisfy $m=n$ and $\nabla \tau^{(m)}(s_n^*)$ is nonsingular, the feasible set cannot consist of isolated points.
This introduces well-defined tangent directions along which the objective can be further decreased while staying feasible.

\subsection{Consistency of the Computational Objective}\label{subsec:comp_consistency}
% By \cite{philippis2013regularity}, for bounded open sets $X$ and $Y$ in $\mathbb R^d$, a set $\Gamma \subset X \times Y$ is $c$-cyclically monotone if:
% \begin{align*}
% \sum_{i=1}^m c(x_i, y_i) \leq \sum_{i=1}^m c(x_i, y_{\sigma(i)}), \quad \forall m \in \mathbb{N}, \, (x_i, y_i) \in \Gamma, \, \sigma \in \Sigma_h,
% \end{align*}
% where $\Sigma_h$ is the set of all permutations of $h$ elements. 
% A function $u : X \to \mathbb{R} \cup \{+\infty\}$ is $c$-convex if it can be written as:
% \begin{align*}
% u(x) = \sup_{y \in Y} \{-c(x, y) + \lambda_y\},
% \end{align*}
% for some constants $\lambda_y \in \mathbb{R} \cup \{-\infty\}$. 
% For a $c$-convex function $u : X \to \mathbb{R} \cup \{+\infty\}$, the $c$-subdifferential of $u$ at $x \in X$ is the (nonempty) set:
% \begin{align*}
% \partial_c u(x) := \{y \in Y : u(z) \geq -c(z, y) + c(x, y) + u(x), \, \forall z \in X\}.
% \end{align*}

% \begin{remark}[Rockafellar Theorem]\label{thm:Rockafellar}\cite{philippis2013regularity}
% A set $\Gamma\subset X$ is c-cyclically monotone if and only if it is included in the c-subdifferential of a c-convex function.   
% \end{remark}
We begin by recording existence and uniqueness results for the Monge problem (P), as well as an almost-everywhere bijectivity characterization of the optimal transport map.
Since our consistency analysis compares the computational objective (CP) to the true objective (P), we first ensure that (P) admits a well-defined, unique optimizer $s^*$.

\begin{assumption}\label{ass:philippis-twist}
For each fixed $y\in \mathrm{supp}(g)$, the map $x\mapsto D_x c(x,y)$ is injective, and for each fixed
$x\in \mathrm{supp}(f)$, the map $y\mapsto D_y c(x,y)$ is injective. % Use D_yc(x,y) for c'
\end{assumption}

\begin{assumption}\label{ass:philippis-nondeg}
The mixed derivative satisfies
\[
D^2_{xy}c(x,y)\neq 0, \qquad \forall\, x\in \mathrm{supp}(f),\ \forall\, y\in \mathrm{supp}(g).
\]
\end{assumption}

\begin{assumption}\label{ass:philippis-Calpha}
The cost satisfies
\[
c \in C^{2,\alpha}_{\mathrm{loc}}\!\big(\mathrm{supp}(f)\times \mathrm{supp}(g)\big) \quad\text{for some }\alpha\in(0,1).
\]
\end{assumption}

\begin{theorem}[\cite{philippis2013regularity}, Theorem 1.28]\label{thm:philippis}
Under Assumptions~\ref{ass:domain-support}, \ref{ass:density-cont}, \ref{ass:philippis-C2}, and \ref{ass:philippis-twist},
the Monge problem (P) admits a unique optimal transport map $s^*$ from $f$ to $g$ (unique $\mu$--a.e.).
Moreover, under the additional regularity Assumptions~\ref{ass:philippis-nondeg}--\ref{ass:philippis-Calpha},
there exist relatively closed sets $\Sigma_f\subset \mathrm{supp}(f)$ and $\Sigma_g\subset \mathrm{supp}(g)$ of measure zero such that
\[
s^*:\mathrm{supp}(f)\setminus \Sigma_f \ \longrightarrow\ \mathrm{supp}(g)\setminus \Sigma_g
\]
is a bijection of class $C^{1,\alpha}_{\mathrm{loc}}$.
\end{theorem}

\begin{theorem}[\cite{GangboMcCann1996}, Theorem 4.5]\label{thm:gangbo-mccann}
Under Assumption~\ref{ass:domain-support} alone, assume further that the cost is of the transportation form
\[
c(x,y)=\tilde c(x-y),
\]
where $\tilde c:\mathbb R\to[0,\infty)$ is strictly convex. Then the Monge problem (P) admits a unique optimal transport map
$s^*$ (unique $f$--a.e.) pushing $f$ forward to $g$. Moreover, letting $t^*$ denote the optimal transport map of the reversed Monge problem
(from $g$ to $f$), one has
\[
t^*(s^*(x)) = x \quad f\text{--a.e.}, 
\qquad 
s^*(t^*(y)) = y \quad g\text{--a.e.}.
\]
In particular, $s^*$ is bijective between full-measure subsets of $\mathrm{supp}(f)$ and $\mathrm{supp}(g)$.
\end{theorem}

Under conditions of either Theorem~\ref{thm:philippis} or Theorem~\ref{thm:gangbo-mccann}, the Monge problem (P) admits a unique optimal transport map $s^*$ (unique $f$--a.e.) pushing $f$ forward to $g$,
and $s^*$ is bijective between full-measure subsets of $\mathrm{supp}(f)$ and $\mathrm{supp}(g)$.

With the well-posedness of (P) stated above,
%we can cleanly separate the two sources of approximation in (CP): the restriction of the transport maps to a finite-dimensional subspace of $\mathcal{S}$, and replacing the functional pushforward constraints by finitely many grid constraints. 
%This leads naturally to
we now consider the intermediate problems (CP-T) and (CP-C) below, each isolating one source of approximation error, to enable a consistency analysis of (CP) $\to$ (P) as the discretization is refined. 

\begin{align*}
    &\min_{s_n \in\Omega_n} J(s_n)=\int_{\mathcal{X}} c(x, s_n(x)) f(x)dx \tag{CP-T}\\
    &\text{s.t.} \\
    & Pf(s_n)-g =0_{\mathcal{Z}}.
\end{align*}

\begin{align*}
    &\min_{s\in\Omega} J(s)=\int_{\mathcal{X}} c(x,s(x))f(x)dx \tag{CP-C}\\
    &\text{s.t.} \\
    & \tau^{(m)}_1(s) = \tau^{(m)}_2(s) =\cdots=\tau^{(m)}_m(s) = 0.
\end{align*}

Denote the feasible set of (P), (CP-T), (CP-C), and (CP) (resp.) as $\mathcal{F}$, $\mathcal{F}_{n\infty}$, $\mathcal{F}_{\infty m}$ and $\mathcal{F}_{nm}$ (resp.). 
Let their solutions be $s^*$, $s_{n\infty}^*$, $s_{\infty m}^*$ and $s_{nm}^*$ (resp.).

\begin{theorem}[Consistency of the Computational Objective]\label{thm:consistency}
Under Assumptions~\ref{ass:domain-support}--\ref{ass:cost-lipschitz},
\ref{ass:mono}, and the conditions of either
Theorem~\ref{thm:philippis} or Theorem~\ref{thm:gangbo-mccann}, the
following hold:
\begin{enumerate}[label={\textnormal{(\roman*)}},itemindent=1em]
    \item\label{thm:consistency-obj}
    \[
    J(s_{nm}^*) \xrightarrow[n,m\to\infty]{} J(s^*).
    \]
    \item\label{thm:consistency-map}
    \[
    s_{nm}^* \xrightarrow[n,m\to\infty]{} s^*
    \qquad \mu\text{--a.e.}
    \]
\end{enumerate}
\end{theorem}

\begin{proof}
We first prove \ref{thm:consistency-obj}. This follows from the two
intermediate consistency results:
\begin{enumerate}[label={(\arabic*)},itemindent=1em]
    \item\label{consistency1}
    \[
    J(s_{nm}^*) \xrightarrow[n\to\infty]{} J(s_{\infty m}^*)
    \quad\text{uniformly in } m
    \]
    (Theorem~\ref{thm:cp-conv-to-cpc} for piecewise linear
    $\mathcal S_n$, and
    Theorem~\ref{thm:cp-conv-to-cpc-poly} for polynomial
    $\mathcal S_n$);
    \item\label{consistency2}
    \[
    J(s_{\infty m}^*) \xrightarrow[m\to\infty]{} J(s^*)
    \]
    (Theorem~\ref{thm: cp-conv-to-cpt-cpc-conv-to-nlp}%
    \ref{thm: cp-conv-to-cpt-cpc-conv-to-nlp2}).
\end{enumerate}
Step~\ref{consistency1} shows that for each fixed $m$, the limit
$\lim_{n\to\infty}J(s_{nm}^*)$ exists and equals $J(s_{\infty m}^*)$,
and that this convergence is uniform in $m$. Step~\ref{consistency2}
identifies the remaining outer limit. Therefore, by the
Moore--Osgood theorem \citep[Theorem~7.11]{rudin1976principles},
\[
J(s_{nm}^*) \xrightarrow[n,m\to\infty]{} J(s^*).
\]

We next prove \ref{thm:consistency-map}. The argument is by sequential compactness and uniqueness of the minimizer.
By Assumption~\ref{ass:mono}, each $s_{nm}^*$ is monotone on
$\mathrm{supp}(f)=[c,d]$ and takes values in the compact interval
$\mathrm{supp}(g)=[a,b]$. Hence $\{s_{nm}^*\}$ is a uniformly bounded
family of monotone functions on $[c,d]$. By the monotone function
special case of Helly's selection theorem
\citep[Chapter~10, Theorem~5]{kolmogorov1970intro}, every sequence in $\{s_{nm}^*\}$
admits a subsequence $\{s_{n_km_k}^*\}$, such that
\[
s_{n_km_k}^*(x)\to \bar s(x)
\qquad\text{for every }x\in[c,d],
\]
for some monotone function $\bar s:[c,d]\to[a,b]$.

Since the feasible sets $\mathcal F_{nm}$ are nested and their limit
recovers the feasible set of the original problem
(Theorem~\ref{thm: feasible-set-conv}), any such subsequential limit is
feasible, i.e.,
\[
\bar s\in\mathcal F.
\]

We now pass to the limit in the objective functional along this
subsequence. For each $x\in[c,d]$, continuity of the cost in its second
argument and the pointwise convergence
$s_{n_km_k}^*(x)\to \bar s(x)$ imply
\[
c(x,s_{n_km_k}^*(x))\,f(x)\to c(x,\bar s(x))\,f(x).
\]
Moreover, since all values of $s_{n_km_k}^*$ and $\bar s$ lie in
$[a,b]$, Assumption~\ref{ass:cost-lipschitz} yields the uniform bound
\[
|c(x,s_{n_km_k}^*(x))|
\le |c(x,a)| + K|s_{n_km_k}^*(x)-a|
\le \sup_{x\in[c,d]}|c(x,a)| + K(b-a),
\]
and therefore
\[
|c(x,s_{n_km_k}^*(x))\,f(x)|
\le \Big(\sup_{x\in[c,d]}|c(x,a)| + K(b-a)\Big) f(x).
\]
The dominating function on the right is integrable because
$f$ is a probability density on the compact interval $[c,d]$.
Hence the dominated convergence theorem gives
\[
J(s_{n_km_k}^*) \to J(\bar s).
\]

On the other hand, \ref{thm:consistency-obj} implies
\[
J(s_{n_km_k}^*) \to J(s^*).
\]
Therefore
\[
J(\bar s)=J(s^*).
\]
Since $s^*$ is the unique minimizer of $(P)$ over $\mathcal F$
under either Theorem~\ref{thm:philippis} or
Theorem~\ref{thm:gangbo-mccann}, it follows that
\[
\bar s = s^*
\qquad \mu\text{--a.e.}
\]

We have thus shown that every subsequence of $\{s_{nm}^*\}$ admits a
further subsequence converging to $s^*$ $\mu$--a.e. Hence the full
family converges:
\[
s_{nm}^* \xrightarrow[n,m\to\infty]{} s^*
\qquad \mu\text{--a.e.}
\]
This proves \ref{thm:consistency-map}.
\end{proof}


\begin{comment}
\begin{theorem}~\label{thm: obj-convex} 
If the cost function $c(x, y)$ is convex, then the objective functional $J(\cdot)$ is convex on $\mathcal{S}$.
\end{theorem}

\begin{proof} 
For any $s_1, s_2\in\mathcal{S}$ and $0\leq\alpha\leq 1$, we have
\begin{align*}
    &J(\alpha s_1 + (1-\alpha) s_2)\\
    =&\int_{\mathcal{X}} c(x,(\alpha s_1(x) + (1-\alpha) s_2(x)))f(x)dx\\
    =&\int_{\mathcal{X}} c(\alpha x + (1-\alpha)x-(\alpha s_1(x) + (1-\alpha) s_2(x)))f(x)dx\\
    \leq&\int_{\mathcal{X}}\alpha c(x,s_1(x)) + (1-\alpha)c(x,s_2(x)) f(x)dx\\
    =&\alpha J(s_1) + (1-\alpha)J(s_2).
\end{align*}

\end{proof}
\end{comment}


\begin{proposition}~\label{prop: obj-continuity} 
Under Assumption~\ref{ass:domain-support} and Assumption~\ref{ass:cost-lipschitz}, the objective functional $J(\cdot)$ is K-Lipschitz on $\mathcal{S}$ under the supremum norm $\|\cdot\|_{\infty}$.
\end{proposition}
\begin{proof}
For $s_1$, $s_2\in\mathcal{S}$, we have
    \begin{align*}
        &\left|J(s_1)-J(s_2)\right|\\
        =&\left|\int_{\mathcal{X}} \left(c(x, s_1(x))-c(x, s_2(x))\right) f(x)dx\right|\\
        \leq&\int_{\mathcal{X}} \left|c(x, s_1(x))-c(x, s_2(x))\right| f(x)dx\\
        \leq&\int_{\mathcal{X}} K\left|s_1(x)-s_2(x)\right| f(x)dx\\
        \leq&K\cdot\|s_1-s_2\|_{\infty}.
    \end{align*}
\end{proof}

\begin{theorem}~\label{thm: feasible-set-conv} 
Under Assumption~\ref{ass:density-cont} and Assumption~\ref{ass:mono}, we have the following:
\begin{enumerate}[label={(\roman*)},itemindent=1em]
\item For any fixed $n>0$, there exists a sequence of sets $\{\mathcal{F}_{nm}: m \geq 1\}$ such that, let $\cap_{m\geq 1} \mathcal{F}_{nm}:= \tilde{\mathcal{F}}_{n\infty}$, then $ \tilde{\mathcal{F}}_{n\infty} = \mathcal{F}_{n\infty}$.\label{thm: feasible-set-conv1} 
\item There exists a sequence sets $\{\mathcal{F}_{\infty m}: m \geq 1\}$ such that, let $\cap_{m\geq 1} \mathcal{F}_{\infty m}:= \tilde{\mathcal{F}}$, then $ \tilde{\mathcal{F}} = \mathcal{F}$. \label{thm: feasible-set-conv2} 
\end{enumerate}
\end{theorem}

\begin{proof}
To show \ref{thm: feasible-set-conv1}, denote $\mathrm{supp}(g)$ as $[a, b]$. Note that all connected sets in $\mathbb R$ are intervals, therefore both $\mathrm{supp}(f)$ and $\mathrm{supp}(g)$ are intervals by Assumption~\ref{ass:density-cont}. Let $\{\mathcal{P}_m\}$ a sequence of partitions of $[a, b]$ such that $\mathcal{P}_1\subseteq\mathcal{P}_2\subseteq\cdots$, and $\lim_{m\rightarrow\infty}\mathcal{P}_m=\mathcal{D}$ where $\mathcal{D}\subset \mathrm{supp}(g)$ is dense in $\mathrm{supp}(g)$. For the fixed $n>0$, Let $\mathcal{F}_{nm}$ be the feasible set for (CP) with $n$ parameters for the transport map and $m$ number of constraints whose grid is defined by $\mathcal{P}_m$. Hence, $\{\mathcal{F}_{nm}: m \geq 1\}$ is a monotone non-increasing sequence of sets, i.e., $\mathcal{F}_{n1}\supseteq \mathcal{F}_{n2}\supseteq \mathcal{F}_{n3},\cdots$. Let $\tilde{\mathcal{F}}_{n\infty}:=\cap_{m\rightarrow\infty}\mathcal{F}_{nm}$. $\mathcal{F}_{nm} \supseteq \mathcal{F}_{n\infty}$ for each $m$, therefore, $ \tilde{\mathcal{F}}_{n\infty}\supseteq \mathcal{F}_{n\infty}$. We further show $\mathcal{F}_{n\infty}=\tilde{\mathcal{F}}_{n\infty}$ by showing $\tilde{\mathcal{F}}_{n\infty}\setminus\mathcal{F}_{n\infty}=\emptyset$.
By Lemma~\ref{lem:pf-continuous-version} in Appendix~\ref{sec:pf-regularity}, under Assumption~\ref{ass:density-cont} and Assumption~\ref{ass:mono}, the function $Pf(s)$ admits a continuous version on $\mathrm{supp}(g)$, which we identify with $Pf(s)$ throughout. Two continuous functions that agree on a dense subset of $\mathrm{supp}(g)$ must also agree on $\mathrm{supp}(g)$. Therefore, there is no such element $s\in\tilde{\mathcal{F}}_{n\infty}\setminus\mathcal{F}_{n\infty}$ such that $Pf(s)$ satisfies the density constraint only on a dense subset of $\mathrm{supp}(g)$, hence we complete the proof.

Now we show \ref{thm: feasible-set-conv2}. With the same sequence of partitions $\{\mathcal{P}_m\}$ on $\mathrm{supp}(g)$ that goes to a dense set $\mathcal{D}$, let $\mathcal{F}_{\infty m}$ be the feasible set of (CP-C) with $m$ number of density constraints on grid points defined by $\mathcal{P}_m$. It is clear that $\{\mathcal{F}_{\infty m}\}$ is also a monotone non-increasing sequence of sets such that $\mathcal{F}_{\infty 1}\supseteq \mathcal{F}_{\infty 1}\supseteq\cdots$. Let $\tilde{\mathcal{F}}:=  \cap_{m\rightarrow\infty}\mathcal{F}_{\infty m}$, $\mathcal{F}_{\infty  m} \supseteq \mathcal{F}$ for each $m$, thus $ \tilde{\mathcal{F}}\supseteq \mathcal{F}$. Similarly as in the proof of Theorem~\ref{thm: feasible-set-conv}\ref{thm: feasible-set-conv1}, we can argue that $\tilde{\mathcal{F}}\setminus\mathcal{F}=\emptyset$ since two continuous functions that agree on a dense subset of $\mathrm{supp}(g)$ must also agree on $\mathrm{supp}(g)$.
\end{proof}

\begin{theorem}~\label{thm: cp-conv-to-cpt-cpc-conv-to-nlp} 
Under Assumptions~\ref{ass:domain-support}--\ref{ass:cost-lipschitz}, \ref{ass:mono}, and the conditions of either Theorem~\ref{thm:philippis} or Theorem~\ref{thm:gangbo-mccann}, we have the following:
\begin{enumerate}[label={(\roman*)},itemindent=1em]
\item For a fixed $n>0$,  $J(s_{nm}^*)\rightarrow J(s_{n\infty}^*)$ as $m\rightarrow\infty$.\label{thm: cp-conv-to-cpt-cpc-conv-to-nlp1} 
\item $J(s_{\infty m}^*)\rightarrow J(s^*)$ as $m\rightarrow\infty$. \label{thm: cp-conv-to-cpt-cpc-conv-to-nlp2} 
\end{enumerate}
\end{theorem}

\begin{proof} 
We begin by showing \ref{thm: cp-conv-to-cpt-cpc-conv-to-nlp1}. Clearly, we have $\mathcal{F}_{n\infty}\subset\mathcal{F}_{nm}$ because any transport map that satisfies the PF constraint would satisfy the set of discretized PF constraints. $\mathcal{F}_{nm}$ is non-increasing in $m$, so $J(\mathcal{F}_{nm})$ is non-decreasing in $m$. Therefore, for any fixed $n>0$, we have $J(s_{n\infty}^*)\geq J(s_{nm}^*)\quad\forall m >0$. 
To show that $J(s_{nm}^*)\rightarrow J(s_{n\infty}^*)$ as $m\rightarrow\infty$, we are left with the other direction, that is $\lim_{m\rightarrow\infty} J(s_{nm}^*)\geq J(s_{n\infty}^*)$ for all $n>0$.
   
For a fixed $n>0$, let $s_{nm}^*\in \mathcal{F}_{nm}$ be the solution to (CP) with $m$ constraints. $\{s_{nm}^*\}$ is then a sequence indexed by $m$. Note that the objective function value $J(\cdot)$ is bounded over $\mathcal{S}$ because the integral of a bounded function $c(x, y)$ with respect to a density function must also be bounded. Hence, that $J(s_{nm}^*)$ converges as  $m\rightarrow\infty$ because it is a bounded non-decreasing sequence in $m$. Let $V_{n\infty}^*:=\lim_{m\rightarrow\infty}J(s_{nm}^*)$. As mentioned above, we need to show that $V_{n\infty}^*\geq J(s_{n\infty}^*)$. By Theorem~\ref{thm: feasible-set-conv}\ref{thm: feasible-set-conv1}, we know that for any $\epsilon>0$, we can choose $M_\epsilon$ such that the Hausdorff distance between $\mathcal{F}_{nm}$ and $\mathcal{F}_{n\infty}$ on the metric space $(\mathcal{S}, \|\cdot\|_{\infty})$ is smaller than $\epsilon$ for $m>M_\epsilon$. By the K-Lipschitz continuity of the objective functional $J(\cdot)$ shown in Proposition~\ref{prop: obj-continuity}, we know that for any $\epsilon>0$, there also exists a $\tilde{M}_{\epsilon}>0$ such that a point $\tilde{s}_{nm}^* \in\mathcal{F}_{n\infty}$ can be  found with $|J(\tilde{s}_{nm}^*)  - J(s_{nm}^*)|<\epsilon$ for all $m\geq \tilde{M}_{\epsilon}$. Now, suppose $\lim_{m\rightarrow\infty}J(s_{nm}^*) =V_{n\infty}^* < J(s_{n\infty}^*)$. Set $\epsilon = J(s_{n\infty}^*)-V_{n\infty}^*$, then, $J(\tilde{s}_{nm}^*)< J(s_{n\infty}^*)$ for $m > \tilde{M}_{\epsilon},$ 
which leads to a contradiction to the definition of $s_{n\infty}^*$.
Therefore, we have
\[\lim_{m\rightarrow\infty} J(s_{nm}^*)=V_{n\infty}^*\geq J(s_{n\infty}^*),
\]
for a fixed $n>0$.
So far, we have shown both directions and established 
$J(s_{nm}^*)\rightarrow J(s_{n\infty}^*)$ as $m\rightarrow\infty$.
 
The proof of \ref{thm: cp-conv-to-cpt-cpc-conv-to-nlp2} follows from Theorem~\ref{thm: feasible-set-conv}\ref{thm: feasible-set-conv2}, similarly to the proof of \ref{thm: cp-conv-to-cpt-cpc-conv-to-nlp1}. Here we consider $s_{\infty m}^*\in\mathcal{F}_{\infty m}$ in place of $s_{nm}^*\in\mathcal{F}_{nm}$ and $s^*\in\mathcal{F}$ in place of $s_{n\infty}^*\in\mathcal{F}_{n\infty}$.
\end{proof}


\begin{theorem}~\label{thm:cp-conv-to-cpc} 
Under Assumptions~\ref{ass:domain-support}--\ref{ass:cost-lipschitz}, \ref{ass:mono}, the conditions of either Theorem~\ref{thm:philippis} or Theorem~\ref{thm:gangbo-mccann}, and assuming $\mathcal{S}_n$ is the space of continuous piecewise linear functions, $J(s_{nm}^*)\xrightarrow{n\rightarrow\infty}J(s_{\infty m}^*)$ uniformly in $m$.
\end{theorem}
\begin{proof}
We first show the point-wise convergence of $J(s_{nm}^*)$ to $J(s_{\infty m}^*)$ fixing a $m\geq 1$. We know that $J(s_{\infty m}^*)\leq J(s_{nm}^*)$ for all $n>0$ because they are the optimal objective functional values of the same problem over the space $\mathcal{F}_{\infty m}$ and $\mathcal{F}_{nm}$ respectively, and $\mathcal{F}_{1m}\subseteq\mathcal{F}_{2m}\subseteq\cdots\subseteq\mathcal{F}_{\infty m}$. Now we show the other direction, i.e., $J(s_{\infty m}^*)\geq \lim_{n\rightarrow\infty}J(s_{nm}^*)$. Let $\mathrm{supp}(g)= [a,b]$ and $\mathrm{supp}(f)= [c,d]$. For $s_1, s_2\in\mathcal{S}$, assume without loss of generality that $s_1, s_2$ are both increasing, we have

\begin{align*}
    &\left|Pf(s_1)(x) - Pf(s_2)(x)\right|\\
    =&\left|\frac{d}{dx}\int_{s_1^{-1}((a, x])} f(y)dy - \frac{d}{dx}\int_{s_2^{-1}((a, x])} f(y)dy\right|\\
    =&\left|\frac{d}{dx}\int_{s_1^{-1}(a)}^{s_1^{-1}(x)} f(y)dy - \frac{d}{dx}\int_{s_2^{-1}(a)}^{s_2^{-1}(x)} f(y)dy\right|\\
    =&\left|f(s_1^{-1}(x))(s_1^{-1})'(x) - f(s_2^{-1}(x))(s_2^{-1})'(x) \right|,
\end{align*}
where the last equality is by Leibniz rule. We know from the last equation that for any $x\in \mathrm{supp}(f)$, if $s_1^{-1}$ and $s_2^{-1}$ have the same function values, and the same first derivatives at $x$, then $\left|Pf(s_1)(x) - Pf(s_2)(x)\right|=0$. For $s_{\infty m}^*\in \mathcal{F}_{\infty m}$, we can always construct a sequence of continuous piecewise linear functions $\{\hat{s}^*_{nm}\}\in\mathcal{S}_n$ in the way that for any $\epsilon>0$, there exists a $N_{\epsilon}$ such that if $n>N_{\epsilon}$, we have 
$\|s_{\infty m}^*-\hat{s}^*_{nm}\|_{\infty}<\epsilon$ and $\left|Pf(s_1) - Pf(s_2)\right|=0$ on $m$ grid points associated with the $m$ constraints, namely, $\hat{s}^*_{nm}\in \mathcal{F}_{nm}$. It is well known that the space of piecewise linear function is dense in the space of continuous functions. Therefore, a piecewise linear continuous approximation of  $s_{\infty m}^*$ that uses the first-order approximation in a small neighborhood around each grid point would suffice.
With such $\hat{s}^*_{nm}$, we have
   \begin{align*}
       J(s_{\infty m}^*) &= \lim_{n\rightarrow\infty} J(\hat{s}^*_{nm}) + (J(s_{\infty m}^*)-J(\hat{s}^*_{nm}))\\
       &= \lim_{n\rightarrow\infty} J(\hat{s}^*_{nm}) + \lim_{n\rightarrow\infty} (J(s_{\infty m}^*)-J(\hat{s}^*_{nm}))\\
       &=\lim_{n\rightarrow\infty} J(\hat{s}^*_{nm})\\
       &\geq \lim_{n\rightarrow\infty} J(s_{nm}^*).
   \end{align*}
The last inequality is due to the definition of $s_{nm}^*$ and the penultimate equation is by Proposition~\ref{prop: obj-continuity}, specifically,
\begin{align*}
    &\lim_{n\rightarrow\infty}|J(s_{nm}^*)-J(\hat{s}^*_{nm})|\\
    &\leq\lim_{n\rightarrow\infty}K\|s_{nm}^*-\hat{s}^*_{nm}\|_{\infty}\\
    &=0.
\end{align*}
Therefore, $\lim_{n\rightarrow\infty}(J(s_{nm}^*)-J(s_{\infty m}^*))=0$ for a fixed $m$.

To show the uniform convergence of $J(s_{nm}^*)$, we use Dini's theorem. Recall that $J(s_{1m}^*)\geq J(s_{2m}^*)\geq\cdots\geq J(s_{\infty m}^*)$ because $\mathcal{F}_{1m}\subseteq\mathcal{F}_{2m}\subseteq\cdots\subseteq\mathcal{F}_{\infty m}$. Now let 
\begin{align*}
l_n(x)
:=
\begin{cases}
J\!\left(s_{n\frac{1}{x}}^*\right), & x \in \left\{1,\frac12,\frac13,\ldots\right\},\\[4pt]
J\!\left(s_{n\infty}^*\right), & x=0,\\[8pt]
J\!\left(s_{n1}^*\right)
+ \dfrac{x-1}{\frac12-1}\Big(J\!\left(s_{n2}^*\right)-J\!\left(s_{n1}^*\right)\Big),
& \frac12 < x < 1,\\[10pt]
J\!\left(s_{n i}^*\right)
+ \dfrac{x-\frac{1}{i}}{\frac{1}{i-1}-\frac{1}{i}}
\Big(J\!\left(s_{n,i-1}^*\right)-J\!\left(s_{n i}^*\right)\Big),
& \frac{1}{i} < x < \frac{1}{i-1},\quad i\ge 3,\\[12pt]
J\!\left(s_{n\infty}^*\right)
+ \dfrac{x-0}{\frac12-0}\Big(J\!\left(s_{n2}^*\right)-J\!\left(s_{n\infty}^*\right)\Big),
& 0 < x < \frac12.
\end{cases}
\end{align*}

\begin{align*}
l(x)
:=
\begin{cases}
J\!\left(s_{\infty\frac{1}{x}}^*\right), & x \in \left\{1,\frac12,\frac13,\ldots\right\},\\[4pt]
J\!\left(s^*\right), & x=0,\\[8pt]
J\!\left(s_{\infty 1}^*\right)
+ \dfrac{x-1}{\frac12-1}\Big(J\!\left(s_{\infty 2}^*\right)-J\!\left(s_{\infty 1}^*\right)\Big),
& \frac12 < x < 1,\\[10pt]
J\!\left(s_{\infty i}^*\right)
+ \dfrac{x-\frac{1}{i}}{\frac{1}{i-1}-\frac{1}{i}}
\Big(J\!\left(s_{\infty,i-1}^*\right)-J\!\left(s_{\infty i}^*\right)\Big),
& \frac{1}{i} < x < \frac{1}{i-1},\quad i\ge 3,\\[12pt]
J\!\left(s^*\right)
+ \dfrac{x-0}{\frac12-0}\Big(J\!\left(s_{\infty 2}^*\right)-J\!\left(s^*\right)\Big),
& 0 < x < \frac12.
\end{cases}
\end{align*}
Now note that for any $x\in [0, 1]$, it is true that $l_1(x) \geq l_2(x)\geq\cdots\geq l(x)$. That is, $l_n$ is point-wise convergent to $l$ monotonically. By construction, $l_n$ and $l$ are continuous on $(0,1]$. By Theorem~\ref{thm: cp-conv-to-cpt-cpc-conv-to-nlp}\ref{thm: cp-conv-to-cpt-cpc-conv-to-nlp1}, $J(s_{nm}^*)\xrightarrow{m\rightarrow\infty}J(s_{n\infty}^*)$ for a fixed $n$, hence $l_n(x)$ is right continuous at $0$. Similarly, $l(x)$ is also right continuous at $0$ because $J(s_{\infty m}^*)\xrightarrow{m\rightarrow\infty}J(s^*)$ by Theorem~\ref{thm: cp-conv-to-cpt-cpc-conv-to-nlp}\ref{thm: cp-conv-to-cpt-cpc-conv-to-nlp2}. Therefore, $l_n\rightarrow l$ uniformly on $[0, 1]$ by Dini's Theorem. This implies $J(s_{nm}^*)\xrightarrow{n\rightarrow\infty}J(s_{\infty m}^*)$ uniformly in $m$.
\end{proof}

\begin{comment}
\begin{theorem}~\label{thm: cpt-conv-to-nlp} 
Under Assumption~\ref{ass:domain-support} and Assumption~\ref{ass:cost-lipschitz}, $J(s_{n\infty}^*)\rightarrow J(s^*)$ as $n\rightarrow\infty$. 
\end{theorem}

\begin{proof} 
We know that $J(s^*)\leq J(s_{n\infty}^*)$ for all $n>0$, because $\mathcal{F}_{1\infty}\subseteq\mathcal{F}_{2\infty}\subseteq\cdots\subseteq\mathcal{F}$. It is well known that the space of piecewise linear function is dense in $\mathcal{F}$, therefore, we have that $\cup _{n=1}^{\infty}\mathcal{F}_{n\infty}=\mathcal{F}$. Denote the closest point in $\mathcal{F}_{n\infty}$ of $s^*$ as $\hat{s}^*_{n\infty}$, we also have
   \begin{align*}
       J(s^*) &= \lim_{n\rightarrow\infty} J(\hat{s}^*_{n\infty}) + (J(s^*)-J(\hat{s}^*_{n\infty}))\\
       &= \lim_{n\rightarrow\infty} J(\hat{s}^*_{n\infty}) + \lim_{n\rightarrow\infty} (J(s^*)-J(\hat{s}^*_{n\infty}))\\
       &=\lim_{n\rightarrow\infty} J(\hat{s}^*_{n\infty})\\
       &\geq \lim_{n\rightarrow\infty} J(s_{n\infty}^*).
   \end{align*}
The last inequality is due to the definition of $s_{n\infty}^*$ and the penultimate equation is by Proposition~\ref{prop: obj-continuity}, specifically,
\begin{align*}
    &\lim_{n\rightarrow\infty}|J(s^*)-J(\hat{s}^*_{n\infty})|\\
    &\leq\lim_{n\rightarrow\infty}K|s^*-\hat{s}^*_{n\infty}|_{\infty}\\
    &=0.
\end{align*}
Therefore, $\lim_{n\rightarrow\infty}(J(s^*)-J(\hat{s}^*_{n\infty}))=0$.
\end{proof}

\begin{theorem}~\label{thm: cp-conv-to-cpt-uniform} 
Under Assumption~\ref{ass:domain-support} and Assumption~\ref{ass:cost-lipschitz}, $J(s_{nm}^*)\xrightarrow{m\rightarrow\infty}J(s_{n\infty}^*)$ uniformly in $n$.
\end{theorem}

\begin{proof}
To show that the convergence is uniform in $n$, recall that $\{\mathcal{F}_{nm}: m \geq 1\}$ is a monotone non-increasing sequence of sets such that $\mathcal{F}_{n1}\supseteq \mathcal{F}_{n2}\supseteq\cdots\supseteq \mathcal{F}_{n\infty}$. This gives us $J(s_{n1}^*)\leq J(s_{n2}^*)\leq\cdots\leq J(s_{n\infty}^*)$ which is again a monotone sequence of functions that converges point-wise for any fixed $n$. Using a similar construction in Theorem~\ref{thm:cp-conv-to-cpc} and Dini's theorem, uniform convergence follows.    
\end{proof}
\end{comment}

\begin{theorem}\label{thm:cp-conv-to-cpc-poly}
Under Assumptions~\ref{ass:domain-support}--\ref{ass:cost-lipschitz}, \ref{ass:mono}, the conditions of either Theorem~\ref{thm:philippis} or Theorem~\ref{thm:gangbo-mccann}, and assuming
$\mathcal{S}_n$ is the intersection of $\mathcal{S}$ and the space of real polynomials on
$\mathrm{supp}(f)$ of degree at most $n-1$, then $J(s_{nm}^*)\xrightarrow{n\rightarrow\infty}J(s_{\infty m}^*)$ uniformly in $m$.
\end{theorem}
assuming $\mathcal{S}_n=\mathcal{S}\cap\mathbb P_{n-1}$, where
$\mathbb P_{n-1}$ denotes the space of real polynomials on
$\mathrm{supp}(f)$ of degree at most $n-1$, we have
\[
J(s_{nm}^*) \xrightarrow[n\to\infty]{} J(s_{\infty m}^*)
\qquad\text{uniformly in }m.
\]
\end{theorem}

\begin{proof}
Let $\mathrm{supp}(f)=[c,d]$ and $\mathrm{supp}(g)=[a,b]$, and denote
by $\mathbb P_k$ the space of real polynomials of degree at most $k$
on $[c,d]$.

We first prove pointwise convergence in $n$ for each fixed $m$, and
then upgrade to uniform convergence in $m$ by Dini's theorem, exactly
as in Theorem~\ref{thm:cp-conv-to-cpc}.

For each fixed $m\ge 1$,
\[
\mathcal F_{1m}\subseteq \mathcal F_{2m}\subseteq \cdots
\subseteq \mathcal F_{\infty m},
\]
hence
\[
J(s_{\infty m}^*)\le J(s_{nm}^*)\qquad \forall n\ge 1.
\]
Therefore it suffices to show
\[
\limsup_{n\to\infty}J(s_{nm}^*)\le J(s_{\infty m}^*).
\]

Fix $m\ge 1$, and let $d_1,\dots,d_m\in\mathrm{supp}(g)$ be the $m$
constraint grid points. Define
\[
y_j:=(s_{\infty m}^*)^{-1}(d_j),\qquad j=1,\dots,m.
\]
By Assumption~\ref{ass:mono}, $s_{\infty m}^*$ is monotone, so its
inverse is well-defined, and the one-dimensional change-of-variables
formula gives
\begin{equation}\label{eq:pf-cov-poly-proof}
g(d_j)
=Pf(s_{\infty m}^*)(d_j)
=\frac{f(y_j)}{(s_{\infty m}^*)'(y_j)},
\qquad j=1,\dots,m.
\end{equation}

We now construct a polynomial approximation that is feasible for the
$m$ discretized constraints. The idea is as follows. We first choose a
polynomial $q_n$ that approximates $s_{\infty m}^*$ and its derivative
globally on $[c,d]$. This polynomial need not satisfy the $m$
discrete constraints exactly. We then amend $q_n$ by adding a small
Hermite interpolation correction $r_n$, chosen so that the corrected
polynomial
\[
\hat s_{nm}:=q_n+r_n
\]
matches both the values and the derivatives of $s_{\infty m}^*$ at the
special nodes $y_1,\dots,y_m$. This guarantees that the discrete
constraints are satisfied exactly, while the correction remains small
because the value and derivative mismatches of $q_n$ at those nodes go
to zero.

By the Weierstrass approximation theorem
\citep[Theorem~7.26]{rudin1976principles}, there exists a sequence of
polynomials $p_n\in\mathbb P_{n-2}$ such that
\[
\|p_n-(s_{\infty m}^*)'\|_\infty\to 0.
\]
Define
\[
q_n(x):=s_{\infty m}^*(c)+\int_c^x p_n(t)\,dt.
\]
Then $q_n\in\mathbb P_{n-1}$, $q_n'=p_n$, and
\[
\|q_n-s_{\infty m}^*\|_\infty
=
\sup_{x\in[c,d]}
\left|\int_c^x\bigl(p_n(t)-(s_{\infty m}^*)'(t)\bigr)\,dt\right|
\le
(d-c)\,\|p_n-(s_{\infty m}^*)'\|_\infty \to 0.
\]
Hence
\[
\|q_n-s_{\infty m}^*\|_\infty
+
\|q_n'-(s_{\infty m}^*)'\|_\infty
\to 0.
\]

We now correct $q_n$ at the finitely many nodes
$y_1,\dots,y_m$. By the classical first-order Hermite interpolation
theorem, for any prescribed values and first derivatives at $m$
distinct nodes, there exists a unique interpolating polynomial in
$\mathbb P_{2m-1}$; equivalently, there exist basis polynomials
$\{A_j,B_j\}_{j=1}^m\subset \mathbb P_{2m-1}$ such that ~\cite[Section 8.2]{hildebrand1987introduction}
\[
A_j(y_i)=\delta_{ij},\quad A_j'(y_i)=0,
\qquad
B_j(y_i)=0,\quad B_j'(y_i)=\delta_{ij},
\qquad i,j=1,\dots,m,
\]
where $\delta_{ij}=1$ if $i=j$ and $\delta_{ij}=0$ otherwise. 

The role of these basis polynomials are: $A_j$ corrects only the
function value at the node $y_j$, while $B_j$ corrects only the
derivative at the node $y_j$. Therefore, if we define
\[
r_n(y)
:=
\sum_{j=1}^m\bigl(d_j-q_n(y_j)\bigr)A_j(y)
+
\sum_{j=1}^m\bigl((s_{\infty m}^*)'(y_j)-q_n'(y_j)\bigr)B_j(y),
\]
and set
\[
\hat s_{nm}:=q_n+r_n,
\]
then $\hat s_{nm}$ agrees with $s_{\infty m}^*$ in both value and
slope at each node $y_j$.

Indeed, by the defining properties of the Hermite basis,
\begin{align*}
\hat s_{nm}(y_i)
&=
q_n(y_i)
+\sum_{j=1}^m\bigl(d_j-q_n(y_j)\bigr)A_j(y_i)
+\sum_{j=1}^m\bigl((s_{\infty m}^*)'(y_j)-q_n'(y_j)\bigr)B_j(y_i)\\
&=
q_n(y_i)+\bigl(d_i-q_n(y_i)\bigr)=d_i,
\end{align*}
and similarly
\begin{align*}
\hat s_{nm}'(y_i)
&=
q_n'(y_i)
+\sum_{j=1}^m\bigl(d_j-q_n(y_j)\bigr)A_j'(y_i)
+\sum_{j=1}^m\bigl((s_{\infty m}^*)'(y_j)-q_n'(y_j)\bigr)B_j'(y_i)\\
&=
q_n'(y_i)+\bigl((s_{\infty m}^*)'(y_i)-q_n'(y_i)\bigr)
=(s_{\infty m}^*)'(y_i).
\end{align*}
Thus
\begin{equation}\label{eq:hermite-match-poly-proof}
\hat s_{nm}(y_j)=d_j=s_{\infty m}^*(y_j),
\qquad
\hat s_{nm}'(y_j)=(s_{\infty m}^*)'(y_j),
\qquad j=1,\dots,m.
\end{equation}

Since $q_n\in\mathbb P_{n-1}$ and $r_n\in\mathbb P_{2m-1}$, we have
$\hat s_{nm}\in\mathbb P_{n-1}$ whenever $n\ge 2m$.

Next we quantify the size of the correction $r_n$. For fixed $m$ and
fixed nodes $\{y_j\}_{j=1}^m$, the Hermite basis polynomials
$A_j,B_j$ are fixed continuous functions on the compact interval
$[c,d]$. Hence their sup norms, and the sup norms of their
derivatives, are finite. Define
\begin{equation}\label{eq:C_m}
C_m
:=
\sum_{j=1}^m
\Bigl(
\|A_j\|_\infty+\|A_j'\|_\infty+\|B_j\|_\infty+\|B_j'\|_\infty
\Bigr)
<\infty.    
\end{equation}

Using the definition of $r_n$ and the triangle inequality, we obtain
\begin{align*}
\|r_n\|_\infty+\|r_n'\|_\infty
&\le
C_m\left(
\max_{1\le j\le m}|d_j-q_n(y_j)|
+
\max_{1\le j\le m}|(s_{\infty m}^*)'(y_j)-q_n'(y_j)|
\right)\\
&=
C_m\left(
\max_{1\le j\le m}|s_{\infty m}^*(y_j)-q_n(y_j)|
+
\max_{1\le j\le m}|(s_{\infty m}^*)'(y_j)-q_n'(y_j)|
\right)\to 0.
\end{align*}
Therefore
\[
\|r_n\|_\infty+\|r_n'\|_\infty\to 0.
\]
In particular,
\[
\|\hat s_{nm}-s_{\infty m}^*\|_\infty
\le
\|q_n-s_{\infty m}^*\|_\infty+\|r_n\|_\infty\to 0,
\]
and
\[
\|\hat s_{nm}'-(s_{\infty m}^*)'\|_\infty
\le
\|q_n'-(s_{\infty m}^*)'\|_\infty+\|r_n'\|_\infty\to 0.
\]

By Assumption~\ref{ass:mono}, $s_{\infty m}^*\in C^1([c,d])$ and is
strictly monotone. Hence $(s_{\infty m}^*)'$ has a fixed sign on
$[c,d]$. Without loss of generality, we treat the increasing case, so
that
\[
(s_{\infty m}^*)'(x)>0,\qquad x\in[c,d].
\]
Since $\hat s_{nm}'\to (s_{\infty m}^*)'$ uniformly on $[c,d]$, it
follows that $\hat s_{nm}'>0$ on $[c,d]$ for all sufficiently large
$n$.

Finally, combining \eqref{eq:hermite-match-poly-proof} with
\eqref{eq:pf-cov-poly-proof}, we obtain
\[
Pf(\hat s_{nm})(d_j)
=
\frac{f(y_j)}{\hat s_{nm}'(y_j)}
=
\frac{f(y_j)}{(s_{\infty m}^*)'(y_j)}
=
g(d_j),
\qquad j=1,\dots,m.
\]
Thus $\hat s_{nm}\in\mathcal F_{nm}$.

By Proposition~\ref{prop: obj-continuity} and the optimality of
$s_{nm}^*$ over $\mathcal F_{nm}$,
\[
J(s_{\infty m}^*)
\le
J(s_{nm}^*)
\le
J(\hat s_{nm})
\le
J(s_{\infty m}^*)+K\|\hat s_{nm}-s_{\infty m}^*\|_\infty.
\]
Since $\|\hat s_{nm}-s_{\infty m}^*\|_\infty\to 0$, we conclude that
\[
J(s_{nm}^*)\to J(s_{\infty m}^*)
\qquad\text{for each fixed }m.
\]

To upgrade from pointwise convergence to uniform convergence in $m$,
define the auxiliary continuous functions $l_n,l:[0,1]\to\mathbb R$
exactly as in the proof of Theorem~\ref{thm:cp-conv-to-cpc}. Since
$\mathcal S_n$ is nested, the sequence $n\mapsto J(s_{nm}^*)$ is
non-increasing for each fixed $m$, so $l_n\downarrow l$ pointwise on
$[0,1]$. By construction, $l_n$ and $l$ are continuous on $(0,1]$,
and right-continuous at $0$ by
Theorem~\ref{thm: cp-conv-to-cpt-cpc-conv-to-nlp}%
\ref{thm: cp-conv-to-cpt-cpc-conv-to-nlp1} and
Theorem~\ref{thm: cp-conv-to-cpt-cpc-conv-to-nlp}%
\ref{thm: cp-conv-to-cpt-cpc-conv-to-nlp2}, respectively.
Therefore, by Dini's theorem on the compact interval $[0,1]$,
$l_n\to l$ uniformly. Evaluating at $x=1/m$ yields
\[
J(s_{nm}^*)\xrightarrow[n\to\infty]{}J(s_{\infty m}^*)
\qquad\text{uniformly in }m.
\]
\end{proof}


\subsection{Rate of Convergence}\label{subsec:rate}

In this section, we establish the rate of convergence of
$J(s_{nm}^*)$ to $J(s^*)$.  All results in this section require
Assumptions~\ref{ass:cost-lipschitz}, and the conditions for general
cost functions in Theorem~\ref{thm:philippis} or conditions for
strictly convex translation cost in
Theorem~\ref{thm:gangbo-mccann}.  The rate decomposes along the
triangle inequality
\[
\bigl|J(s_{nm}^*)-J(s^*)\bigr|
\leq
\underbrace{\bigl|J(s_{nm}^*)-J(s_{\infty m}^*)\bigr|}_{\text{approximation error in }n}
+
\underbrace{\bigl|J(s_{\infty m}^*)-J(s^*)\bigr|}_{\text{discretization error in }m}.
\]
The second term is bounded by $O(1/m)$ in
Proposition~\ref{prop:rate2} and is independent of the choice of
approximation space $\mathcal{S}_n$.  The first term depends on the choice of $\mathcal{S}_n$.  For
continuous piecewise linear functions, we establish an $O(1/n)$ bound
with a constant that is uniform in $m$ in Proposition~\ref{prop:rate1}, yielding the combined rate
$O(1/n+1/m)$ in Theorem~\ref{thm:rate-main}.  For polynomial
subspaces, the analogous bound carries an $m$-dependent constant
(see Remark~\ref{rem:poly-rate}), so the combined rate result is
stated only for piecewise linear approximation spaces.

\begin{proposition}\label{prop:rate1}
 Under Assumptions~\ref{ass:domain-support},~\ref{ass:cost-lipschitz}, the conditions of either Theorem~\ref{thm:philippis} or Theorem~\ref{thm:gangbo-mccann}, and assuming $\mathcal{S}_n$ is the space of continuous piecewise linear functions on $\mathrm{supp}(f)$, there exists a constant $W_1>0$ independent of $m$, such that for every fixed $m\ge 1$ and all $n\ge 1$,
    \[
    \bigl|J(s_{nm}^*)-J(s_{\infty m}^*)\bigr|\le \frac{W_1}{n}.
    \]
    Equivalently, for each fixed $m$, $\bigl|J(s_{nm}^*)-J(s_{\infty m}^*)\bigr|=O(1/n)$ with constants that are independent of $m$.
\end{proposition}
\begin{proof}
First, let us note that for a fixed $m>1$, $\mathcal{F}_{nm}\subseteq\mathcal{F}_{\infty m}$. Therefore, we have
\begin{align*}
    J(s_{\infty m}^*)\leq J(s_{nm}^*)\leq J(s_{nm}),
\end{align*}
for any $s_{nm}\in\mathcal{F}_{n m}$. Hence, the proof strategy is to construct a sequence $\{\hat{s}_{nm}\}_{n\geq 1}\in\mathcal{F}_{nm}$, such that $J(\hat{s}_{nm})\rightarrow J(s_{\infty m}^*)$ with linear rate, which ensures that $J(s_{nm}^*)\rightarrow J(s_{\infty m}^*)$ with at least linear rate. Consider the even partition $\{x_1, x_2\cdots, x_n\}$ of $\mathrm{supp}(f)=[c,d]$ where $t_n:=\frac{d-c}{n}$ is the length between two consecutive points in the partition. Here we fix by convention that $x_1 = c$, $x_n = d$. For any $x\in \mathrm{supp}(f)$, we know $x_i\leq x\leq x_{i+1}$ for some $1\leq i\leq m-1$. Consider $\hat{s}_{nm}$ that:
\begin{align*}
    \hat{s}_{nm}(x) = s_{\infty m}^*(x_{i+1})\frac{x-x_i}{t_n} + s_{\infty m}^*(x_i)\frac{x_{i+1}-x}{t_n}.
\end{align*}

This choice of $\hat{s}_{nm}$ is essentially a $n$-pieces piecewise linear approximation of $s_{\infty m}^*$ such that it is in space $\mathcal{F}_{nm}$. Note that both $\hat{s}_{nm}$ and $s_{\infty m}^*$ are Lipschitz, denote the larger Lipschitz constant of the two as $L$.
Now, we have 
\begin{align*}
    &J(s_{nm}^*)-J(s_{\infty m}^*)\\
    \leq&|J(\hat{s}_{nm})-J(s_{\infty m}^*)|\\
    \leq& K\|\hat{s}_{nm}-s_{\infty m}^*\|_{\infty}\\
    \leq& K\sup_{x\in \mathrm{supp}(f)}\left(\frac{L(x-x_i)}{t_n}\left(x-x_i\right) + \frac{L(x_{i+1}-x)}{t_n}\left(x_{i+1}-x\right)\right)\\
    =& K\sup_{x\in \mathrm{supp}(f)}\left(\frac{L}{t_n}((x-x_i)^2+(x_{i+1}-x)^2)\right)\\
    \leq& K\sup_{x\in \mathrm{supp}(f)}\left(\frac{L}{t_n}2t_n^2\right)\\
    \leq& 2KLt_n\\
    \leq& \frac{2KL(c-d)}{n},
\end{align*}
for all $m>0$, where the numerator is a constant.
\end{proof}

\begin{remark}\label{rem:poly-rate}
In the proof of Theorem~\ref{thm:cp-conv-to-cpc-poly}, the bound
$\|\hat{s}_{nm}-s_{\infty m}^*\|_\infty\to 0$ depends on the
constant $C_m$ in~\eqref{eq:C_m}, which depends on $m$.  This
prevents a convergence rate in $n$ that is uniform in $m$, in
contrast with the piecewise linear case
(Proposition~\ref{prop:rate1}), where the interpolation error is
intrinsically local and independent of~$m$.  For this reason, the
combined rate $O(1/n+1/m)$ in Theorem~\ref{thm:rate-main} is stated
only for piecewise linear approximation spaces.
\end{remark}

\begin{proposition}\label{prop:rate2}
Under Assumptions~\ref{ass:domain-support},~\ref{ass:cost-lipschitz} and the conditions of either Theorem~\ref{thm:philippis} or Theorem~\ref{thm:gangbo-mccann}, there exists a constant $W_2>0$ such that for all $m\ge 1$,
\[
\bigl|J(s_{\infty m}^*)-J(s^*)\bigr|\le \frac{W_2}{m}.
\]
Equivalently, $\bigl|J(s_{\infty m}^*)-J(s^*)\bigr|=O(1/m)$.
\end{proposition}

\begin{proof}
First, note that we have $\mathcal{F}\subseteq\cdots\subseteq\mathcal{F}_{\infty m+1}\subseteq\mathcal{F}_{\infty m}$. Therefore, for any $\hat{s}\in\mathcal{F}$ we have 
\begin{align*}
    J(s_{\infty m}^*)\leq J(s_{\infty m+1}^*)\leq\cdots\leq J(\hat{s}).
\end{align*}  
We know that $J(s_{\infty m}^*)\xrightarrow{m\rightarrow\infty}J(s^*)$ monotonically from below, therefore $\{J(s_{\infty m}^*)\}_{m\geq 1}$ is a Cauchy sequence. First, observe that we have
\begin{align*}
    0\leq J(s_{\infty m+1}^*) - J(s_{\infty m}^*)\leq J(s_{\infty m+1}) - J(s_{\infty m}^*),
\end{align*}
for any $s_{\infty m+1}\in\mathcal{F}_{\infty m+1}$. Moreover,
\begin{align*}
    &|J(s^*)-J(s_{\infty m}^*)|\\
    =&J(s^*)-J(s_{\infty m}^*)\\
    =&\lim_{i\rightarrow\infty}J(s_{\infty i}^*)-J(s_{\infty m}^*)\\
    =&\lim_{i\rightarrow\infty}(-J(s_{\infty m}^*)+J(s_{\infty m+1}^*)-J(s_{\infty m+1}^*)+\\
    &\cdots+J(s_{\infty i-1}^*)-J(s_{\infty i-1}^*)+J(s_{\infty i}^*))\\
    =&\lim_{j\rightarrow\infty}\sum_{i=m+1}^{j} J(s_{\infty i}^*)-J(s_{\infty i-1}^*).
\end{align*}
The last limit exists because we proved the second equality in Theorem~\ref{thm: cp-conv-to-cpt-cpc-conv-to-nlp}(ii). Therefore, we can write
\begin{align*}
    &|J(s^*)-J(s_{\infty m}^*)|\\
    =&\lim_{j\rightarrow\infty}\sum_{i=m+1}^{j} J(s_{\infty i}^*)-J(s_{\infty i-1}^*)\\
    =&\sum_{i=m+1}^{\infty} J(s_{\infty i}^*)-J(s_{\infty i-1}^*)\\
    \leq&\sum_{i=m+1}^{\infty} J(\hat{s}_{\infty i})-J(s_{\infty i-1}^*)
\end{align*}
for any sequence $\{\hat{s}_{\infty i}\}_{i\geq 1}\subseteq \cup_{i\geq 1}\mathcal{F}_{\infty i}$.

Our proof strategy, therefore, is to
construct a sequence $\{\hat{s}_{\infty i}\}_{i\geq 1}\subseteq \cup_{i\geq 1}\mathcal{F}_{\infty i}$ such that 
\[
\left\{\sum_{i=m+1}^{\infty}J(\hat{s}_{\infty i})-J(s_{\infty i-1}^*)\right\}_{m\geq 1}
\]
goes to $0$ as $m\rightarrow\infty$ at a certain rate.
Now, let us construct such a sequence. Recall the sequence of partitions in the proof of Theorem~\ref{thm: feasible-set-conv} where $\{\mathcal{P}_i:=\{d_1, d_2, \cdots,d_i\}\}$ is a sequence of partitions of $\mathrm{supp}(g)$ such that $\mathcal{P}_1\subseteq\mathcal{P}_2\subseteq\cdots$, and $\lim_{i\rightarrow\infty}\mathcal{P}_i=\mathcal{D}$ where $\mathcal{D}\subset \mathrm{supp}(g)$ is dense. We assume without loss of generality that each $\mathcal{P}_i$ is a uniform partition such that $d_{n+1}-d_n=h_i,\quad\forall 1\leq n\leq i-1$ with $d_1=a$, $d_2=b$ given that $\mathrm{supp}(g) = [a,b]$.


Consider two consecutive terms $\mathcal{P}_i$ and $\mathcal{P}_{i+1}$ in the sequence of partitions on which the density constraint must be satisfied, let $D:= \mathcal{P}_{i+1}\setminus\mathcal{P}_i$. First, note that for $s_{\infty i}^*$, we have 
\begin{align*}
Pf(s_{\infty i}^*)(d_1)&=g(d_1),\\
&\vdots\\
Pf(s_{\infty i}^*)(d_i)&=g(d_i).
\end{align*}
Now let $\hat{s}_{\infty i+1}=s_{\infty i}^*$ so we have
\begin{align*}
Pf(\hat{s}_{\infty i+1})(d_1)&=g(d_1),\\
&\vdots\\
Pf(\hat{s}_{\infty i+1})(d_i)&=g(d_i).
\end{align*}
Now we want to perturb $\hat{s}_{\infty i+1}$ such that the density constraint is satisfied at all points in $\mathcal{D}$ as well as the other grid points where the constraint is already satisfied. Recall that by Leibniz rule, we have $Pf(\hat{s}_{\infty i+1})(d)=f((\hat{s}_{\infty i+1})^{-1}(d))(\hat{s}_{\infty i+1}^{-1})'(d)$. Note that for any $d\in D$ it must be true that $d\in(d_n, d_{n+1})$ for some $1\leq n\leq i-1$. Now, we redefine $\hat{s}_{\infty i+1}$ in a small neighborhood around $d$ as a piecewise linear function with 3 pieces. Since $\hat{s}_{\infty i+1}$ is a bijection, we redefine $\hat{s}_{\infty i+1}$ by redefining $\hat{s}^{-1}_{\infty i+1}$ for simplicity. 
\begin{figure}[hbt!]
    \centering
    \includegraphics[scale=.5]{graphs/construction.png}
    \caption{}
    \label{fig:2d}
\end{figure}
Consider
\begin{align*}
    \hat{s}^{-1}_{\infty i+1}(y)= 
\begin{cases}
    s_{\infty i}^{*-1}(y),& \text{if } y\in\{d-\frac{1}{i}(d-d_n),d,d+\frac{1}{i}(d_{n+1}-d)\}\\
    s_{\infty i}^{*-1}(y)-\frac{1}{2i}(d-d_{n})\frac{g(d)}{f(s_{\infty i}^{*-1}(d))},& \text{if } y=d-\frac{1}{2i}(d-d_n)\\
    s_{\infty i}^{*-1}(y)+\frac{1}{2i}(d_{n+1}-d)\frac{g(d)}{f(s_{\infty i}^{*-1}(d))},& \text{if } y=d+\frac{1}{2i}(d_{n+1}-d)\\
    \text{linear interpolated value}, & \text{else if } y\in(d-\frac{1}{i}(d-d_n),d+\frac{1}{i}(d_{n+1}-d)).
\end{cases}
\end{align*}
$\frac{g(d)}{f(s_{\infty i}^{*-1}(d))}$ is the slope we impose on the modified $\hat{s}^{-1}_{\infty i+1}$ around $d$ while maintaining $\hat{s}^{-1}_{\infty i+1}(d) = s_{\infty i}^{*-1}(d)$. Therefore, we have
\begin{align*}
    Pf(\hat{s}_{\infty i+1})(d)&=f((\hat{s}_{\infty i+1})^{-1}(d))(\hat{s}_{\infty i+1}^{-1})'(d)\\
    &=f(s_{\infty i}^{*-1}(d))\frac{g(d)}{f(s_{\infty i}^{*-1}(d))}\\
    &=g(d),
\end{align*}
and we can confirm that the constraint is satisfied on $d$ while the constraint at other locations is not disrupted. By construction, $\|s_{\infty i}^{*}-\hat{s}_{\infty i+1}\|_{\infty}\leq \frac{1}{i}(h_i)=\frac{b-a}{i^2}$ and 
\begin{align*}
    &|J(s^*)-J(s_{\infty m}^*)|\\
    \leq&\sum_{i=m+1}^{\infty} J(\hat{s}_{\infty i})-J(s_{\infty i-1}^*)\\
    \leq&\sum_{i=m+1}^{\infty}K \|\hat{s}_{\infty i}-s_{\infty i-1}^{*}\|_{\infty}\\
    \leq&K(b-a)\sum_{i=m+1}^{\infty}\frac{1}{i^2}\\
    \sim&\frac{K(b-a)}{m}\quad\text{for large } m.
\end{align*}
Where the numerator is a constant.
\end{proof}

\begin{theorem}\label{thm:rate-main}
Under Assumptions~\ref{ass:domain-support}--\ref{ass:cost-lipschitz}, \ref{ass:mono}, the conditions of either Theorem~\ref{thm:philippis} or Theorem~\ref{thm:gangbo-mccann}, and assuming $\mathcal{S}_n$ is the space of continuous piecewise linear functions on $\mathrm{supp}(f)$, there exist constants $W_1,W_2>0$, independent of $n$ and $m$, such that for all $n,m\ge 1$,
\[
\bigl|J(s_{nm}^*)-J(s^*)\bigr|
\;\le\;
\frac{W_1}{n}+\frac{W_2}{m}
\;=\;
O\!\left(\frac1n+\frac1m\right).
\]
In particular, $J(s_{nm}^*)\to J(s^*)$ as $n,m\to\infty$.
\end{theorem}

\begin{proof}
By the triangle inequality,
\[
\bigl|J(s_{nm}^*)-J(s^*)\bigr|
\le
\bigl|J(s_{nm}^*)-J(s_{\infty m}^*)\bigr|
+
\bigl|J(s_{\infty m}^*)-J(s^*)\bigr|.
\]
By Proposition~\ref{prop:rate1} and Proposition~\ref{prop:rate2}, the two terms on the right admit upper bounds of order $O(1/n)$ and $O(1/m)$ respectively, with constants $W_1$, $W_2$. Combining them yields the result.
\end{proof}


For the computational problem (CP), where the constraint is discretized to $m$ grid points, the local geometry of feasibility depends on the relationship between $n$ and $m$. If $n = m$ and the constraint Jacobian is invertible, feasible points remain isolated. However, in the standard approximation regime where $n > m$, Proposition~\ref{prop:discrete-LICQ} guarantees LICQ for appropriate grid placement, making the feasible set $\mathcal{F}_{nm}$ locally an $(n - m)$-dimensional manifold. In this case, feasible points are not isolated, and first-order feasible directions exist. The difficulty is that the manifold $\mathcal{F}_{nm}$ is defined by a complex nonlinear system. Moving along it while strictly maintaining feasibility requires solving a nonlinear root-finding problem at every continuous step. Sequential quadratic programming avoids this computational bottleneck by linearizing the pushforward constraint at each iterate, approximating the manifold locally by its tangent space and reducing each step to a structured quadratic program with linear constraints. This linearization is what makes SQP computationally viable and well-suited to the geometry of (CP).

To summarize, Section~\ref{sec:cp} has established three things. First, the optimal objective value of the computational problem (CP) converges to the optimal value of the infinite-dimensional problem (P) at rate $O(1/n + 1/m)$ as the approximation dimension $n$ and constraint grid size $m$ are refined (Theorem~\ref{thm:rate-main}). Second, under LICQ, the feasible set $\mathcal{F}_{nm}$ of (CP) is locally an $(n-m)$-dimensional manifold, providing a well-defined geometric setting for constrained optimization. Third, the nonlinear geometry of the feasible sets motivates the use of a method that linearizes the pushforward constraint at each iterate rather than attempting to enforce it exactly. Motivated by these geometric insights, Section~\ref{sec:SQP} formally develops the Sequential Quadratic Programming framework, assembling our variational derivatives into a tractable algorithm with convergence guarantees.

\section{Sequential Quadratic Programming}\label{sec:SQP}
\subsection{The Functional SQP}\label{subsec:SQP}
As discussed, the Monge formulation yields a constrained nonconvex problem, while local optima are still characterized by the standard KKT conditions. From a computational standpoint, any implementable method must restrict the search to a finite-dimensional subspace of $\mathcal{S}$. In this setting, feasibility can be poorly conditioned, so it is natural to use a method that allows infeasible iterates while controlling the constraint residual. Sequential quadratic programming (SQP) is well suited for this purpose. Moreover, the Fr\'echet derivatives derived in Section~\ref{sec:deriv} provide the required local structures: the first and second variations of the objective to form the quadratic approximation, and the first variation of the pushforward constraint to build its linear approximation.
In this section we present the SQP method by defining the local quadratic subproblem (QP\textsuperscript{k}) for (P) and its discrete analogue (CPQP\textsuperscript{k}) for (CP), and we summarize the corresponding convergence properties in both the infinite-dimensional and computational settings.

We begin with the infinite-dimensional setting. Given the first- and second-order information, we now consider the SQP procedure for (P). Let $s_k$ represent the transport map computed at the $k$th step of the procedure and let $h\in\mathcal{S}$ be a perturbation direction. We approximate $J(s_k+h)$ up to second order by $\Tilde{J}(s_k+h)$, where for any $h\in\mathcal{S}$,
\begin{align*}
    \Tilde{J}(s_k+ h) := J(s_k) + J_{s_k}(h)+ \frac{1}{2}H_{s_k}(h)(h),
\end{align*}
and $J_{s_k}(h)$ and $H_{s_k}(h)(h)$ are the first and second Fr\'echet derivatives defined in Theorems~\ref{thm: J1deriv} and~\ref{thm: J2deriv}.

Next, we approximate the constraint up to first order. Recall the original constraint $Pf(s_k)-g = 0$, and define
\begin{align*}
    \Tilde{P}f(s_k+h)(x) := Pf(s_k)(x) + Pf_{s_k}(h)(x),
\end{align*}
which is a first-order linearization of the Perron--Frobenius operator $Pf(s_k)$. Here, $Pf_{s_k}(h)(x)$ is the Fr\'echet derivative of the F--P operator as derived in Theorem~\ref{thm: PF1deriv}. We now propose to solve the following quadratic problem:
\begin{align*}
      &\min_{h\in\Omega}\Tilde{J}(s_k + h)\tag{QP\textsuperscript{k}}\\
    \text{s.t.}\quad &\Tilde{P}f(s_k+h)-g = 0.  
\end{align*}

Consider the Lagrangian of (QP\textsuperscript{k}):
\begin{align*}
   \mathcal{L}(h, z_k^*)
   &= \Tilde{J}(s_k+ h) + \big\langle z_k^*,\,\Tilde{P}f(s_k+h)-g\big\rangle_{\mathcal{Z}^*,\mathcal{Z}}\\
   &=  J(s_k) + J_{s_k}(h)+ \frac{1}{2}H_{s_k}(h)(h)
    + \big\langle z_k^*,\,Pf(s_k) + Pf_{s_k}(h)-g\big\rangle_{\mathcal{Z}^*,\mathcal{Z}},
\end{align*}
where $\mathcal{Z}:=L^1(\mathcal{X})$ and $\mathcal{Z}^*\cong L^\infty(\mathcal{X})$. In particular, when $z_k^*$ is represented by a function $\lambda_k\in L^\infty(\mathcal{X})$ under $\mathcal{Z}^*\cong L^\infty(\mathcal{X})$, the Lagrangian admits the integral form
\begin{align*}
   \mathcal{L}(h, \lambda_k)
   &= \Tilde{J}(s_k+ h) + \int_{\mathcal{X}} \lambda_k(x)\big(\Tilde{P}f(s_k+h)-g\big)(x)dx\\
   &=  J(s_k) + J_{s_k}(h)+ \frac{1}{2}H_{s_k}(h)(h)
    + \int_{\mathcal{X}} \lambda_k(x)\big(Pf(s_k)(x) + Pf_{s_k}(h)(x)-g(x)\big)dx,
\end{align*}
where $\lambda_k$ is the Lagrange multiplier at the current step.

To derive the first-order optimality conditions for the SQP subproblem, we compute the directional derivative of each term in $\Tilde{\mathcal{L}}$ with respect to $h$ along a direction $u\in\mathcal{S}$. The first term $J(s_k)$ is independent of $h$, so its directional derivative is zero. Since $J_{s_k}(\cdot)$ is linear, we have $J_{s_k,h}(h)(u)=J_{s_k}(u)$. For the quadratic term, using the bilinear form representation
\[
H_{s_k}(a)(b)=\int_{\mathcal{X}} a(x)b(x)\,D^2_{yy}c\big(x,s_k(x)\big)\,f(x)dx,
\]
we obtain
\begin{align*}
    &\left(\frac{1}{2}H_{s_k}(h)(h)\right)_h(u)
    =\lim_{\epsilon\rightarrow 0}\frac{\frac{1}{2}H_{s_k}(h+\epsilon u)(h+\epsilon u)-\frac{1}{2}H_{s_k}(h)(h)}{\epsilon}\\
    &=\lim_{\epsilon\rightarrow 0}\frac{1}{\epsilon}\int_{\mathcal{X}}\Big(h(x)\epsilon u(x)+\tfrac{1}{2}\epsilon^2u(x)^2\Big)\,D^2_{yy}c\big(x,s_k(x)\big)\,f(x)dx\\
    &=\int_{\mathcal{X}} h(x)u(x)\,D^2_{yy}c\big(x,s_k(x)\big)\,f(x)dx\\
    &=H_{s_k}(h)(u).
\end{align*}

Finally, for the constraint term, since $Pf_{s_k}(\cdot)$ is linear, we have $Pf_{s_k,h}(h)(u)=Pf_{s_k}(u)$ and hence
\[
\left(\int_{\mathcal{X}} \lambda_k(x)\big(Pf(s_k)(x)+Pf_{s_k}(h)(x)-g(x)\big)dx\right)_h(u)
=\int_{\mathcal{X}} \lambda_k(x)\,Pf_{s_k}(u)(x)dx.
\]

Putting the pieces together, the first-order optimality conditions for (QP\textsuperscript{k}) are:
\begin{align*}
     J_{s_k}(\cdot) + H_{s_k}(h)(\cdot) + \int_{\mathcal{X}} \lambda_k(x)\, Pf_{s_k}(\cdot)(x)dx &= 0,\\
     \Tilde{P}f(s_k+h)(x)-g(x) &= 0,\qquad \forall x\in\mathcal{X}.
\end{align*}
Let $(h_k, z_k^*)$ be the solution of the first-order optimality conditions above, where $z_k^*$ is represented by $\lambda_k\in L^\infty(\mathcal{X})$ under $\mathcal{Z}^*\cong L^\infty(\mathcal{X})$. We then update
\begin{align*}
    s_{k+1} = s_k + \alpha_k h_k,
\end{align*}
where $\alpha_k$ is a suitable step size.



\subsection{Local Convergence}\label{subsec:loc_conv}

In this section, we establish the local convergence of the SQP procedure by verifying the assumptions of Alt~\cite{alt1990lagrange} at a local solution $s_0$ of (P).Recall that our constraint space is
\[\mathcal{Z}_0=\Big\{ \eta \in L^1(\mathcal{X}) : \int_{\mathcal{X}}\eta(x)\,dx = 0 \Big\},
\]
whose dual is $\mathcal{Z}_0^* \cong L^\infty(\mathcal{X})/\mathbb R$. By identifying any abstract multiplier $z^*\in \mathcal{Z}_0^*$ with a functional representative $\lambda\in L^\infty(\mathcal{X})$ (determined up to an additive constant), we can express the Lagrangian $\mathcal{L}(s,\lambda)$ directly in integral form:
\[
\mathcal{L}(s,\lambda) := J(s)+\int_{\mathcal{X}}\lambda(x)\mathcal{H}(s)(x)\,dx,
\]
where $\mathcal{H}(s):=Pf(s)-g\in \mathcal{Z}_0$.Accordingly, we will henceforth refer to a local KKT pair by its representative $(s_0,\lambda_0)$, which satisfies primal feasibility, $\mathcal{H}(s_0)=0_{\mathcal{Z}_0}$, and the stationarity condition:
\[
J_{s_0}(h)+\int_{\mathcal{X}}\lambda_0(x)\,\mathcal{H}_{s_0}(h)(x)\,dx=0, \qquad \forall h\in\mathcal{S}.
\]

\begin{proposition}\label{prop:verify-alt-condition}Fix a local KKT pair $(s_0,\lambda_0)$ of (P).Under Assumptions~\ref{ass:domain-support}, \ref{ass:density-cont}, \ref{ass:cost-second}, \ref{ass:mono}, \ref{ass:hardy}, assuming $\lambda_0 \in W^{1,\infty}(\mathcal{X})$, and the conditions of Theorem~\ref{thm: PF1deriv} and Theorem~\ref{thm: PF2deriv}, the following conditions hold at $s_0$.
\begin{itemize}
\item[(B1)] The objective functional $J:\mathcal{S}\to\mathbb R$ and the Perron--Frobenius operator $Pf:\mathcal{S}\to\mathcal{Z}$ are twice continuously Fr\'echet differentiable on a neighborhood $N_1(s_0)\subset \mathcal{S}$.
\item[(B2)] The second-variation operators $J'':\mathcal{S}\to \mathcal{S}^{**}$ and $Pf'':\mathcal{S}\to \mathcal{B}(\mathcal{S},\mathcal{B}\big(\mathcal{S},\mathcal{Z}_0)\big)$ are Lipschitz continuous on $N_1(s_0)\subset \mathcal{S}$.
\item [(B3)] $0_{\mathcal{Z}_0} \in \mathrm{int}\Big\{Pf_{s_0}(h)\mid h\in\mathcal{S}\Big\}$ in $\mathcal{Z}_0$, where $0_{\mathcal{Z}_0}$ is the zero \item [(B4)] There exists an $\alpha>0$ such that $\mathcal{L}_{s_0 s_0}(h)(h)\geq\alpha\|h\|^2$
for all $h\in\mathcal{S}$ with $Pf_{s_0}(h)=0_{\mathcal{Z}_0}$.
\end{itemize}
\end{proposition}
\begin{proof}
\noindent Proof of (B1) and (B2):
Since \(s_0\in\Omega\) and \(\Omega\) is open in \(W^{1,\infty}(\mathcal{X})\) (see Section~\ref{subsec:notation}), there exists \(\epsilon>0\) such that
\[
B^{W^{1,\infty}}_\epsilon(s_0)\subset \Omega.
\]
We therefore work on the neighborhood
\[
N_1(s_0):=\mathcal{S}\cap B^{W^{1,\infty}}_\epsilon(s_0).
\]
By construction, every \(s\in N_1(s_0)\) remains in the admissible class, hence remains monotone and nondegenerate. In particular, \(Pf(s)\) is well defined throughout \(N_1(s_0)\) (see Lemma~\ref{lem:pf-continuous-version}), and Theorem~\ref{thm: PF1deriv}, and Theorem~\ref{thm: PF2deriv} apply on this neighborhood.

We first verify the required properties for \(J\). For any \(s_1,s_2\in N_1(s_0)\),
\begin{align*}
&\|J'(s_1)-J'(s_2)\|_{\mathcal{S}^*}\\
=& \sup_{\substack{h\in\mathcal{S}\\ \|h\|=1}}
\left|
\int_{\mathcal{X}} h(x)\big(D_yc(x,s_2(x))-D_yc(x,s_1(x))\big)f(x)\,dx
\right|\\
\leq&
\sup_{\substack{h\in\mathcal{S}\\ \|h\|=1}}
\int_{\mathcal{X}} |h(x)|\,L\,\|s_1-s_2\|_\infty f(x)\,dx\\
\leq& L\|s_1-s_2\|_\infty,
\end{align*}
where \(L\) is the Lipschitz constant of \(D_yc(x,\cdot)\) from Assumption~\ref{ass:cost-second}. Likewise,
\begin{align*}
&\|J''(s_1)-J''(s_2)\|_{\mathcal{S}^{**}}\\
=& \sup_{\substack{h,\tilde h\in\mathcal{S}\\ \|h\|=\|\tilde h\|=1}}
\left|
\int_{\mathcal{X}} h(x)\tilde h(x)\big(D_{yy}c(x,s_1(x))-D_{yy}c(x,s_2(x))\big)f(x)\,dx
\right|\\
\leq&
\sup_{\substack{h,\tilde h\in\mathcal{S}\\ \|h\|=\|\tilde h\|=1}}
\int_{\mathcal{X}} |h(x)|\,|\tilde h(x)|\,M\,\|s_1-s_2\|_\infty f(x)\,dx\\
\leq& M\|s_1-s_2\|_\infty,
\end{align*}
where \(M\) is the Lipschitz constant of \(D_{yy}c(x,\cdot)\). Hence \(J\) is twice continuously Fr\'echet differentiable on \(N_1(s_0)\), and both \(J'\) and \(J''\) are Lipschitz continuous there.

For the Perron--Frobenius operator, the corresponding properties on \(N_1(s_0)\) follow from Theorem~\ref{thm: PF1deriv} and Theorem~\ref{thm: PF2deriv}. More precisely, \(Pf:\mathcal{S}\to\mathcal{Z}\) is twice continuously Fr\'echet differentiable on \(N_1(s_0)\), while
\[
Pf':\mathcal{S}\to\mathcal{B}(\mathcal{S},\mathcal{Z}_0)
\qquad\text{and}\qquad
Pf'':\mathcal{S}\to\mathcal{B}\bigl(\mathcal{S},\mathcal{B}(\mathcal{S},\mathcal{Z}_0)\bigr)
\]
are Lipschitz continuous on \(N_1(s_0)\).

Therefore, conditions (B1) and (B2) hold at \(s_0\).
\end{proof}

\noindent Proof of (B3):
By Theorem~\ref{thm:pf-surjectivity}, under the stated assumptions, the operator $Pf_{s_0}:\mathcal{S}\to\mathcal{Z}_0$ is surjective. By the Open Mapping Theorem \citep[Theorem~2.11]{Rudin1974FunctionalAnalysis}, $Pf_{s_0}$ is an open map, so $0_{\mathcal{Z}_0} = Pf_{s_0}(0)$ is an interior point of $\{Pf_{s_0}(h) \mid h \in \mathcal{S}\}$ in $\mathcal{Z}_0$.

\noindent Proof of (B4):
Recall that the Lagrangian of problem (P) at a local minimizer $(s_0,\lambda_0)$ is
\[
\mathcal{L}(s,\lambda_0)
=
J(s)+\int_{\mathcal{X}}\lambda_0(x)\big(Pf(s)(x)-g(x)\big)\,dx,
\]
where $\lambda_0\in L^\infty(\mathcal{X})$ represents the multiplier.

Its second Fr\'echet derivative with respect to $s$ at $s_0$ is
\[
\mathcal{L}_{s_0s_0}(h)(h)
=
H_{s_0}(h)(h)
+\int_{\mathcal{X}}\lambda_0(x)\,Qf_{s_0}(h)(h)(x)\,dx.
\]

To verify (B4), we only need to consider directions $h\in\mathcal{S}$ satisfying the linearized constraint
\[
Pf_{s_0}(h)=0_{\mathcal{Z}_0}.
\]
By Proposition~\ref{prop:injective}, the operator $Pf_{s_0}$ is injective. Hence $Pf_{s_0}(h)=0_{\mathcal{Z}_0}$ implies $h=0$.
Therefore, the set of feasible critical directions in (B4) is the singleton $\{0\}$.
Evaluating the second derivative in the only such direction gives
\[
\mathcal{L}_{s_0s_0}(0)(0)=0.
\]
Thus, for any $\alpha>0$,
\[
\mathcal{L}_{s_0s_0}(h)(h)
=
0
\ge \alpha \|h\|^2
\qquad
\text{whenever } Pf_{s_0}(h)=0_{\mathcal{Z}_0},
\]
since necessarily $h=0$. Hence condition (B4) holds.
\end{proof}

\begin{theorem}[Local Quadratic Convergence]
\label{thm: local-conv}
Let $s_0$ be a local solution of (P), and let $\lambda_0$ denote the Lagrange multiplier function associated with $s_0$.
Assume that the regularity and stability conditions \emph{(B1)--(B4)} in Proposition~\ref{prop:verify-alt-condition} hold at $s_0$.

Then there exist constants $0<\delta<1$, and $\rho>0$, and a ball
$B_{\rho\delta}\big((s_0,\lambda_0)\big)$ such that $s\in N_1(s_0)$ for every
$(s,\lambda)\in B_{\rho\delta}\big((s_0,\lambda_0)\big)$.
Moreover, for any starting point
$(s_1,\lambda_1) \in B_{\rho\delta}\big((s_0,\lambda_0)\big)$, the SQP method generates a
unique sequence $\{(s_k,\lambda_k)\}_{k\ge 1}$ satisfying, for all $k\ge 2$,
\[
\big\|(s_k,\lambda_k)-(s_0,\lambda_0)\big\|_{\mathcal{S}\times (L^\infty/\mathbb{R})}
:=
\left( \|s_k-s_0\|^2 + \|\lambda_k-\lambda_0\|_{L^\infty(\mathcal{X})/\mathbb{R}}^2 \right)^{1/2}
\;\le\;
\rho\,\delta^{2^{k-1}}.
\]
In particular, the iterates converge \emph{quadratically} to $(s_0,\lambda_0)$.
\end{theorem}

\begin{proof}
Proposition~\ref{prop:verify-alt-condition} verifies (B1)--(B4) for our problem.
Therefore the conclusion follows by applying Alt~(1990, Theorem~3.3) to the SQP mapping written in the $(s,\lambda)$ variables.
\end{proof}

\subsection{Global Convergence}\label{subsec:glo_conv}
Consider the merit functional:
\[
M(s,\eta)=J(s)+\int_{\mathcal{X}}\lambda(x)(Pf(s)(x)-g(x))dx+\frac{\eta}{2}\int_{\mathcal{X}}|Pf(s)(x)-g(x)|dx
\]
where $\eta$ is the penalty parameter of the merit functional.
\begin{theorem}~\label{thm: Merit} 
    For the above-defined merit functional $M(s,\eta)$, under conditions for Theorem~\ref{thm: PF1deriv}, we have the following,
    \begin{enumerate}
        \item $s^*$ is a local min of (P) if and only if it is a local min of the merit functional with sufficiently large $\eta$.
        \item If $s_k$ is not a local min of (P), then $h_k$, the descent direction computed in the $(QP\textsuperscript{k})$ iterations at $s_k$ must also be a descent direction for the merit functional with sufficiently large $\eta$.
    \end{enumerate}
\end{theorem}

\begin{proof}
    Proof of $(1)$: Denote the feasible set of (P) as $\mathcal{F}$. For the forward direction, if $s^*$ is a local min of (P), i.e., $\exists \epsilon$ such that for a neighbourhood $\mathcal{N}_{\epsilon}(s^*)$ of $s^*$, we have 
    \[
    J(s^*)\leq J(s)\quad\forall s\in \mathcal{F}\cap\mathcal{N}_{\epsilon}(s^*).
    \]
    Now suppose that $\exists\Tilde{s}\in \mathcal{N}_{\epsilon}(s^*)$ and $\Tilde{\eta}>0$ such that 
    \[
    M(\Tilde{s}, \eta)<M(s^*, \eta),\quad\forall\eta>\Tilde{\eta}.
    \]
    Since we know that $s^*$ is feasible, we could simplify the former inequality to:
    \[
    J(\Tilde{s})+\frac{\eta}{2}\int_{\mathcal{X}}|Pf(\Tilde{s})(x)-g(x)|dx<J(s^*),\quad\forall\eta>\Tilde{\eta}.
    \]
    Here and throughout the proof of $(1)$, we choose to omit the second term in the merit functional because it will be dominated by the third term for large enough $\eta$. Consider two cases. In the first case, $\Tilde{s}\in\mathcal{F}$. Note that the former inequality implies
    \[
    J(\Tilde{s})<J(s^*),\quad\forall\eta>\Tilde{\eta}
    \]
    which leads to a contradiction. In the second case, $\Tilde{s}\notin\mathcal{F}$, then $\int_{\mathcal{X}}|Pf(\Tilde{s})(x)-g(x)|dx>0$. We could increase $\eta$ until $J(\Tilde{s})+\frac{\eta}{2}\int_{\mathcal{X}}|Pf(\Tilde{s})(x)-g(x)|dx<J(s^*)$ is violated. The second case also leads to a contradiction. Hence we have shown the forward direction of $(1)$. For the backward direction, let's first observe that, if $s^*$ is a local min of the merit functional, then it must be feasible for (P). To see this, first note that the first order necessary condition for the merit functional must be satisfied at $s^*$. Recall the merit functional after omitting the second term:
    \begin{align*}
        M(s,\eta)=J(s)+G(s,\eta),
    \end{align*}
    where $G(s,\eta)=\frac{\eta}{2}\int_{\mathcal{X}}|Pf(s)(x)-g(x)|dx$. The absolute value function $|x|$ is not differentiable at 0, therefore, we use the function $2\epsilon log(1+e^{x/\epsilon})-x-2\epsilon log(2)$ to approximate $|x|$. The derivative of $2\epsilon log(1+e^{x/\epsilon})-x-2\epsilon log(2)$ is $2\frac{e^{x/\epsilon}}{1+e^{x/\epsilon}}-1$. Note that both the approximating function and its derivative converge to $|x|$ and its derivative pointwise as $\epsilon\to 0$ except when $x=0$ since $|x|$ is not differentiable at 0. The approximating function, on the other hand, is differentiable at 0, and its derivative at 0 is 0. Now, applying the approximation of $|x|$ to the merit functional, we obtain the approximating merit functional:
    \begin{align*}
        M^\epsilon(s^*,\eta)(\cdot)&=J(s^*)+G^\epsilon(s^*,\eta)\\
        &= J(s^*)+\frac{\eta}{2}\int_{\mathcal{X}}2\epsilon log(1+e^{Pf(s^*)(x)-g(x)/\epsilon})-(Pf(s^*)(x)-g(x))-2\epsilon log(2)dx.
    \end{align*}
The first variation with respect to $s$ of the approximating merit functional is then
 \begin{align*}
        M_{s^*}^\epsilon(s^*,\eta)(\cdot)&=J_{s^*}(\cdot)+G_{s^*}^\epsilon(s^*,\eta)(\cdot)\\
        &= J_{s^*}(\cdot)+\frac{\eta}{2}\int_{\mathcal{X}}\left(2\frac{e^{(Pf(s^*)(x)-g(x))/\epsilon}}{1+e^{(Pf(s^*)(x)-g(x))/\epsilon}}-1\right)Pf_{s^*}(\cdot)(x)dx.
    \end{align*}
    Note that $G_{s^*}^\epsilon(s^*,\eta)(h)\to G_{s^*}(s^*,\eta)(h)$ for any $h\in\mathcal{S}$ by bounded convergence theorem. 
    Setting the first variation to 0, we get the first necessary condition of the merit functional:
    \begin{align*}
        M_{s^*}^\epsilon(s^*,\eta)(\cdot)=J_{s^*}(\cdot)+\frac{\eta}{2}\int_{\mathcal{X}}\left(2\frac{e^{(Pf(s^*)(x)-g(x))/\epsilon}}{1+e^{(Pf(s^*)(x)-g(x))/\epsilon}}-1\right)Pf_{s^*}(\cdot)(x)dx = 0
    \end{align*}
    where $=0$ means that the first variation is a zero-operator, i.e., $M_{s^*}^\epsilon(s^*,\eta)(h)=0$ for any $h\in\mathcal{S}$.
    In the proof of (B3) in Proposition~\ref{prop:verify-alt-condition}, we showed that $Pf_{s^*}(h)(x)=0$ if and only if $h=0$, that $Pf_{s^*}(\cdot)(x)=0$ is not a zero operator. In order for $\frac{\eta}{2}\int_{\mathcal{X}}\left(2\frac{e^{(Pf(s^*)(x)-g(x))/\epsilon}}{1+e^{(Pf(s^*)(x)-g(x))/\epsilon}}-1\right)Pf_{s^*}(\cdot)(x)dx=0$ for all $\eta$. We need to set $2\frac{e^{(Pf(s^*)(x)-g(x))/\epsilon}}{1+e^{(Pf(s^*)(x)-g(x))/\epsilon}}-1=0$, and this only happens when $Pf(s^*)(x)-g(x)=0$, namely, when $s^*$ is feasible. So far, we have shown that $s^*$ is necessarily a feasible point of (P) if it is a local min of the merit functional. Now suppose that $\exists \Tilde{s}\in\mathcal{F}\cap\mathcal{N}_{\epsilon}(s^*)$ such that $J(\Tilde{s})<J(s^*)$. Equivalently, we have $M(\Tilde{s},\eta)<M(s^*,\eta)$ which contradicts the fact that $s^*$ is a local min of the merit functional. Now we have completed the proof of (1).\\
    Proof of $(2)$: To show that $h_k$ is a descent direction for the merit functional, note that $h_k$ is the solution of the $k-th$ step of (QP\textsuperscript{k}), hence it must satisfy the first-order necessary condition of (QP\textsuperscript{k}). Namely, we have
    \begin{align*}
        J_{s_k}(\cdot)+H_{s_k}(h_k)(\cdot)+\int_{\mathcal{X}}\lambda_k(x)Pf_{s_k}(\cdot)(x)dx&=0\\
        Pf(s_k)(x)+Pf_{s_k}(h_k)(x)-g(x)&=0.
    \end{align*}
    Recall the merit functional:
    \[
    M(s_k,\eta)=J(s_k)+\int_{\mathcal{X}}\lambda_k(x)(Pf(s_k)(x)-g(x))dx+\frac{\eta}{2}\int_{\mathcal{X}}|Pf(s_k)(x)-g(x)|dx.
    \]
    $h_k$ is a descent direction if the first variation of the merit functional along $h_k$ is negative. The first variation of the merit functional along $h_k$ is
    \begin{align*}
        M_{s_k}(h_k)&=J_{s_k}(h_k)+\int_{\mathcal{X}}\lambda_k(x)Pf_{s_k}(h_k)(x)dx+ G_{s_k}(h_k)\\
        &=-H_{s_k}(h_k)(h_k)+G_{s_k}(h_k)\\
        &=-H_{s_k}(h_k)(h_k)+\frac{\eta}{2}\int_{\mathcal{X}} \frac{Pf(s_k)(x)-g(x)}{|Pf(s_k)(x)-g(x)|}Pf_{s_k}(h_k)(x) dx\\
        &=-H_{s_k}(h_k)(h_k)+\frac{\eta}{2}\int_{\mathcal{X}} \frac{-Pf_{s_k}(h_k)(x)}{|Pf_{s_k}(h_k)(x)|}Pf_{s_k}(h_k)(x) dx\\
        &=-H_{s_k}(h_k)(h_k)-\frac{\eta}{2}\int_{\mathcal{X}}\left|Pf_{s_k}(h_k)(x)\right|dx,
    \end{align*}
    where the second equality is due to the first equation in the first-order condition of $(QP\textsuperscript{k})$ and the penultimate equality is due to the second equation in the first-order condition of $(QP\textsuperscript{k})$. Observe the final expression to conclude that, $ M_{s_k}(h_k)<0$ for large enough $\eta$ since $Pf_{s_k}(h_k)(x)\neq 0$ for $h_k\neq 0$. 
\end{proof}    

For each iteration of (QP\textsuperscript{k}), we compute a descent direction $h_k$ for the merit functional. We could choose the step size of $\alpha_k$ according to suitable line search criteria, for example, the Armijo rule:
\begin{align*}
    M(s_k+\alpha_k h_k,\eta)\leq M(s_k,\eta)+c\alpha_k M_{s_k}(h_k),
\end{align*}
for some $0<c<1$.
\begin{theorem}~\label{thm: global-convergence} 
 For large enough $\eta$, the SQP iterations with step size chosen according to the Armijo rule are guaranteed to converge to a critical point of the merit functional.
\end{theorem}
\begin{proof}
    As shown in the proof of Theorem~\ref{thm: Merit}, $M_{s_k}(h_k)$ is negative for any non-negative direction $h_k$. By imposing the Armijo rule, we have
    \begin{align*}
         M(s_k+\alpha_k h_k,\eta)\leq M(s_k,\eta)+c\alpha_k M_{s_k}(h_k).
    \end{align*}  
    By definition of the first variation of the merit functional, we have: 
    \begin{align*}
         M(s_k+\alpha_k h_k,\eta)-M(s_k,\eta) = M_{s_k}(\alpha_k h_k)+o(\|\alpha_k h_k\|).
    \end{align*}
    Recall that the first variation $M_{s_k}(\alpha_k h_k)$ is linear in $\alpha_k$ for a fixed $h_k$, we obtain:
    \begin{align*}
        M(s_k+\alpha_k h_k,\eta)
        &= M(s_k,\eta) + M_{s_k}(\alpha_k h_k)+o(\|\alpha_k h_k\|)\\
        &= M(s_k,\eta) + \alpha_kM_{s_k}(h_k)+o(\|\alpha_k h_k\|)\\
        &= M(s_k,\eta) + c \alpha_kM_{s_k}(h_k)+ (1-c) \alpha_kM_{s_k}(h_k)+o(\|\alpha_k h_k\|).
    \end{align*}
    Observe the last two terms, $(1-c) \alpha_kM_{s_k}(h_k)<0$ and $\epsilon \alpha_k \|h_k\|>0$. Since $\epsilon\rightarrow 0$ as $\|\alpha_k h_k\|\rightarrow 0$, we know that for any descent direction $h_k$, there exists a $\alpha_k$ small enough such that $(1-c) \alpha_kM_{s_k}(h_k)+o(\|\alpha_k h_k\|)<0$. Therefore, an $\alpha_k$ that satisfies the Armijo rule can always be found.
    Now, we have
    \begin{align*}
        \sum_{k=0}^{K} M(s_k+\alpha_k h_k,\eta)- M(s_k,\eta) \leq c \sum_{k=0}^{K}\alpha_kM_{s_k}(h_k).
    \end{align*}
    This implies that:
    \begin{align*}
        M(s_K+\alpha_K h_K,\eta)-M(s_0,\eta)\leq c \sum_{k=0}^{K}\alpha_k M_{s_k}(h_k).
    \end{align*}
    Note that $\alpha_k M_{s_k}(h_k)< 0$ for all $k$. Since the left-hand side is negative and bounded from below, the right-hand side summation must be convergent, therefore, we necessarily have either $\alpha_k \rightarrow 0$ or $M_{s_k}(h_k) \rightarrow 0$ and either implies that the merit functional converges to a critical point.
\end{proof}
The convergence results we have discussed so far are at the level of the original problem (P) and its SQP iterates. In practice, however, the infinite-dimensional formulation cannot be solved directly. Any implementable algorithm must replace the transport map with a finite-dimensional parameterization and approximate the pushforward constraint accordingly. The next section makes this transition explicit by introducing the computational problem and discussing the consistency of these approximations.

\subsection{The Discrete SQP}\label{subsec:discr-SQP}

At the $k$th iteration of SQP, and at the current iterate $s_{n,k} \in \mathcal{S}_n$, we construct the $k$th quadratic subproblem of (CP):
\begin{align*}
    &\min_{\delta_{s_n,k}\in\Omega_n} \nabla J(s_{n,k})^T\delta_{s_n,k} + \frac{1}{2}\delta_{s_n,k}^T H_k \delta_{s_n,k} \tag{CPQP\textsuperscript{k}}\\
    &\text{s.t.} \\
    & \nabla \tau^{(m)}(s_{n,k})^T \delta_{s_n,k} + \tau^{(m)}(s_{n,k}) = 0,
\end{align*}
where $s_{n,k+1} = s_{n,k} + \delta_{s_n,k}$. The KKT condition of (CPQP\textsuperscript{k}) is given by the linear system:
\begin{align*}
    G(s_{n,k}, \lambda_{m,k})\delta^k = -\Psi(s_{n,k}, \lambda_{m,k}),
\end{align*}
where
\begin{align*}
\Psi(s_{n,k}, \lambda_{m,k}) = 
\begin{bmatrix}
\nabla \mathcal{L}(s_{n,k}, \lambda_{m,k}) \\
\tau^{(m)}(s_{n,k})
\end{bmatrix}, \quad
G(s_{n,k}, \lambda_{m,k}) = 
\begin{bmatrix}
H\mathcal{L}(s_{n,k}, \lambda_{m,k}) & \nabla \tau^{(m)}(s_{n,k}) \\
\nabla \tau^{(m)}(s_{n,k})^T & 0
\end{bmatrix}.
\end{align*}
The solution vector $\delta^k = [\delta_{s_{n,k}}^T, \delta_{\lambda_{m,k}}^T]^T$ contains the descent directions for the optimization variables and the updates for the Lagrange multipliers.

To control the step size, consider the merit function with penalty parameter $\eta > 0$:
\begin{align*}
    M(s_{n,k}, \lambda_{m,k}, \eta) = \mathcal{L}(s_{n,k}, \lambda_{m,k}) + \frac{1}{2}\eta\|\tau^{(m)}(s_{n,k})\|^2_2.
\end{align*}
Starting from $\alpha_k = 1$, we decrease the step size until it satisfies the Armijo line search criterion:
\begin{align*}
    M(s_{n,k+1}, \lambda_{m,k+1}, \eta) \leq M(s_{n,k}, \lambda_{m,k}, \eta) + c\alpha_k \nabla M(s_{n,k}, \lambda_{m,k}, \eta)^T \delta_{s_{n,k}}.
\end{align*}

\begin{definition}[Second-Order Sufficient Condition]
\label{def:SOSC}
Let $(s_{n,0}, \lambda_{m,0})$ satisfy the first-order KKT system:
\begin{align*}
\nabla \mathcal{L}(s_{n,0}, \lambda_{m,0}) = 0,
\qquad
\tau^{(m)}(s_{n,0}) = 0,
\end{align*}
and assume the Jacobian $\nabla\tau^{(m)}(s_{n,0})$ has full row rank (LICQ). Denote by
\begin{align*}
\mathcal{C}(s_{n,0}, \lambda_{m,0})
\;=\;
\Bigl\{\,w\in \mathcal{S}_{n}
       :\;
       \nabla\tau^{(m)}(s_{n,0})^{\top}w = 0
\Bigr\}
\end{align*}
the critical cone (null-space of the Jacobian at the solution). We say that \textnormal{(CP)} satisfies the second-order sufficient condition at $(s_{n,0}, \lambda_{m,0})$ if there exists a constant $\alpha > 0$ such that:
\begin{align*}
w^{\top} H\mathcal{L}(s_{n,0}, \lambda_{m,0})\,w
\;\;\ge\;\;
\alpha\,\|w\|^{2},
\qquad
\forall\,w\in\mathcal{C}(s_{n,0}, \lambda_{m,0})\setminus\{0\}.
\end{align*}
\end{definition}

\begin{proposition}\label{prop:discrete-LICQ}
For any $m$ and any $n > m$, and any feasible transport map $s_n \in \mathcal{S}_n$ for \textnormal{(CP)}, there exists a partition $\mathcal{P}_m$ for the discretized pushforward constraints such that the computational objective \textnormal{(CP)} satisfies the linearly independent constraint qualification (LICQ) at $s_n$.
\end{proposition}

\begin{proof}
Recall that the first variation of the PF operator at a feasible $s_n$ along $h$ is:
\begin{align*}
Pf_{s_n}(h)(x) =& \frac{1}{2\pi} \int_{\mathbb{R}} -e^{itx}\,e^{-it\,\mathbb E[s_n(Y)]}\mathbb E[ite^{it(s_n(Y)-\mathbb E[s_n(Y)]}(h(Y)-\mathbb E[h(Y)])] dt\\
    -&g_{s_n}'(x)\mathbb E [h(Y)],    
\end{align*}
and $Pf_{s_n}(h)$ is Fr\'echet, meaning that it is linear in $h$. For the computational objective \textnormal{(CP)} with an $n$-dimensional approximation to the solution space and $m$-dimensional discretization of the functional constraint, the direction $h \in \mathcal{S}_n$ is parameterized by $b:=(b_0,\cdots,b_{n-1})\in\mathbb{R}^n$ such that $h = \sum_{j=0}^{n-1} b_j \chi_j$, where $\{\chi_0, \dots, \chi_{n-1}\}$ form an orthonormal function basis. Combining these facts, we have $Pf_{s_n}(h) = \sum_{j=0}^{n-1} b_j Pf_{s_n}(\chi_j)$. Therefore, the Jacobian matrix of the discretized constraints $\tau^{(m)}$ is:
\begin{align*}
\nabla\tau^{(m)}(s_n) =
\begin{bmatrix}
Pf_{s_n}(\chi_0)(d_1) & \cdots & Pf_{s_n}(\chi_{n-1})(d_1) \\
\vdots & \ddots & \vdots \\
Pf_{s_n}(\chi_0)(d_m) & \cdots & Pf_{s_n}(\chi_{n-1})(d_m)
\end{bmatrix}.
\end{align*}


Since $\{\chi_0, \dots, \chi_{n-1}\}$ are orthonormal, they are linearly independent. We can always find a partition $\mathcal{P}_m$ such that 
\begin{align*}
\begin{bmatrix}
\chi_0(d_1) & \cdots & \chi_{n-1}(d_1) \\
\vdots & \ddots & \vdots \\
\chi_0(d_m) & \cdots & \chi_{n-1}(d_m)
\end{bmatrix}
\end{align*}
is row independent—observe that if $\sum_{j}\alpha_j \chi_j \equiv 0$, then for all $i\in\{0,\cdots,n-1\}$, we have:
\begin{align*}
    &\langle \chi_i, \sum_{j}\alpha_j \chi_j \rangle = \alpha_i \|\chi_i\|= 0,
\end{align*}
which implies that $\alpha_i = 0$. By Proposition~\ref{prop:injective} and the fact that injective linear transformations preserve linear independence \citep[Theorem 8, p.~80]{hoffmann1971linear}, we can conclude row independence of the constraint Jacobian, i.e., the LICQ condition is satisfied.
\end{proof}
\begin{comment}
Note that the showing linearly independent constraint qualification is satisfied at $s$ is equivalent to showing that for all $j\in\{1,\cdots,m\}$,
\begin{align*}
    \sum_{i=1}^{n-1}\alpha_i Pf_s(h_{b_{i}})(d_j) = Pf_s(\sum_{i=1}^{n-1}\alpha_ih_{b_{i}})(d_j) = 0
\end{align*}
if and only if $\alpha_i = 0$ for all $i\in\{0,\cdots,n-1\}$. We have shown previously that $Pf_s(h)$ is a zero function if and only if $h$ is a zero function. That is, there is one point $\tilde{d}$ on which 
\end{comment}

\begin{theorem}[Local quadratic convergence of discrete SQP]
\label{thm:SQP_quad}
Let $\{(s_{n,k}, \lambda_{m,k})\}$ be the iterates generated by the discrete SQP algorithm using the exact Hessian. Suppose $(s_{n,0}, \lambda_{m,0})$ satisfies the second-order sufficient condition of \textnormal{(CP)}. If the initial point $(s_{n,1}, \lambda_{m,1})$ is sufficiently close to $(s_{n,0}, \lambda_{m,0})$, then there exists a partition $\mathcal{P}_m$ such that:
\begin{align*}
\|s_{n,k+1}-s_{n,0}\| + \|\lambda_{m,k+1}-\lambda_{m,0}\| = O\bigl( \|s_{n,k}-s_{n,0}\| + \|\lambda_{m,k}-\lambda_{m,0}\| \bigr)^2.
\end{align*}
i.e. the method converges quadratically to the solution pair.
\end{theorem}

\begin{proof}
The original objective functional and the constraint are shown to be twice continuously Fr\'echet differentiable with Lipschitz continuous second variations. For \textnormal{(CP)}, we consider transport maps from the $n$-degree polynomial basis where the differentiability remains. Combining with Proposition~\ref{prop:discrete-LICQ},
the result follows from \cite[Theorem 18.4]{wright1999numerical}.
\end{proof}
% The results in this section provide a principled solver for the finite-dimensionalized Monge OT problem (CP), together with conditions that ensure local convergence.
% We next evaluate the proposed discretization-optimization method on representative one-dimensional transport problems and compare with available analytic solutions.

\zihe{
\begin{remark}
The consistency results of Section~\ref{sec:cp} and the convergence results of 
Section~\ref{sec:SQP} are related but not directly compositional. 
Theorem~\ref{thm:rate-main} guarantees that if (P) admits a global solution $s^*$, the objective value $J(s^*_{nm})$ at the global solution of (Cp) satisfies $|J(s^*_{nm}) - J(s^*)| = O(1/n + 1/m)$, while 
Theorem~\ref{thm:SQP_quad} guarantees that the SQP iterates converge quadratically 
to a local solution $s^*_{nm}$ of (CP) for fixed $n$ and $m$. However, to combine
these into an end-to-end guarantee, that is the convergence of the solution of (CP) under discrete SQP method to the the global solution of (P), requires the global convergence of discrete SQP. This is a potential direction for future work.
\end{remark}}

\section{Experiments}\label{sec:experiments}
We demonstrate the utility of the SQP method on several experiments. All experiments transport one-dimensional distributions. 
%To test the generalization of our method, We include one example with concave transport cost from \cite{McCann1999Exact}. 

\refstepcounter{experiment}\label{exp:exp-to-gaus}
\subsection*{Experiment~\ref{exp:exp-to-gaus}: Exponential to Gaussian}
In the first experiment, we consider the classic Brenier setting of a quadratic displacement transport cost $c(x, y) = (x - y) ^2$. For reference, the general optimal transport map for 1D distributions with strictly convex transport costs is
\begin{align*}
s(x) \;=\; F_{\nu}^{-1}\!\big(F_{\mu}(x)\big), 
\end{align*}
where $F_{\mu}$ stands for the CDF for the source measure and $F_{\nu}$ stands for the target measure.
%For all of the experiments, we compare the analytical solution with the SQP solution.

\paragraph{\it Solver configuration.}
We consider transport of an exponential $(\lambda=1)$ source density $f(x)=e^{-x}$, to a standard Gaussian target density $g(x) = \frac{1}{\sqrt{2\pi}}e^{-x^2/2}$. The map $s(x)$ is parameterized as a piecewise linear function on a grid based on the quantile of the source distribution. We use a quantile based grid on the support of $g$ for the constraint evaluation. Specifically, we choose a set of uniformly spaced probability levels
\begin{align*}
v_j = p_{\mathrm{lo}} +
\frac{j-1}{m_d-1}\bigl(p_{\mathrm{hi}}-p_{\mathrm{lo}}\bigr),
\qquad j=1,\ldots,m_d,    
\end{align*}
with $0<p_{\mathrm{lo}}\ll 1$ and $p_{\mathrm{hi}}\approx 1$ and the corresponding constraint grid on the target domain is then defined by
\begin{align*}
d_j := F_{\nu}^{-1}(v_j)
\end{align*}
This way, each constraint point represents an equal portion of probability mass
under the target distribution. This construction concentrates resolution in the rapidly decaying tail of the source and in the extreme Gaussian tail, yielding uniform accuracy in probability space. 

In addition, we want to accurately capture the nonuniform tail behavior of the exponential source. Hence, we also adopt a quantile based grid for the piecewise linear transport map and the objective function quadrature. Recall the CDF of the exponential$(\lambda=1)$
distribution, $F_{\mu}(x) = 1-e^{-x}$ for $x \ge 0$. We select probability levels
\begin{align*}
u_i = p_{\mathrm{lo}} +
\frac{i-1}{m_x-1}\bigl(p_{\mathrm{hi}}-p_{\mathrm{lo}}\bigr),
\qquad i=1,\ldots,m_x ,
\end{align*}
and the corresponding source grid is
\begin{align*}
x_i := F_{\mu}^{-1}(u_i) = -\log(1-u_i),
\qquad w_i = u_{i}-u_{i-1}
\end{align*}
with $u_0 =0$ and  $w_i$ acting as probability weights in the objective.

The analytic reference map in this setting is
\begin{align*}
s^*(x)=F_{\nu}^{-1}(1-e^{-x}).
\end{align*}
We use this as the ground truth for comparison. The learned map reproduces this function with high accuracy over the entire range of $x$, and the pushforward distribution matches the target distribution within numerical tolerance. Deviations near the truncation boundary at large $x$ are primarily due to the finite grid and can be reduced by extending the domain or increasing grid resolution.

\begin{figure}[ht]
    \centering
    \includegraphics[width=0.55\linewidth]{graphs/E2G_params101_200pieces_601_401_500ite1.png}
    \caption{Source $\nu(x)$ and target densities $\mu(y)$ for the Exponential$\to$Gaussian experiment, comparing the learned pushforward $P_f(s)(y)$ with the true target.}
    \label{fig:e2g-densities}
\end{figure}

\begin{figure}[ht]
    \centering
    \includegraphics[width=0.8\linewidth]{graphs/E2G_params101_200pieces_601_401_500ite2.png}
    \caption{Computed transport map (SQP, piecewise linear) versus the analytic quantile map for the Exponential$\to$Gaussian case.}
    \label{fig:e2g-map}
\end{figure}

\newpage
%I believe this so-called "concave" experiment is completely wrong
\refstepcounter{experiment}\label{exp:nonmono}
\input{exp5_concave_updated}

\newpage

\refstepcounter{experiment}\label{exp:mccann}
\input{mccann_experiment}

\newpage
% \subsection*{Experiment~\ref{exp:mccann}: Concave Transport Cost}
% In this section, we numerically solve the one-dimensional concave problem in \cite{McCann1999Exact}. We recall the precise setting: consider the line segment $[-10,10]$ and the signed excess mass
% \begin{align*}
% d\rho(x)=\sin\!\bigl(\tfrac{\pi}{5}\,x\bigr)dx .    
% \end{align*}
% Let $\mu=\rho_{+}$ (surplus) and $\nu=\rho_{-}$ (deficit); the total masses agree since
% $\int_{-10}^{10}\sin(\tfrac{\pi}{5}x)dx=0$.  We transport $\mu$ to $\nu$ under the strictly
% concave cost
% \begin{align*}
% c(x,y)=\sqrt{\,2\,|x-y|\,}.
% \end{align*}
% In this example the optimal map $s$ is piecewise affine with discontinuity points at $x\in\{-9,-1,1,9\}$:
% \begin{align*}
% s(x)=
% \begin{cases}
% -x-10, & -9<x<-1,\\[2pt]
% -x,    & |x|<1\ \text{or}\ |x|>9,\\[2pt]
% -x+10, & 1<x<9.
% \end{cases}
% \end{align*}
% \begin{comment}
% Thus mass in the central and far–end slivers reflects across the origin ($s(x)=-x$), while mass in
% the two middle bands is shifted by $\pm10$ with slope $-1$ to pair each surplus interval with the
% adjacent deficit interval.  This piecewise structure matches the “no–crossing arcs’’ property
% characteristic of strictly concave transport costs on the line. 
% \end{comment}



% \paragraph{\it Piecewise Linear Map with Optimal Partition}
% %We consider the concave cost line transport setting in which the source and target measures arise from a sinusoidal excess mass on the interval $[-10,10]$; the optimal map is known to be piecewise affine with grid points at $x\in\{-9,-1,1,9\}$. 
% In this experiment, we fix the piece boundaries to the known discontinuity points and parameterize $s(x)$ as a piecewise linear function with one slope and intercept per interval. The constraint is enforced on a uniform grid over $[-10,10]$, and we solve the resulting constrained problem with SQP. 

% Because the partition matches the true solution, the learned map overlays the analytical piecewise–affine solution almost exactly, and the pushforward density closely matches the target across the whole domain. This setup shows that SQP solver fully recover the analytical optimal map when the optimal partition for grid points is supplied.
% \begin{figure}[ht]
%         \centering
%         \includegraphics[width=0.55\linewidth]{graphs/concave optimal partition.png}
%         \label{fig:sub-linear}
%     \caption{SQP Solution vs. Analytical Solution with Optimal Partition}
%     \label{fig:linear-vs-quadratic}
% \end{figure}
% \begin{figure}[ht]
%         \centering
%         \includegraphics[width=\linewidth]{graphs/concave optimal partition target density.png}
%         \label{fig:sub-linear}
%     \caption{SQP Target Density with Optimal Partition}
%     \label{fig:linear-vs-quadratic}
% \end{figure}
% %\clearpage

% \paragraph{\it Piecewise Linear Map with Uniform Partition}
% To assess our method's performance when the piecewise linear optimal transport map's partition for pieces are unknown, we repeat the experiment using an uniform partition of $[-10,10]$ into equal-length intervals, again fitting a piecewise linear $s(x)$ with independent slope and intercept on each interval under the same grid. 

% Without access to the true grid points, the recovered map follows the overall trend of the analytical solution but exhibits mismatches near the true kinks. %By plotting the pushforward density of the transported measure, we find that it is not close to any feasible solution, either. 

% \begin{figure}[ht]
%         \centering
%         \includegraphics[width=0.55\linewidth]{graphs/concave uniform partition.png}
%         \label{fig:sub-linear}
%     \caption{SQP Solution vs. Analytical Solution with Uniform  Partition}
%     \label{fig:linear-vs-quadratic}
% \end{figure}

% \begin{figure}[ht]
%         \centering
%         \includegraphics[width=\linewidth]{graphs/concave uniform partition target density.png}
%         \label{fig:sub-linear}
%     \caption{SQP Target Density with Uniform Partition}
%     \label{fig:linear-vs-quadratic}
% \end{figure}

% =============================================================================
% Experiment: Non-Monotone Optimal Transport with Concave Cost
% Drop into Section 8 of the paper (Numerical Experiments)
% =============================================================================


% \subsection*{Experiment 5: Concave Cost and Non-Monotone Optimal Maps}\label{sec:exp-concave}

% We now consider a setting that lies squarely outside the Brenier regime: a concave transport cost for which the optimal map is \emph{non-monotone}. This experiment demonstrates three features of the SQP framework that cannot be exhibited by methods restricted to monotone maps or by Kantorovich solvers: (i) the ability to compute non-monotone optimal maps; (ii) the landscape of distinct local optima that arises from the nonconvexity of the Monge problem; and (iii) the gap between the Monge and Kantorovich optimal values for concave costs.

% %\paragraph{\it Problem setup.}
% Let $\mu = \text{Beta}(2,5)$ and $\nu = \text{Beta}(5,2)$ on $[0,1]$, so that the source distribution concentrates mass near zero while the target concentrates mass near one. We consider the cost
% \begin{equation}\label{eq:concave-cost}
%     c(x,y) = -|x - y|^p, \qquad p = 1/2.
% \end{equation}
% This cost is concave in $|x-y|$, so it \emph{rewards} large displacement: the optimizer prefers to send mass far from its origin. Since $c$ is concave rather than convex, Brenier's theorem does not apply; there is no reason to expect the optimal transport map to be the gradient of a convex function, nor even to be monotone increasing. Indeed, as we will see, the global optimum is non-monotone.

% We verify that $c$ satisfies the assumptions of our framework. The cost is smooth on $\{(x,y) : x \neq y\}$; we use the smoothed version $c_\varepsilon(x,y) = -(|x-y|^2 + \varepsilon)^{p/2}$ with $\varepsilon = 10^{-8}$ to ensure differentiability at $x = y$. The mixed derivative $D^2_{xy} c_\varepsilon$ is nonzero away from $x = y$, consistent with Assumption~7.

% \paragraph{\it Solver configuration.}
% We parameterize the transport map as a piecewise linear function with $n = 30$ pieces (i.e., $n+1 = 31$ parameters), without imposing a monotonicity constraint on the parameter vector. The pushforward constraint $s_\# \mu = \nu$ is enforced via a histogram-based discretization with $m = 20$ bins: for each target bin $B_j = [y_j, y_{j+1})$, we require $\mu(s^{-1}(B_j)) = \nu(B_j)$. This histogram formulation naturally accommodates non-monotone maps, unlike the change-of-variables formulation $f(s^{-1}(y))/|s'(s^{-1}(y))| = g(y)$ used in Experiments~1--4, which requires $s$ to be a diffeomorphism.

% We solve the resulting finite-dimensional problem via an augmented Lagrangian method, which shares the same KKT structure as the SQP subproblem. The inner minimization is performed by L-BFGS-B with box constraints $a_j \in [10^{-4}, 1-10^{-4}]$ (but no monotonicity constraint). We run the outer loop for up to 40 iterations with a convergence tolerance of $\|\tau\|_\infty < 5 \times 10^{-3}$.

% \paragraph{\it Multiple initializations.}
% Because the Monge problem with concave cost is nonconvex, different initializations may converge to different local optima. We exploit this by running the solver from four distinct starting maps:
% \begin{enumerate}
%     \item \emph{Quantile map} (monotone increasing): $s_0(x) = F_\nu^{-1}(F_\mu(x))$, the Brenier-optimal map for quadratic cost.
%     \item \emph{Anti-quantile map} (monotone decreasing): $s_0(x) = F_\nu^{-1}(1 - F_\mu(x))$, which sends low-$x$ mass to high-$y$ values.
%     \item \emph{V-shape} (non-monotone): $s_0(x) = F_\nu^{-1}(|2F_\mu(x) - 1|)$, which sends both tails of $\mu$ to one tail of $\nu$.
%     \item \emph{Tent} (non-monotone): $s_0(x) = F_\nu^{-1}(1 - |2F_\mu(x) - 1|)$, the complement of the V-shape.
% \end{enumerate}
% Each initialization is constructed to be measure-preserving at the level of quantile functions, ensuring an approximately feasible starting point.

% \paragraph{\it Results.}

% \begin{table}[t]
%     \centering
%     \caption{Concave cost $c(x,y) = -|x-y|^{1/2}$, transporting $\text{Beta}(2,5)$ to $\text{Beta}(5,2)$. Four initializations converge to four distinct local optima. The tent initialization (non-monotone) achieves the lowest Monge objective, beating both monotone candidates. The Kantorovich LP lower bound is $-0.6823$.}
%     \label{tab:concave-results}
%     \begin{tabular}{lcccc}
%         \toprule
%         Initialization & Monotonicity & $J(s)$ & $\|\tau\|_\infty$ & Optimality \\
%         \midrule
%         Quantile       & monotone $\uparrow$    & $-0.6560$ & $5.2 \times 10^{-3}$ & local opt. \\
%         Anti-quantile  & monotone $\downarrow$   & $-0.6491$ & $5.9 \times 10^{-3}$ & local opt. \\
%         V-shape        & non-monotone             & $-0.6337$ & $1.6 \times 10^{-2}$ & local opt. \\
%         {\bf Tent}           & {\bf non-monotone}             & $\mathbf{-0.6565}$ & $8.7 \times 10^{-3}$ & {\bf best found} \\
%         \midrule
%         \multicolumn{2}{l}{Kantorovich LP lower bound} & $-0.6823$ & --- & --- \\
%         \bottomrule
%     \end{tabular}
% \end{table}

% Table~\ref{tab:concave-results} reports the results. Several observations are noteworthy.

% First, the four initializations converge to four \emph{distinct} local optima, with objective values spanning a range of $3.5\%$ (from $-0.634$ to $-0.657$). This directly demonstrates the nonconvexity of the Monge problem under concave costs and underscores the importance of the globalization theory in Section~\ref{subsec:glo_conv}: without a mechanism for escaping poor local optima, the solver's output depends critically on initialization.

% Second, the best map found is \emph{non-monotone}. The tent initialization converges to a map with $J = -0.6565$, which is lower than both the quantile map ($-0.6560$, monotone increasing) and the anti-quantile map ($-0.6491$, monotone decreasing). The optimal map rises steeply from $s(0) \approx 0.6$ to $s(0.2) \approx 1.0$, then decreases gradually to $s(0.8) \approx 0.25$ (see Figure~\ref{fig:concave-maps}). This non-monotone structure reflects the concave cost's preference for large displacement: low-$x$ mass (which is abundant under $\mu = \text{Beta}(2,5)$) is sent to high-$y$ values (which are abundant under $\nu = \text{Beta}(5,2)$), maximizing $|x - y|$, while the remaining mass is distributed to fill out the target density.

% Third, the Monge--Kantorovich gap is clearly visible: the best Monge map achieves $J = -0.6565$, while the Kantorovich LP relaxation attains $-0.6823$, a gap of approximately $3.8\%$. This gap is expected for concave costs~\citep{McCann1999Exact,GangboMcCann1996}: the LP relaxation achieves a lower objective by splitting mass across multiple target locations, which is feasible for couplings but not for deterministic transport maps. The SQP framework correctly solves the \emph{Monge} problem---finding the best map, not the best coupling.

% Fourth, the constraint violations are of order $10^{-3}$ to $10^{-2}$, consistent with the histogram-based discretization. The pushforward check (Figure~\ref{fig:concave-pushforward}) confirms that the histogram of $s(x)$ values, weighted by the source density, closely matches the target density $g(y)$.

% \paragraph{\it Discussion.}
% This experiment demonstrates that the SQP framework operates effectively outside the Brenier regime. In the classical setting of convex costs on $\mathbb{R}^d$, Brenier's theorem guarantees that the optimal map is the gradient of a convex function; in one dimension, this reduces to a monotone increasing map, uniquely determined as the quantile coupling $F_\nu^{-1} \circ F_\mu$. For concave costs, this structural result fails entirely: the optimal map can be monotone decreasing, non-monotone, or even discontinuous, and there may be multiple local optima with genuinely different structure.

% The SQP framework handles this setting naturally because it treats the Monge problem as a constrained optimization problem over the space of transport maps, without imposing structural assumptions (such as monotonicity or convexity of a potential) on the optimizer. The augmented Lagrangian solver navigates infeasible regions while systematically restoring the pushforward constraint, and the multiple-initialization strategy provides a practical mechanism for exploring the landscape of local optima.

% We note that the Monge--Kantorovich gap observed here (${\sim}3.8\%$) is a genuine structural gap, not a discretization artifact: for concave costs with absolutely continuous measures in one dimension, the Kantorovich relaxation is strictly tighter than the Monge problem whenever the optimal coupling is not induced by a map. This gap motivates the need for Monge-side solvers---such as the one developed in this paper---that directly optimize over transport maps rather than over couplings.

% \begin{figure}[t]
%     \centering
%     \includegraphics[width=\textwidth]{graphs/fig2_local_optima.png}
%     \caption{Transport maps from four initializations under concave cost $c(x,y) = -|x-y|^{1/2}$, transporting $\text{Beta}(2,5)$ to $\text{Beta}(5,2)$. Each panel shows the converged map $s(x)$ from one initialization. The tent initialization (bottom right, starred) achieves the best objective and is non-monotone, with both increasing and decreasing segments. The dashed gray line is the identity $s(x) = x$.}
%     \label{fig:concave-maps}
% \end{figure}

% \begin{figure}[t]
%     \centering
%     \includegraphics[width=0.7\textwidth]{graphs/fig3_pushforward_conv.png}
%     \caption{Left: Pushforward check. Histogram of $s(x)$ values (weighted by source density) for the best map, compared against the target density $g(y) = \text{Beta}(5,2)$. Right: Constraint violation $\|\tau\|_\infty$ versus outer iteration for the best initialization.}
%     \label{fig:concave-pushforward}
% \end{figure}

% \begin{figure}[t]
%     \centering
%     \includegraphics[width=0.7\textwidth]{graphs/fig4_setup_gaps.png}
%     \caption{Left: Source and Target densities. Right: Objective values for different local optima, compared to the Kantorovich LP relaxation.}
%     \label{fig:concave-pushforward}
% \end{figure}


% % =============================================================================
% % Experiment: 1D Freeform Lens Design via Optimal Transport
% % Drop into Section 8 of the paper (Numerical Experiments)
% % =============================================================================

% % =============================================================================
% % Experiment: Non-Quadratic Costs from Physics
% % Drop into Section 8 of the paper (Numerical Experiments)
% % =============================================================================

\refstepcounter{experiment}\label{exp:non-quad-phys}
\subsection*{Experiment~\ref{exp:non-quad-phys}: Non-Quadratic Costs from Physics}\label{sec:exp-physics}
This experiment applies the SQP framework to cost functions arising from physical design problems in optics, serving two purposes.  First, it \emph{verifies} the framework against exact solutions: for twisted costs \zihe{whats twisted cost} in 1D, the optimal map is the quantile coupling $F_\nu^{-1} \circ F_\mu$ regardless of the exact cost function (Section~\ref{sec:exp-physics-1d}), so the SQP output can be compared to a known ground truth.  Second, it demonstrates the framework on the \emph{reflector antenna cost} $c(x,y) = -\log|x-y|$, posed on its natural domain $S^1$ where the Gangbo--McCann and Wang--Loeper regularity theory applies (Section~\ref{sec:exp-physics-s1}).


\begin{figure}[ht]
    \centering
    \includegraphics[width=\textwidth]{graphs/fig_lens_schematic.pdf}
    \caption{Physical setups motivating the non-quadratic costs.
    (a)~Cylindrical lens: a uniform beam $f(x)$ is refracted into target intensity $g(y)$ on a screen at distance~$D$; cost is $c_{\textup{lens}}(x,y) = \sqrt{(x-y)^2 + D^2}$.
    (b)~Reflector antenna: a point source is redirected to a prescribed far-field pattern $g(y)$; cost is $c_{\textup{refl}}(x,y) = -\log|x-y|$.  Both shown in 1D cross-section.}
    \label{fig:lens-schematic}
\end{figure}

\subsubsection{Twisted costs on the interval}\label{sec:exp-physics-1d}

\paragraph{\it Setup.}
Consider the following one-dimensional analogue of the freeform illumination problem~\citep{yadav2019monge,romijn2019freeform,prins2015least}.  A cylindrical lens redirects a collimated beam with intensity $f(x)$ into a target intensity $g(y)$ on a screen at distance $D$ (see Figure~\ref{fig:lens-schematic}(a)).  A ray entering at position $x$ and arriving at $y$ travels a distance $\sqrt{(x-y)^2 + D^2}$, yielding the Monge problem with cost
\begin{equation}\label{eq:lens-cost}
    c_{\textup{lens}}(x,y) = \sqrt{(x-y)^2 + D^2}.
\end{equation}
This cost interpolates between two regimes: for $|x-y| \ll D$ it reduces to $D + (x-y)^2/(2D)$ (the Brenier regime), while for $|x-y| \gg D$ it approaches $|x-y|$.  We test it alongside the quadratic cost $c(x,y) = (x-y)^2$ as a baseline.

The source is $\mu = \tfrac{1}{2}\mathcal{L}_{[-1,1]}$ (uniform) and we consider three targets: bimodal (ring), concentrated (Gaussian), and asymmetric.\zihe{may want to specify the target measaures}

\paragraph{\it A 1D degeneracy.}
Both costs satisfy the twist condition: $D^2_{xy}c_{\textup{quad}} = -2$ and $D^2_{xy}c_{\textup{lens}} = -D^2/[(x-y)^2 + D^2]^{3/2} < 0$.  By the Gangbo--McCann theorem~\citep{gangbo1996geometry}, the optimal map exists, is unique, and is monotone increasing. \zihe{I cannot locate this theorem}  But in one dimension, there is exactly one monotone increasing map pushing $\mu$ to $\nu$: the quantile coupling $s(x) = F_\nu^{-1}(F_\mu(x))$.  It is a measure-theoretic fact that the unique monotone rearrangement is independent of the cost function.

The consequence is striking: \emph{all twisted costs produce the same optimal map in 1D.}  The cost affects only the optimal \emph{value} $J^* = \int c(x, s^*(x))\,d\mu(x)$, not the map $s^*$ itself.  This degeneracy is specific to one dimension; in $d \geq 2$, the set of ``monotone'' maps (gradients of convex functions, or more generally $c$-convex maps) is infinite-dimensional, and the cost genuinely selects among them through the $c$-convex potential.

\paragraph{\it Verification results.}
We compute the exact quantile coupling and its objective for each cost, then run the SQP solver (piecewise linear maps with $n = 40$ pieces, $m = 20$ \zihe{Are those 40, 20 or 80, 40. Table 2 says 80, 40} histogram bins, monotonicity enforced) and compare.  Table~\ref{tab:lens-verification} reports the results.  The SQP recovers the exact map to within discretization error for both costs across all three targets, with objective gaps below $0.29\%$.  Figure~\ref{fig:lens-maps} overlays the SQP solutions (colored) on the exact quantile coupling (dashed black), confirming they are visually indistinguishable.

This experiment validates the SQP machinery on a problem with a known closed-form solution, providing confidence in the solver before deploying it on costs where no such solution is available (Experiment~\ref{sec:exp-lagrangian}).

\begin{figure}[t]
    \centering
    \includegraphics[width=\textwidth]{graphs/lensfig/fig1_lens_maps.png}
    \caption{Verification: SQP maps (colored) vs.\ exact quantile coupling (dashed black) for quadratic and lens costs on $[-1,1]$.  The maps are visually indistinguishable, confirming the 1D degeneracy: all twisted costs produce the same optimal map.  Bottom row: pushforward histograms match the target densities.}
    \label{fig:lens-maps}
\end{figure}

\begin{table}[t]
    \centering
    \caption{Verification: SQP vs.\ exact quantile coupling for twisted costs on $[-1,1]$.
    $n=80$ PL pieces, $m=40$ histogram bins.}
    \label{tab:lens-verification}
    \small
    \begin{tabular}{llccc}
        \toprule
        Cost & Target & Exact $J^*$ & SQP $J$ & Gap \\
        \midrule
        Quadratic & ring        & 0.01598 & 0.01598 & ${<}0.01\%$ \\
        Quadratic & concentrated & 0.15591 & 0.15589 & $0.02\%$ \\
        Quadratic & asymmetric  & 0.06080 & 0.06062 & $0.29\%$ \\
        Lens $D=1$ & ring       & 1.00819 & 1.00819 & ${<}0.01\%$ \\
        Lens $D=1$ & concentrated & 1.07380 & 1.07380 & ${<}0.01\%$ \\
        Lens $D=1$ & asymmetric & 1.02971 & 1.02964 & $0.01\%$ \\
        \bottomrule
    \end{tabular}
\end{table}


\subsubsection{Reflector antenna on $S^1$}\label{sec:exp-physics-s1}

\paragraph{\it Physical setup.}
The far-field reflector antenna problem~\citep{wang2004design,glimm2003optical} asks, for a given point source with intensity $f$ over directions $x \in S^{n-1}$, to design a reflector surface that produces a prescribed far-field intensity $g$ over directions $y \in S^{n-1}$.  This can be posed as a Monge optimal transport problem with cost
\begin{equation}\label{eq:reflector-cost}
    c_{\textup{refl}}(x,y) = -\log|x - y|,
\end{equation}
where $|x-y|$ is the Euclidean (chord) distance on the sphere~\citep{wang2004design}.  See Figure~\ref{fig:lens-schematic}(b).

This cost has a logarithmic singularity at $x = y$ but is well-posed on the sphere when the source and target occupy \emph{disjoint} regions---as is physically natural, since the reflector redirects light from one hemisphere to another.  Wang~\citep{wang2004design} established existence and uniqueness of the optimal map via a Monge--Amp\`ere formulation, and Loeper~\citep{loeper2011regularity} showed that the cost-sectional curvature is uniformly positive on $S^{n-1}$, yielding smooth optimal maps for smooth data.  The Ma--Trudinger--Wang condition (A3) is satisfied, making this one of the best-understood non-quadratic costs in the regularity literature~\citep{ma2005regularity,loeper2009regularity}.

\paragraph{\it 1D cross-section on $S^1$.}
We consider the 1D analogue problem that seeks to transport on the unit circle $S^1$, with source and target on disjoint arcs.  Parametrizing $S^1$ by angle $\theta \in [0, 2\pi)$, the chord distance is $|x-y| = 2\sin(|\theta_x - \theta_y|/2)$ and the cost becomes
\[
    c(\theta_x, \theta_y) = -\log\bigl(2\sin(|\theta_x - \theta_y|/2)\bigr).
\]
We place the source on the upper arc $[\delta, \pi - \delta]$ and the target on the lower arc $[\pi + \delta, 2\pi - \delta]$ with $\delta = 0.15$, ensuring a minimum chord distance of $\approx 0.30$ between arcs. The cost is therefore bounded in $[-0.69, 1.21]$.

\paragraph{\it Results.}
Since the twist condition holds on $S^1$ ($D^2_{\theta_x \theta_y} c < 0$ away from the diagonal) and the source is absolutely continuous, the optimal map is again the quantile coupling---now between densities on disjoint arcs.  The SQP solver (with monotonicity enforcement) recovers it to within $0.04\%$ of the exact objective for all three targets; see Figure~\ref{fig:reflector-s1} (b, c, d).  The solver converges rapidly, typically in 1--3 outer iterations, reflecting the benign optimization landscape when the cost is bounded and the initialization is close to optimal.

Figure~\ref{fig:reflector-s1} (a) visualizes the transport on the circle: rays (green) connect source positions on the upper arc (blue) to target positions on the lower arc (red), showing the reflector redirecting intensity across the circle.

\begin{figure}[t]
    \centering
    \includegraphics[width=\textwidth]{graphs/lensfig/fig6_reflector_s1.png}
    \caption{Reflector antenna on $S^1$.  (a)~Transport rays on the circle connecting source arc (blue) to target arc (red).  (b--d)~SQP maps (green) vs.\ exact quantile coupling (dashed) for ring, concentrated, and asymmetric targets.  Gaps to exact objective are below $0.04\%$.}
    \label{fig:reflector-s1}
\end{figure}

\refstepcounter{experiment}\label{exp:Bilevel}
\subsection*{Experiment~\ref{exp:Bilevel}: Lagrangian Obstacle Cost: Bilevel Optimal Transport}\label{sec:exp-lagrangian}
Our next experiment addresses optimal transport with a \emph{Lagrangian cost}, where the cost function is itself defined through an inner variational problem. This is a setting identified in recent literature as lacking classical solvers~\citep{fan2021scalable}.  This experiment demonstrates that the SQP framework extends naturally to costs that are not given in closed form, a capability that has so far been available only through neural-network-based methods~\citep{pooladian2024neural,asadulaev2024neural,fan2021scalable}.

\paragraph{\it Problem setup.}
Consider transporting mass from $x$ to $y$ along a path $\gamma:[0,1] \to \mathbb{R}$ through a potential landscape $V(z)$.  The cost of this transport is the minimum-action path:
\begin{equation}\label{eq:lagrangian-cost}
    c(x,y) = \inf_{\substack{\gamma \in C^1([0,1]) \\ \gamma(0)=x,\, \gamma(1)=y}}
    \int_0^1 \bigl[\lambda |\dot\gamma(t)|^2 + V(\gamma(t))\bigr]\, dt,
\end{equation}
where $\lambda > 0$ controls the kinetic energy penalty and $V \geq 0$ is a potential barrier.
This is the action functional of classical mechanics; the optimal path satisfies the Euler--Lagrange equation $2\lambda\ddot\gamma = V'(\gamma)$, and the cost $c(x,y)$ is the corresponding Hamilton--Jacobi value function at unit time.
The Monge problem is then a \emph{bilevel optimization}: the outer problem chooses the transport map $s$, while the inner problem~\eqref{eq:lagrangian-cost} computes the cost of each assignment.

\paragraph{\it Why this cost is interesting.}
When $V = 0$, the optimal path is a straight line and $c(x,y) = \lambda(x-y)^2$, recovering the quadratic cost.  A nonzero barrier $V$ deforms the cost in a way that depends on the potential landscape between $x$ and $y$, not merely on the displacement $|x-y|$.  An important subtlety revealed by our GPU-accelerated solver is that the inner path optimization~\eqref{eq:lagrangian-cost} is itself nonconvex near the barrier. Specifically, for source--destination pairs $(x,y)$ that both lie near the barrier, the action functional has multiple local minima corresponding to qualitatively different paths.  In particular, a ``straight-through'' path that traverses the barrier incurs high potential cost but low kinetic cost, while a ``detour'' path that swings away from the barrier and returns incurs low potential cost but higher kinetic cost.  For a single Gaussian barrier with height $h = 5$ and width $\sigma = 0.15$ at $\lambda = 0.5$, the detour path gives $c(0,0) \approx 1.64$, substantially less than the stationary path cost $V(0) = 5.0$.  Solving the inner problem to global optimality is therefore essential for obtaining the correct cost table.

We numerically verify the twist condition $D^2_{xy}c < 0$  on the GPU-computed cost table using finite differences on a $160 \times 160$ grid.  The twist condition holds for approximately $99\%$ of the $(x,y)$ domain for both barrier configurations, with a small violation region (${\sim}1\%$) near the diagonal where the optimal path transitions between the straight-through and detour regimes.  Since the twist condition holds almost everywhere, the Gangbo--McCann theory~\citep{GangboMcCann1996} applies generically, and the optimal map is expected to be monotone increasing---consistent with our numerical results.

The significance of this experiment is not that classical theory fails, but rather that classical \emph{solvers} are inapplicable.  Evaluating $c(x,y)$ at a single point requires solving a variational problem; there is no closed-form expression for the cost, its gradient, or its Hessian.  Standard Monge--Amp\`ere PDE methods require analytic derivatives of the cost function, which are unavailable here.  The SQP framework handles this setting by precomputing $c(x,y)$ on a grid, fitting a spline interpolant, and using the spline's analytic derivatives in place of those of the true cost.  This ``precompute--interpolate--optimize'' architecture is generic: the inner variational problem could be replaced by any black-box simulator (for instance, molecular dynamics, ray tracing, or fluid solve), and the same pipeline applies.

This class of problems has been studied in the recent machine learning literature as ``optimal transport with Lagrangian costs''~\citep{pooladian2024neural}, where both the transport map and the inner paths are parameterized by neural networks.  Prior computational work on Lagrangian costs relied on stochastic methods (neural SDEs, Schr\"odinger bridges) that do not produce deterministic transport maps; Pooladian et al.\ provide the first neural framework yielding both deterministic maps and optimal paths.  Our experiment provides the first \emph{classical} (non-neural) optimization framework for this problem class, with the convergence guarantees of Theorems~\ref{thm: local-conv}--\ref{thm: global-convergence}.

\paragraph{\it GPU-accelerated cost precomputation.}
We solve all $n_g^2$ inner path problems simultaneously on a GPU using batched Jacobi iteration on the Euler--Lagrange equation.  The EL equation for interior path point $\gamma_k$ is
\[
    2\lambda\frac{2\gamma_k - \gamma_{k-1} - \gamma_{k+1}}{\Delta t} + \tfrac{1}{2}\Delta t\bigl(V'(\text{mid}_{k-1}) + V'(\text{mid}_k)\bigr) = 0,
\]
discretized with $n_{\text{path}} = 25$ interior points.  This is a diagonally dominant tridiagonal system, so the Jacobi update $\gamma_k \leftarrow \gamma_k - F_k / (4\lambda/\Delta t)$ converges geometrically.  All $n_g^2 = 25{,}600$ path problems are updated simultaneously as a single tensor operation, requiring approximately 500 iterations to reach residuals below $10^{-10}$.  On an NVIDIA GPU, the full $160 \times 160$ cost table is computed in roughly 10 seconds---a substantial speedup over the sequential CPU solver, which requires several minutes.  The resulting cost table is fitted with a bicubic spline to provide smooth interpolation and analytic partial derivatives for the outer transport solver.

\paragraph{\it Experimental setup.}
We consider two potential barriers:
\begin{equation}\label{eq:barriers}
    V_{\textup{single}}(z) = 5\exp\Bigl(-\frac{z^2}{2 \cdot 0.15^2}\Bigr), \qquad
    V_{\textup{double}}(z) = 4\exp\Bigl(-\frac{(z-0.3)^2}{2 \cdot 0.12^2}\Bigr) + 4\exp\Bigl(-\frac{(z+0.3)^2}{2 \cdot 0.12^2}\Bigr).
\end{equation}
The source is $\mu = \text{Uniform}([-1,1])$ and we test three target distributions: bimodal (ring), concentrated (Gaussian), and asymmetric.

The transport solver uses piecewise linear maps with $n = 40$ pieces, histogram pushforward with $m = 20$ bins, and the SQP method of Section~\ref{subsec:discr-SQP} with Armijo line search.

\paragraph{\it Multiple initializations.}
Because the Monge problem with Lagrangian cost is nonconvex, we run the solver from seven distinct initializations:
\begin{enumerate}
    \item \emph{Quantile} (monotone increasing): $s_0 = F_\nu^{-1} \circ F_\mu$, the Brenier-optimal map for quadratic cost.
    \item \emph{Anti-quantile} (monotone decreasing): $s_0 = F_\nu^{-1}(1 - F_\mu)$.
    \item \emph{Tent} and \emph{V-shape} (non-monotone): as in Experiment~5.
    \item \emph{LP-barycentric}: the barycentric projection of the Kantorovich optimal plan, defined by $T(x_i) = \sum_j y_j \pi^*_{ij} / \sum_j \pi^*_{ij}$, where $\pi^*$ is the solution of the Kantorovich LP. This extracts a Monge map from the optimal coupling by computing the conditional mean destination for each source point.
    \item \emph{LP-monotone}: the sorted (monotonically increasing) version of the LP-barycentric map, which maintains the LP's assignment structure while imposing monotonicity.
    \item \emph{Cost-aware}: for each PL knot $x_k$, select the destination $y$ minimizing $c(x_k, y) - w \log g(y)$, where $g$ is the target density and $w$ is a scaling parameter. This balances transport cost against the desire to place mass in high-density target regions.
\end{enumerate}
The LP-based initializations use information from the Kantorovich relaxation to seed the Monge solver, providing a warm start that can substantially outperform the generic quantile-based initializations.

\paragraph{\it Results.}

\begin{table}[t]
    \centering
    \caption{Lagrangian obstacle cost: best objectives across barriers and targets.  The LP-barycentric initialization, which uses the Kantorovich optimal plan as a warm start, achieves the best objective in most cases.  Kantorovich gaps are small (0.2\%--3.0\%), indicating near-optimal Monge maps. The cost table is computed on a $160 \times 160$ GPU-accelerated grid with the inner path problems solved to global optimality.\zihe{actual global optimum?, and why for the first row, the J(s) is lower than the Kant. LB?}}
    \label{tab:lagrangian-results}
    \small
    \begin{tabular}{llcccl}
        \toprule
        Barrier & Target & Best init.\ & $J(s)$ & Kant.\ LB & Gap \\
        \midrule
        \multirow{3}{*}{Single}
        & ring         & quantile & $0.1852$ & $0.1910$ & $3.0\%$ \\
        & concentrated & LP-bary  & $0.6399$ & $0.6230$ & $2.7\%$ \\
        & asymmetric   & LP-bary  & $0.3904$ & $0.3845$ & $1.5\%$ \\
        \midrule
        \multirow{3}{*}{Double}
        & ring         & LP-bary  & $0.4025$ & $0.4018$ & $0.2\%$ \\
        & concentrated & LP-bary  & $0.7814$ & $0.7785$ & $0.4\%$ \\
        & asymmetric   & LP-mono  & $0.5291$ & $0.5256$ & $0.7\%$ \\
        \bottomrule
    \end{tabular}
\end{table}

Table~\ref{tab:lagrangian-results} summarizes the results.  Several findings stand out.

\emph{Near-optimal Monge maps.}  The Kantorovich gaps are small across all configurations, ranging from $0.2\%$ (double barrier, ring target) to $3.0\%$ (single barrier, ring target).  This is a substantial improvement over preliminary results obtained with the CPU cost solver, where gaps regularly exceeded $30\%$ due to the inner path problems being trapped at local minima near the barrier.  The GPU solver's ability to find globally optimal paths, including the detour paths that avoid the barrier, produces a more accurate cost table and, consequently, tighter transport maps.

\emph{LP-based initializations dominate.}  The LP-barycentric initialization achieves the best objective in four of six configurations, and the LP-monotone initialization wins in a fifth.  The quantile map, which is optimal for the quadratic cost, wins only for the single-barrier ring target.  This demonstrates the value of using the Kantorovich relaxation to warm-start the Monge solver: the LP plan encodes structural information about the cost (in particular, the barrier avoidance structure) that the generic quantile initialization lacks.

\emph{Multiple local optima.}  Figure~\ref{fig:lagrangian-optima} shows that the seven initializations converge to distinct local optima, with objective values spanning a wide range.  For the double barrier with ring target, the anti-quantile initialization reaches $J = 1.89$, which is nearly five times worse than the best map ($J = 0.40$).  This multiplicity confirms the nonconvexity of the Monge problem with Lagrangian cost and motivates the use of diverse initializations.

\emph{Barrier avoidance.}  The optimal maps exhibit a pronounced step or plateau near the barrier (Figure~\ref{fig:lagrangian-maps}).  For the single barrier, the map compresses near $s(x) \approx 0$, routing mass from the left of the barrier to the left of the target and mass from the right to the right, with a rapid transition through the barrier zone.  The double barrier produces a two-step structure.  Figure~\ref{fig:lagrangian-vs-quad} overlays the barrier-cost maps with the quadratic (Brenier) maps, showing that the two coincide away from the barriers but diverge sharply near them.

\emph{Bilevel perspective.}  Figure~\ref{fig:lagrangian-paths} shows the optimal transport paths connecting source points to their assigned destinations.  Paths whose source and destination lie on the same side of the barrier travel nearly straight, while paths that must cross the barrier deflect away from it, swinging through low-potential regions.

\emph{Convergence.}  The convergence plots (Figure~\ref{fig:lagrangian-convergence}) show that the constraint violation $\|\tau\|_\infty$ reaches the tolerance $5 \times 10^{-3}$ within 10--20 SQP iterations for the best initializations.  Some initializations (particularly the quantile map for the double barrier) converge slowly or stagnate at higher constraint violations, reflecting the difficulty of the optimization landscape.

\paragraph{\it Discussion.}
The computational architecture in this experiment is generic: the cost $c(x,y)$ is computed by solving the inner variational problem on a GPU, fitted with a spline, and passed to the SQP framework.  The same pipeline extends to any setting where the cost is defined through simulation or PDE solve.  The GPU acceleration is essential at two levels: first, for solving the $n_g^2$ inner path problems in parallel (a task that is embarrassingly parallel but would take minutes sequentially on CPU); and second, for the SQP derivatives (Jacobian, gradient, Hessian) at each outer iteration.

An important lesson of this experiment concerns the inner problem's nonconvexity.  The action functional~\eqref{eq:lagrangian-cost} has multiple local minima near the barrier, corresponding to qualitatively different path topologies (straight-through vs.\ detour).  A sequential CPU solver using L-BFGS-B from straight-line initialization misses the globally optimal detour paths, producing an incorrect cost table and consequently incorrect transport maps.  The GPU Jacobi solver, by updating all paths simultaneously, provides implicit regularization that helps escape local minima.  This bilevel nonconvexity, where the cost precomputation itself requires global optimization, is a distinctive feature of Lagrangian transport that does not arise in standard optimal transport with closed-form costs.

A second lesson concerns the initialization strategy.  The LP-barycentric initialization consistently outperforms the generic quantile-based initializations because it leverages the solution of the Kantorovich relaxation to warm-start the nonconvex Monge solver. The Kantorovich optimal plan $\pi^*$ encodes the cost structure of the problem. Roughly speaking, for the Lagrangian cost, the plan ``knows'' about the barrier and routes mass around it.  The barycentric projection $s(x_i) = \sum_j y_j \pi^*_{ij} / \sum_j \pi^*_{ij}$ extracts a deterministic map from this plan by computing the conditional mean destination for each source point.  This map is generally not feasible for the Monge problem (its pushforward does not exactly match the target), but it is a far better starting point than the quantile coupling, which encodes only the marginals and ignores the cost entirely.  This suggests a general computational strategy for nonconvex Monge problems: solve the Kantorovich relaxation first, project the optimal plan to a map via barycentric projection, and use the result to initialize a Monge-side solver such as SQP. The Kantorovich LP thus serves a dual role, providing both a lower bound for certifying near-optimality (Table~\ref{tab:lagrangian-results}) and a warm start for accelerating convergence.

\begin{figure}[t]
    \centering
    \includegraphics[width=\textwidth]{graphs/lagrangian cost/fig_lagrangian_0_cost_landscape.png}
    \caption{Lagrangian cost landscape.  Top left: the single and double barrier potentials $V(z)$.  Top center/right: the GPU-computed cost matrices $c(x,y)$ on a $160 \times 160$ grid.  Note the ``valley'' along the diagonal near the barrier, where the detour paths produce costs substantially below $V(x)$.  Bottom left: cost slice at $x = -0.5$ compared to the quadratic cost.  Bottom center/right: the mixed derivative $c_{xy}$, showing the twist condition holds for ${\sim}99\%$ of the domain, with a narrow violation band along the diagonal near the barrier.}
    \label{fig:lagrangian-cost}
\end{figure}

\begin{figure}[t]
    \centering
    \includegraphics[width=\textwidth]{graphs/lagrangian cost/fig_lagrangian_1_local_optima.png}
    \caption{{Local optima for the Lagrangian cost.  Each panel shows all seven initializations overlaid, with the best marked by a star.  Rows: single and double barrier.  Columns: ring, concentrated, and asymmetric targets.  The LP-barycentric initialization (orange) achieves the best objective in most cases, outperforming the generic quantile-based initializations.}}
    \label{fig:lagrangian-optima}
\end{figure}

\begin{figure}[t]
    \centering
    \includegraphics[width=\textwidth]{graphs/lagrangian cost/fig_lagrangian_2_maps_on_cost.png}
    \caption{Best transport maps overlaid on the Lagrangian cost matrix.  The map traces a path through the low-cost regions, avoiding the barrier.  Titles report the Monge objective and Kantorovich gap, which ranges from $0.2\%$ to $3.0\%$.}
    \label{fig:lagrangian-maps}
\end{figure}

\begin{figure}[t]
    \centering
    \includegraphics[width=\textwidth]{graphs/lagrangian cost/fig_lagrangian_3_vs_quadratic.png}
    \caption{Obstacle cost vs.\ quadratic cost.  Away from the barrier, the optimal maps nearly coincide.  Near the barrier, the Lagrangian maps develop a step that has no analogue in the Brenier solution.  The single and double barriers (pink and purple) produce distinct step structures.}
    \label{fig:lagrangian-vs-quad}
\end{figure}

\begin{figure}[t]
    \centering
    \includegraphics[width=\textwidth]{graphs/lagrangian cost/fig_lagrangian_4_convergence.png}
    \caption{Convergence of the SQP method for the Lagrangian cost.  The LP-barycentric and LP-monotone initializations (orange, brown) typically reach feasibility within 10--20 iterations, while the generic initializations (quantile, tent) may stagnate at higher constraint violations.}
    \label{fig:lagrangian-convergence}
\end{figure}

\begin{figure}[t]
    \centering
    \includegraphics[width=\textwidth]{graphs/lagrangian cost/fig_lagrangian_5_paths.png}
    \caption{Optimal transport paths through the barrier for the ring target.  Circles mark source positions $x$; squares mark destinations $s(x)$.  Paths near the barrier (orange shading) deflect away from it, while paths far from the barrier travel nearly straight.  Left: single barrier.  Right: double barrier.}
    \label{fig:lagrangian-paths}
\end{figure}



\refstepcounter{experiment}\label{exp:electrons}
\subsection*{Experiment~\ref{exp:electrons}: Coulomb Optimal Transport and Strictly Correlated Electrons}\label{sec:exp-coulomb}

We now consider the Coulomb cost $c(x,y) = 1/|x-y|$ from electronic structure theory in which the transport cost is \emph{singular} and \emph{repulsive}. In particular, (i) this cost is singular on the diagonal; (ii) it is convex and decreasing in $|x-y|$ (repulsive), rather than concave or increasing; (iii) the mixed partial $D^2_{xy}c > 0$, which is the wrong sign for the standard twist condition; and (iv) the transport is \emph{self-transport} ($\mu = \nu$), a degenerate setting for classical uniqueness results. This experiment demonstrates the SQP framework's ability to handle costs that violate essentially every structural assumption of Brenier--Gangbo--McCann theory.   

\paragraph{\it Physical motivation.}
In the strong-interaction limit of Density Functional Theory (DFT), the ground-state energy of $N$ electrons with one-body density $\rho$ is determined by a multi-marginal optimal transport problem with Coulomb repulsive cost~\citep{Seidl1999,ButCDP2012,ColomboDeP2015}. For $N = 2$ electrons in one dimension, this reduces to finding a map $f$ pushing $\rho/2$ to itself (self-transport) that minimizes the Coulomb repulsion energy
\begin{equation}\label{eq:coulomb-sce}
    \min_{f:\; f_\#(\rho/2) = \rho/2} \int \frac{1}{|x - f(x)|}\,\frac{\rho(x)}{2}\,dx.
\end{equation}
The exact solution is the \emph{co-motion function} of Seidl~\citep{Seidl1999}:
\begin{equation}\label{eq:comotion}
    f(x) = F^{-1}(1 - F(x)),
\end{equation}
where $F$ is the CDF of $\rho$ \zihe{$\rho$ or $\rho/2$}. This is the anti-quantile coupling: a monotone decreasing map that sends each electron to its ``opposite'' in the density distribution. Physically, for a diatomic molecule with atoms at positions $\pm d/2$, the co-motion function maps electrons on atom~A to atom~B and vice versa.

\paragraph{\it Why Gangbo--McCann does not apply.}
The cost $c(x,y) = 1/|x-y|$ violates the hypotheses of the Gangbo--McCann existence and uniqueness theory~\citep{GangboMcCann1996} in multiple ways. First, the cost is singular: $c(x,y) \to +\infty$ as $x \to y$, so it is not continuous on the diagonal. Second, $c$ is convex and decreasing in $|x-y|$, the opposite regime from the convex costs $c(x,y) = h(x-y)$ with $h$ strictly convex for which Gangbo--McCann guarantees optimal maps are monotone increasing. Third, the mixed partial derivative
\[
    D^2_{xy} c(x,y) = \frac{2}{|x-y|^3} > 0,\quad x\neq y
\]
has the wrong sign for the standard twist condition \zihe{The standard twist condition only asks $y\mapsto c_x(x,y)$ to be injective\citep[Definition 1.16]{Santambrogio2015}}, which requires $D^2_{xy}c < 0$ to ensure the optimal map is monotone increasing. The positive sign instead produces a ``reverse twist'' consistent with a monotone \emph{decreasing} optimal map. Fourth, the self-transport setting $\mu = \nu$ is degenerate. Specifically, the identity map is always feasible but infinitely costly under the Coulomb cost, so the problem has a fundamentally different character from the generic $\mu \neq \nu$ setting.

\paragraph{\it Problem setup.}
We model a one-dimensional diatomic molecule with density
\[
    \rho(x) \propto \exp\!\Big(\!-\!\Big(\frac{x - d/2}{\sigma}\Big)^{\!2}\Big) + \exp\!\Big(\!-\!\Big(\frac{x + d/2}{\sigma}\Big)^{\!2}\Big),
\]
a symmetric bimodal Gaussian with atom separation~$d$ and width~$\sigma$. We consider three configurations of increasing difficulty:
\begin{enumerate}
    \item \emph{Wide molecule}: $d = 5.0$, $\sigma = 0.3$ (well-separated atoms, easiest);
    \item \emph{Medium molecule}: $d = 3.5$, $\sigma = 0.4$;
    \item \emph{Tight molecule}: $d = 2.5$, $\sigma = 0.4$ (overlapping tails, hardest).
\end{enumerate}
\zihe{Comment to the next line:
For a symmetric density, the anti-quantile map
\[
f(x)=F^{-1}(1-F(x))
\]
has a fixed point at the median \(x_0\), since \(F(x_0)=1/2\) implies
\[
f(x_0)=F^{-1}(1-F(x_0))=F^{-1}(1/2)=x_0.
\]
Hence one cannot claim \(\min_x |x-f(x)|>0\) on the full support.}
\zihe{Although one does not have \(\min_x |x-f(x)|>0\) on the full support, the anti-quantile map in the symmetric bimodal setting satisfies
\[
|x-f(x)|>0
\quad\text{for almost every }x,
\]
with equality only at the median \(x_0\), where \(f(x_0)=x_0\). Thus the map transports almost all mass away from the diagonal, pairing the two modes against each other. This avoids the Coulomb singularity almost everywhere.}
The bimodal structure ensures that $\min_x |x - f(x)| > 0$ on the support of $\rho$, so that the Coulomb singularity is avoided by the optimal map. As the atom separation~$d$ decreases, the minimum gap $|x - f(x)|$ shrinks and the problem becomes numerically more challenging.

We use the smoothed cost $c_\varepsilon(x,y) = 1/(|x-y| + \varepsilon)$ with $\varepsilon = 10^{-3}$ to regularize the singularity.

\paragraph{\it Solver configuration.}
We parameterize the transport map as a piecewise linear function with $n = 400$ pieces ($n + 1 = 401$ parameters) on the domain $[-6, 6]$, without imposing a monotonicity constraint. The pushforward constraint $f_\#(\rho/2) = \rho/2$ is enforced via a histogram discretization with $m = 200$ bins. This is a substantially finer discretization than the earlier experiments ($n = 40$, $m = 20$), motivated by the sharp transitions in the co-motion function between the two atomic modes.

We solve via the SQP method of Section~\ref{subsec:discr-SQP} with Armijo line search on the merit function $M = \mathcal{L} + (\sigma/2)\|\tau\|^2$. The objective gradient and Hessian diagonal are computed analytically using the closed-form derivatives of the Coulomb cost (Theorems~\ref{thm: J1deriv}--\ref{thm: J2deriv}), and the constraint Jacobian is computed via batched finite differences on the GPU. The initial merit penalty is set to $\sigma_0 = 200$ (higher than in other experiments) to reflect the large cost magnitudes near the singularity. Up to 120 SQP iterations are performed with tolerance $\|\tau\|_\infty < 5 \times 10^{-3}$.

\paragraph{\it Multiple initializations.}
We run the solver from four distinct starting maps, each constructed via quantile functions of $\rho$:
\begin{enumerate}
    \item \emph{Anti-quantile} (monotone decreasing): $f_0(x) = F^{-1}(1 - F(x))$, the exact co-motion function evaluated on the PL grid.
    \item \emph{Quantile} (monotone increasing): $f_0(x) = F^{-1}(F(x)) = x$, the identity.
    \item \emph{Tent} (non-monotone): $f_0(x) = F^{-1}(1 - |2F(x) - 1|)$.
    \item \emph{V-shape} (non-monotone): $f_0(x) = F^{-1}(|2F(x) - 1|)$.
\end{enumerate}

\paragraph{\it Results.}

\begin{table}[t]
    \centering
    \caption{Coulomb self-transport $c(x,y) = 1/|x-y|$ for three diatomic molecule configurations, solved by SQP with $n = 400$ PL pieces and $m = 200$ histogram bins. The anti-quantile initialization recovers the exact co-motion function in all cases, with relative errors below $0.05\%$ for the medium and wide molecules.}
    \label{tab:coulomb-results}
    \begin{tabular}{lcccc}
        \toprule
        Density & Best init. & $J(f)$ & $J_{\mathrm{exact}}$ & Rel.\ error \\
        \midrule
        Tight ($d{=}2.5$)  & anti-quantile & $0.2125$ & $0.2126$ & $0.05\%$ \\
        Medium ($d{=}3.5$) & anti-quantile & $0.1469$ & $0.1469$ & $< 0.01\%$ \\
        Wide ($d{=}5.0$)   & anti-quantile & $0.1007$ & $0.1007$ & $< 0.01\%$ \\
        \bottomrule
    \end{tabular}
\end{table}

Table~\ref{tab:coulomb-results} reports the results across the three molecule configurations. Several observations are noteworthy.

First, the anti-quantile initialization recovers the exact co-motion function to high precision in all three cases. For the medium and wide molecules, the SQP objective matches the exact value to four or five significant figures; for the tight molecule, the relative error is $0.05\%$. 

Second, the solver recovers the monotone decreasing structure of the co-motion function without any monotonicity being imposed \emph{a priori}; the decreasing character emerges from the optimization. The other initializations (quantile, tent, V-shape) all converge to local optima with substantially higher objective values (Figure~\ref{fig:coulomb-maps}). The quantile (identity) map is particularly poor, with objectives $2$--$6\times$ worse than the anti-quantile, reflecting the Coulomb cost's strong preference for large displacement.

Third, the convergence plots (Figure~\ref{fig:coulomb-convergence}) show that the anti-quantile initialization converges to feasibility ($\|\tau\|_\infty < 5 \times 10^{-3}$) within the first few SQP iterations for all three molecule configurations, reaching constraint violations of order $10^{-3}$--$10^{-4}$. 
%This contrasts with the augmented Lagrangian, where the constraint violation stagnated at ${\sim}10^{-1}$ for the same problem. 
The other initializations reach higher constraint violations ($10^{-2}$--$10^{-1}$) but still show active progress, with the tent and V-shape initializations gradually approaching feasibility over 120 iterations.

Fourth, the cost landscape visualization (Figure~\ref{fig:coulomb-cost-maps}) reveals the structure of the problem: the Coulomb cost matrix has a bright ridge along the diagonal $x = y$ (where $c \to \infty$), and the optimal map traces a path through the low-cost region far from this ridge. The SQP solution (cyan) overlays the exact co-motion function (white dashed) nearly exactly across the support of $\rho$. \harsha{Yeah, but the experiment started at the optimal solution...}

\paragraph{\it Cost analysis.}
Figure~\ref{fig:coulomb-analysis} compares the Coulomb cost to the quadratic and concave costs used in earlier experiments. The left panel shows $c$ as a function of distance: unlike the quadratic cost (which grows with distance) or the concave cost (which rewards distance), the Coulomb cost is convex and decreasing---it penalizes proximity. The right panel plots the mixed partial $D^2_{xy}c$, which is singular at $x = y$ with large magnitude of both signs due to the non-differentiability of $|x-y|$ at the diagonal. Away from the singularity, the analytical formula $c_{xy} = 2/|x-y|^3 > 0$ confirms the ``reverse twist'': the positive sign is consistent with the anti-quantile (monotone decreasing) structure of the optimal map.

\paragraph{\it Discussion.}
This experiment demonstrates the SQP framework's applicability to a physically important class of optimal transport problems that lies entirely outside classical theory. The Coulomb cost combines three features that individually pose challenges for existing methods (singularity, repulsiveness, and self-transport), and yet the SQP solver recovers the correct co-motion function from a suitable initialization, with relative errors below $0.05\%$.

%The improvement from the augmented Lagrangian to the SQP solver is dramatic: constraint violations decrease from $10^{-1}$ (stagnation) to $10^{-3}$--$10^{-4}$ (convergence), and relative errors in the objective decrease from $0.8\%$--$6.4\%$ to below $0.05\%$. This improvement comes from two sources: the finer discretization ($n = 400$ vs.\ $40$), which resolves the sharp transitions in the co-motion function; and the SQP solver, which drives feasibility and optimality simultaneously through the KKT system rather than sequentially through penalty updates.

The connection to electronic structure theory is interesting since the strictly correlated electrons (SCE) formalism~\citep{Seidl1999,ButCDP2012} provides exact exchange-correlation functionals for DFT in the strong-interaction limit, and efficient computation of co-motion functions is an active area of research~\citep{ColomboDeP2015}. The SQP framework offers a direct approach to these problems that does not require the specialized structure exploited by existing SCE solvers (such as the ODE-based method for 1D problems), and extends naturally to the multi-marginal setting ($N > 2$ electrons) where analytical solutions are unavailable.

\begin{figure}[t]
    \centering
    \includegraphics[width=\textwidth]{graphs/coulumb/fig_coulomb_1_maps.png}
    \caption{Top row: Bimodal electron densities for three diatomic molecule configurations (tight, medium, wide). Bottom row: Transport maps from four initializations. The anti-quantile initialization (pink, starred) recovers the exact co-motion function (black dashed) in all cases, with the SQP objective matching the exact value to four or five significant figures. The other initializations converge to distinct local optima with substantially higher objective values. Brackets~$[\,!\,]$ indicate solutions that did not meet the feasibility tolerance.}
    \label{fig:coulomb-maps}
\end{figure}

\begin{figure}[t]
    \centering
    \includegraphics[width=\textwidth]{graphs/coulumb/fig_coulomb_2_cost_maps.png}
    \caption{Coulomb cost matrix $c(x,y) = 1/|x-y|$ (clipped at the 95th percentile for visibility) overlaid with the exact co-motion function (white dashed) and the SQP solution (cyan). The bright ridge along the diagonal represents the Coulomb singularity; the optimal map avoids this region by mapping each atomic mode to the opposite mode.}
    \label{fig:coulomb-cost-maps}
\end{figure}

\begin{figure}[t]
    \centering
    \includegraphics[width=\textwidth]{graphs/coulumb/fig_coulomb_3_convergence.png}
    \caption{Constraint violation $\|\tau\|_\infty$ versus SQP iteration for all initializations and molecule configurations. The anti-quantile initialization (pink) converges to ${\sim}10^{-3}$--$10^{-4}$ within the first few iterations, well below the tolerance (red dotted line). This is a marked improvement over the augmented Lagrangian, which stagnated at ${\sim}10^{-1}$ for the same problem.}
    \label{fig:coulomb-convergence}
\end{figure}

\begin{figure}[t]
    \centering
    \includegraphics[width=\textwidth]{graphs/coulumb/fig_coulomb_4_analysis.png}
    \caption{Left: Comparison of cost functions as a function of distance $|x-y|$. The Coulomb cost (red) is convex and decreasing, unlike the quadratic (blue) or concave (green) costs. Right: Mixed partial $D^2_{xy}c$ computed via finite differences on the smoothed cost, showing the singularity at $x = y$. The analytical formula $c_{xy} = 2/|x-y|^3 > 0$ confirms the ``reverse twist''---the wrong sign for the standard Gangbo--McCann condition.}
    \label{fig:coulomb-analysis}
\end{figure}

\section{Conclusions}
This paper developed a sequential quadratic programming (SQP) framework for the Monge optimal transport problem under general cost functions. Working directly in the infinite-dimensional space of transport maps, we derived explicit first- and second-order Fr\'echet derivatives for both the Monge objective functional and the Perron--Frobenius pushforward constraint operator. These variational formulas served as the foundation for a functional SQP iteration, for which we established local quadratic convergence by verifying the conditions of \cite{alt1990lagrange}, and global convergence to critical points of a merit functional via an Armijo line search.

A key structural insight of our analysis concerns the local geometry of the feasible set: injectivity of the first-variation operator of the pushforward constraint implies that feasible points are isolated in any finite-dimensional approximation space, yet this isolation breaks down in infinite dimensions due to the possible vanishing of the minimum modulus. This distinction has direct algorithmic consequences, as it necessitates methods---such as SQP---that can navigate infeasible regions while systematically restoring feasibility.

To bridge the gap between the infinite-dimensional formulation and computation, we introduced a finite-dimensional computational problem obtained by restricting transport maps to a parametric family and discretizing the pushforward constraint on a finite grid. We proved that the optimal objective value of the computational problem converges to the true optimal value as both the approximation dimension and the grid resolution are refined, with explicit rates of $O(1/n + 1/m)$. Numerical experiments on one-dimensional problems---including quadratic and concave transport costs---validated the framework and demonstrated close agreement with known analytical solutions.



\subsection{Limitations and Future Directions}

Several important directions remain open.

\paragraph{\bf Extension to higher dimensions.}
The most significant limitation of the present work is the restriction to one-dimensional domains, $\mathcal{X}, \mathcal{Y} \subset \mathbb{R}$. Extending the framework to $\mathbb{R}^d$ with $d \geq 2$ poses fundamental challenges. In one dimension, the pushforward constraint admits a convenient pointwise representation through the change-of-variables formula involving $s^{-1}$ and its derivative, which we exploit for both the theoretical analysis and the grid-based discretization. In higher dimensions, the analogous object is the Jacobian determinant of the transport map, and the Perron--Frobenius operator takes the form
\[
Pf(s)(y) = \frac{f(s^{-1}(y))}{|\det \nabla s(s^{-1}(y))|},
\]
which is well-defined only when $s$ is a diffeomorphism. Computing Fr\'echet derivatives of this operator with respect to $s$ requires differentiating through the Jacobian determinant and the inverse map simultaneously, leading to substantially more complex expressions. Moreover, the characteristic-function-based derivation of the first and second variations (Theorems~4 and~5) relies on one-dimensional Fourier analysis and does not directly generalize to multivariate settings. Developing analogous variational formulas and verifying the SQP regularity conditions in $\mathbb{R}^d$ is a challenging but important direction for future work.

\paragraph{\bf Richer parameterizations and scalability.}
Our computational experiments employed piecewise linear and low-degree polynomial parameterizations, which suffice for one-dimensional problems but may be inadequate for higher-dimensional or large-scale applications. Neural network parameterizations of transport maps, as used in normalizing flows and related generative models, offer far greater expressiveness but introduce additional nonconvexity through the parameterization itself. Integrating SQP-type constraint enforcement with neural network training, potentially through augmented Lagrangian or penalty methods that approximate the SQP structure, is a promising avenue for making the framework applicable to modern machine learning settings.

\paragraph{\bf Constraint qualification and degeneracy.} \harsha{perhaps we should remove this...might invite revision work...}
Our convergence results rely on regularity conditions such as LICQ and second-order sufficiency at the solution. While Proposition~8 shows that LICQ can be ensured by appropriate grid placement when $n > m$, a more thorough understanding of when second-order sufficient conditions hold---or fail---for specific cost functions and distribution pairs would strengthen the practical applicability of the framework. In particular, characterizing the landscape of local optima for nonconvex Monge problems, and understanding the basins of attraction of the SQP iterates, remain largely unexplored.

% \paragraph{\bf Convergence rates for the SQP iterates.}
% While we established convergence rates for the \emph{objective values} of the computational approximation ($O(1/n + 1/m)$), the convergence of the transport maps themselves---in appropriate function space norms---is only guaranteed under strict convexity of the cost. Establishing map-level convergence rates under weaker assumptions, or quantifying the interaction between SQP iteration count, approximation dimension $n$, and grid resolution $m$, would provide a more complete computational complexity picture.

\paragraph{\bf Discrete and singular source measures.}
Many applications involve empirical or discrete source measures $\mu = \frac{1}{N}\sum_{i=1}^N \delta_{x_i}$, for which the absolute continuity assumption underlying our formulation does not hold. Adapting the SQP framework to handle singular measures---potentially through regularization, semi-discrete formulations, or relaxation of the pushforward constraint---would significantly broaden the scope of applicability.

\clearpage
% References here (outcomment the appropriate case)
\bibliographystyle{unsrtnat}
\bibliography{lib}
\newpage
\appendix
\section{Regularities of Pushforward Density}\label{sec:pf-regularity}
Our consistency analysis in Section~\ref{sec:cp} relies on pointwise constraint discretization. To guarantee that solutions to this discretized problem converge to the true continuous solution, the pushforward density $Pf(s)$ must be continuous. This appendix proves that $Pf(s)$ admits a continuous representative, providing the justification for our grid-based approximation.
Recall that
\begin{align}
    Pf(s)(x) := \frac{d}{dx}\int_{s^{-1}((-\infty, x])} f(y)\,dy .
    \label{eq:Pf-def}
\end{align}
Let $\mu_s$ be the pushforward of $\mu$ under $s$, i.e.,
\[
\mu_s(A)=\mu\big(s^{-1}(A)\big), \qquad \forall\text{ measurable sets }  A\subseteq \mathrm{supp}(g).
\]
Let $\mu_L$ denote the Lebesgue measure on $\mathrm{supp}(g)$.
Whenever $\mu_s\ll \mu_L$, the Radon--Nikodym theorem guarantees the existence of
a $\mu_L$--a.e.\ unique density $\frac{d\mu_s}{d\mu_L}\in L^1(\mu_L)$ satisfying
\[
\mu_s(A)=\int_A \frac{d\mu_s}{d\mu_L}(x)\,d\mu_L(x),\qquad \forall\text{ measurable sets }  A\subseteq \mathrm{supp}(g).
\]

A measurable function $u:\mathrm{supp}(g)\to\mathbb R$ is called a \emph{version} of
$\frac{d\mu_s}{d\mu_L}$ if $u=\frac{d\mu_s}{d\mu_L}$ holds $\mu_L$--a.e.
In particular, whenever the derivative in~\eqref{eq:Pf-def} exists, $Pf(s)$ can be
viewed as a version of $\frac{d\mu_s}{d\mu_L}$ (see, e.g., \cite[Ch.\ 4]{boyarsky2012laws}).
A \emph{continuous version} of $\frac{d\mu_s}{d\mu_L}$ means a continuous function on
$\mathrm{supp}(g)$ that agrees with $\frac{d\mu_s}{d\mu_L}$ $\mu_L$--a.e.

We will use the following likelihood-ratio criterion to obtain a continuous version.
Fix $x_0\in \mathrm{supp}(g)$ and let $\mathcal{U}(x_0)$ be the collection of open neighborhoods of $x_0$.
For any nested sequence $A_1\supset A_2\supset \cdots$ with each $A_n\in\mathcal{U}(x_0)$ and
$\bigcap_{n\ge1}A_n=\{x_0\}$, define the likelihood ratio (when it exists) by
\begin{align}
    \lim_{n\to\infty}\frac{\mu_s(A_n)}{\mu_L(A_n)}.
    \label{eq:likelihood-def}
\end{align}

The following theorem of \cite{piccioni1983continuous} links the pointwise existence of
likelihood ratios to the existence of a continuous Radon--Nikodym representative.

\begin{theorem}[\cite{piccioni1983continuous}]\label{thm:likelihood-ratio}
If for every $x_0\in \mathrm{supp}(g)$ the limit~\eqref{eq:likelihood-def} exists (for all nested neighborhood bases as above),
then $\frac{d\mu_s}{d\mu_L}$ admits a continuous version $\widetilde p$ on $\mathrm{supp}(g)$.
Moreover, $\widetilde p(x_0)$ equals the likelihood ratio at $x_0$ for every $x_0\in \mathrm{supp}(g)$.
\end{theorem}

\begin{lemma}\label{lem:pf-continuous-version}
Under Assumption~\ref{ass:density-cont} and Assumption~\ref{ass:mono},
the function $Pf(s)$ can be identified as a continuous version of $\frac{d\mu_s}{d\mu_L}$.
\end{lemma}

\begin{proof}
Fix any $x_0\in \mathrm{supp}(g)$. It suffices, by Theorem~\ref{thm:likelihood-ratio},
to show that the likelihood ratio exists at $x_0$ and can be identified pointwise.
For $r>0$, set $A_r:=(x_0-r,x_0+r)\cap \mathrm{supp}(g)$.
Then, using the definition of pushforward,
\begin{align*}
\frac{\mu_s(A_r)}{\mu_L(A_r)}
&= \frac{\mu\big(s^{-1}(A_r)\big)}{\mu_L(A_r)}
 = \frac{\int_{\mathrm{supp}(f)} \mathbbm{1}\{y:\, x_0-r<s(y)<x_0+r\}\, f(y)\,dy}{\mu_L(A_r)} .
\end{align*}
Since $\mu_L(A_r)=2r$ for all sufficiently small $r$ (because $\mathrm{supp}(g)$ is an open interval),
we have for small $r$,
\begin{align}
\frac{\mu_s(A_r)}{\mu_L(A_r)}
= \frac{1}{2r}\int_{\mathrm{supp}(f)} \mathbbm{1}\{y:\, x_0-r<s(y)<x_0+r\}\, f(y)\,dy .
\label{eq:likelihood-ball}
\end{align}

By Assumption~\ref{ass:mono}, $s:\mathrm{supp}(f)\to\mathrm{supp}(g)$ is a bijection that is strictly monotone.
Hence $s$ is continuous and admits a continuous inverse $s^{-1}$ on $\mathrm{supp}(g)$.
Moreover, $s^{-1}$ is monotone and therefore has finite left and right derivatives everywhere, and is differentiable a.e.\ on $\mathrm{supp}(g)$.

Assume without loss of generality that $s$ is strictly increasing (the decreasing case is analogous).
Then for all sufficiently small $r>0$,
\[
\{y:\, x_0-r<s(y)<x_0+r\} = \big(s^{-1}(x_0-r),\, s^{-1}(x_0+r)\big),
\]
and~\eqref{eq:likelihood-ball} becomes
\begin{align}
\frac{\mu_s(A_r)}{\mu_L(A_r)}
= \frac{1}{2r}\int_{s^{-1}(x_0-r)}^{s^{-1}(x_0+r)} f(y)\,dy .
\label{eq:likelihood-interval}
\end{align}

We first consider points $x_0$ at which $s^{-1}$ is differentiable.
Define
\[
F(r):=\int_{s^{-1}(x_0)}^{s^{-1}(x_0+r)} f(y)\,dy,
\qquad r \ \text{in a neighborhood of }0.
\]
Then
\[
\int_{s^{-1}(x_0-r)}^{s^{-1}(x_0+r)} f(y)\,dy
= F(r)-F(-r),
\]
so
\[
\frac{1}{2r}\int_{s^{-1}(x_0-r)}^{s^{-1}(x_0+r)} f(y)\,dy
= \frac{F(r)-F(-r)}{2r}.
\]
By the Leibniz rule and continuity of $f$ (Assumption~\ref{ass:density-cont}),
$F$ is differentiable at $0$ with
\[
F'(0)= f\big(s^{-1}(x_0)\big)\,(s^{-1})'(x_0).
\]
Therefore,
\[
\lim_{r\rightarrow0}\frac{F(r)-F(-r)}{2r} = F'(0)= f\big(s^{-1}(x_0)\big)\,(s^{-1})'(x_0),
\]
which shows the likelihood ratio exists at such $x_0$.

At points $x_0$ where $s^{-1}$ is not differentiable, monotonicity still guarantees the existence
of the one-sided derivatives $\partial^- s^{-1}(x_0)$ and $\partial^+ s^{-1}(x_0)$.
Applying the Leibniz rule in one-sided form yields
\[
\partial^- F(0)= f\big(s^{-1}(x_0)\big)\,\partial^- s^{-1}(x_0),
\qquad
\partial^+ F(0)= f\big(s^{-1}(x_0)\big)\,\partial^+ s^{-1}(x_0),
\]
and hence
\[
\lim_{r\rightarrow 0}\frac{F(r)-F(-r)}{2r}
= \frac{1}{2}\partial^-F(0)+\frac{1}{2}\partial^+F(0)
= \frac{f(s^{-1}(x_0))}{2}\Big(\partial^- s^{-1}(x_0)+\partial^+ s^{-1}(x_0)\Big),
\]
so the likelihood ratio still exists for every $x_0\in \mathrm{supp}(g)$.

Thus, for each $x_0\in \mathrm{supp}(g)$ the likelihood ratio
$\lim_{r\rightarrow0}\mu_s(A_r)/\mu_L(A_r)$ exists.
By Theorem~\ref{thm:likelihood-ratio}, $\frac{d\mu_s}{d\mu_L}$ admits a continuous version on $\mathrm{supp}(g)$.
Since $Pf(s)$ is a version of $\frac{d\mu_s}{d\mu_L}$,
we may identify $Pf(s)$ with this continuous version, completing the proof.
\end{proof}


\section{Proofs from Section\ref{sec:pre}}\label{proof:pre}

\subsection{Proof of Theorem~\ref{thm: J1deriv}}
Recall that in Theorem~\ref{thm: J1deriv} we define
\[
J_s(h)=\int_{\mathcal{X}} h(x)\,D_y c\big(x,s(x)\big)\,f(x)dx .
\]
Then
\begin{align*}
    &\left|J(s+h)-(J(s)+J_s(h))\right|\\
    &\quad\quad=\Bigg|\int_{\mathcal{X}} c\big(x,s(x)+h(x)\big)f(x)dx
          -\int_{\mathcal{X}} c\big(x,s(x)\big)f(x)dx
          -\int_{\mathcal{X}} h(x)\,D_y c\big(x,s(x)\big)f(x)dx\Bigg|.
\end{align*}
By the fundamental theorem of calculus, for each fixed $x\in\mathcal{X}$,
\begin{align*}
c\big(x,s(x)+h(x)\big)
&=c\big(x,s(x)\big)+\int_0^1 D_y c\big(x,s(x)+\tilde{\epsilon} h(x)\big)\,h(x)\,d\tilde{\epsilon}\\
&=c\big(x,s(x)\big)+D_y c\big(x,s(x)\big)h(x)
    +\int_0^1\!\Big(D_y c\big(x,s(x)+\tilde{\epsilon} h(x)\big)-D_y c\big(x,s(x)\big)\Big)h(x)\,d\tilde{\epsilon} .
\end{align*}
Hence,
\begin{align*}
\left|c\big(x,s(x)+h(x)\big)-c\big(x,s(x)\big)-D_y c\big(x,s(x)\big)h(x)\right|
&\le \int_0^1 \left|D_y c\big(x,s(x)+\tilde{\epsilon} h(x)\big)-D_y c\big(x,s(x)\big)\right|\,|h(x)|\,d\tilde{\epsilon}\\
&\le \int_0^1 L\,|\tilde{\epsilon} h(x)|\,|h(x)|\,d\tilde{\epsilon} \qquad \text{(by Assumption~\ref{ass:cost-second})}\\
&=\frac{L}{2}|h(x)|^2 .
\end{align*}
Integrating against $f(x)dx$ yields
\begin{align*}
    \left|J(s+h)-(J(s)+J_s(h))\right|
    &\le \int_{\mathcal{X}}\frac{L}{2}|h(x)|^2 f(x)dx
    \le \frac{L}{2}\|h\|
    =o(\|h\|),
\end{align*}
which completes the proof by the definition of the Fr\'echet derivative.

\subsection{Proof of Theorem~\ref{thm: J2deriv}}
Fix $h,\Tilde h\in\mathcal{S}$. Recall
\[
J_s(h)=\int_{\mathcal{X}} h(x)\,D_y c\big(x,s(x)\big)\,f(x)dx,
\qquad
H_s(h)(\Tilde h)=\int_{\mathcal{X}}\Tilde h(x)h(x)\,D^2_{yy}c\big(x,s(x)\big)\,f(x)dx .
\]
We consider
\begin{align*}
\left|J_{s+\epsilon\Tilde h}(h)-\big(J_s(h)+\epsilon H_s(h)(\Tilde h)\big)\right|
&=\Bigg|\int_{\mathcal{X}} h(x)\Big(D_y c\big(x,s(x)+\epsilon\Tilde h(x)\big)-D_y c\big(x,s(x)\big)\Big)f(x)dx\\
&\qquad\qquad\qquad\qquad\qquad
-\epsilon\int_{\mathcal{X}}\Tilde h(x)h(x)\,D^2_{yy}c\big(x,s(x)\big)f(x)dx\Bigg|.
\end{align*}
By the fundamental theorem of calculus applied to $y\mapsto D_y c(x,y)$,
\[
D_y c\big(x,s+\epsilon\Tilde h\big)-D_y c(x,s)
=\epsilon\,D^2_{yy}c(x,s)\,\Tilde h
+\int_0^1\epsilon\Big(D^2_{yy}c\big(x,s+\tilde{\epsilon}\epsilon\Tilde h\big)-D^2_{yy}c(x,s)\Big)\Tilde h\,d\tilde{\epsilon} .
\]
Therefore,
\begin{align*}
&\left|J_{s+\epsilon\Tilde h}(h)-\big(J_s(h)+\epsilon H_s(h)(\Tilde h)\big)\right|\\
&\le \int_{\mathcal{X}}|h(x)|\,f(x)\int_0^1 |\epsilon|\,
\left|D^2_{yy}c\big(x,s(x)+\tilde{\epsilon}\epsilon\Tilde h(x)\big)-D^2_{yy}c\big(x,s(x)\big)\right|\,|\Tilde h(x)|\,d\tilde{\epsilon}dx\\
&\le \int_{\mathcal{X}}|h(x)|\,f(x)\int_0^1 |\epsilon|\, M\,|\tilde{\epsilon}\epsilon\Tilde h(x)|\,|\Tilde h(x)|\,d\tilde{\epsilon}dx
\qquad\text{(by Assumption~\ref{ass:cost-second})}\\
&=\frac{M}{2}\epsilon^2\int_{\mathcal{X}}|h(x)|\,|\Tilde h(x)|^2 f(x)dx
=o(|\epsilon|),
\end{align*}
which implies
\[
\frac{1}{\epsilon}\left(J_{s+\epsilon\Tilde h}(h)-J_s(h)\right)-H_s(h)(\Tilde h)\to 0
\quad\text{as }\epsilon\to 0.
\]
Hence $H_s(h)(\cdot)$ is the Fr\'echet derivative of $J_s(h)$ at $s$.

\subsection{Proof of Theorem~\ref{thm: PF1deriv}}
Fix $s\in\mathcal{S}$ and $h\in\mathcal{S}$. For $\epsilon\in\mathbb R$ define
\[
X_\epsilon:=(s+\epsilon h)(Y),\qquad m_\epsilon:=\mathbb E[X_\epsilon],
\qquad Z_\epsilon:=X_\epsilon-m_\epsilon.
\]
Let $\kappa_\epsilon(t):=\mathbb E[e^{itZ_\epsilon}]$ denote the characteristic function of $Z_\epsilon$.
Recall that $Pf(s)(x)$ is the density of $X_0=s(Y)$ on $\mathcal{X}$ and that $g_s(x):=Pf(s)(x)$ is assumed
continuously differentiable with $g_s'\in L^1(\mathcal{X})$.
Let $\hat G_s$ and $\hat g_s$ denote the CDF and pdf of $Z_0=s(Y)-m_0$, respectively.
By definition,
\begin{align*}
Pf_s(h)(x)
:=\lim_{\epsilon\to 0}\frac{Pf(s+\epsilon h)(x)-Pf(s)(x)}{\epsilon}
=\lim_{\epsilon\to 0}\frac{1}{\epsilon}\frac{d}{dx}\Big(\mathbb P(X_\epsilon\le x)-\mathbb P(X_0\le x)\Big),
\end{align*}
whenever the limit exists.
Write
\[
\mathbb P(X_\epsilon\le x)=\mathbb P(Z_\epsilon\le x-m_\epsilon),
\qquad
\mathbb P(X_0\le x)=\mathbb P(Z_0\le x-m_0).
\]
Add and subtract $\mathbb P(Z_0\le x-m_\epsilon)=\hat G_s(x-m_\epsilon)$ to obtain the identity
\begin{align*}
\mathbb P(X_\epsilon\le x)-\mathbb P(X_0\le x)
&=\Big(\mathbb P(Z_\epsilon\le x-m_\epsilon)-\mathbb P(Z_0\le x-m_\epsilon)\Big)\\
&\quad-\Big(\hat G_s(x-m_\epsilon)-\hat G_s(x-m_0)\Big).
\end{align*}
Therefore,
\begin{align}\label{eq:split_two_terms}
\frac{1}{\epsilon}\frac{d}{dx}\Big(\mathbb P(X_\epsilon\le x)-\mathbb P(X_0\le x)\Big)
&=
T_{1,\epsilon}(x)+T_{2,\epsilon}(x),
\end{align}
where
\begin{align*}
T_{1,\epsilon}(x)
&:=
\frac{1}{\epsilon}\frac{d}{dx}\Big(\mathbb P(Z_\epsilon\le x-m_\epsilon)-\mathbb P(Z_0\le x-m_\epsilon)\Big),\\
T_{2,\epsilon}(x)
&:=
-\frac{1}{\epsilon}\frac{d}{dx}\Big(\hat G_s(x-m_\epsilon)-\hat G_s(x-m_0)\Big).
\end{align*}

By assumption, for all $|\epsilon|<K$ the characteristic functions $\kappa_\epsilon$ and $\kappa_0$ are integrable.
Hence, by \cite[Lemma~3.4.19]{durrett2019probability}, for each fixed $\epsilon$,
\begin{align*}
\mathbb P(Z_\epsilon\le x-m_\epsilon)-\mathbb P(Z_0\le x-m_\epsilon)
=\frac{1}{2\pi}\int_{\mathbb R}-e^{it(x-m_\epsilon)}\frac{\kappa_\epsilon(t)-\kappa_0(t)}{it}\,dt.
\end{align*}
Differentiate with respect to $x$. Since $\kappa_\epsilon,\kappa_0\in L^1(\mathbb R)$ and
\[
\left|\frac{d}{dx}\left(-e^{it(x-m_\epsilon)}\frac{1}{it}\right)\right|
=\left| -e^{it(x-m_\epsilon)}\right|=1,
\]
differentiation under the integral sign is justified and yields
\begin{align*}
\frac{d}{dx}\Big(\mathbb P(Z_\epsilon\le x-m_\epsilon)-\mathbb P(Z_0\le x-m_\epsilon)\Big)
=
\frac{1}{2\pi}\int_{\mathbb R}-e^{it(x-m_\epsilon)}\big(\kappa_\epsilon(t)-\kappa_0(t)\big)\,dt.
\end{align*}
Therefore,
\begin{align}\label{eq:T1_eps}
T_{1,\epsilon}(x)
=\frac{1}{2\pi}\int_{\mathbb R}-e^{itx}\,e^{-itm_\epsilon}\,\frac{\kappa_\epsilon(t)-\kappa_0(t)}{\epsilon}\,dt.
\end{align}
Observe that
\[
Z_\epsilon
= (s(Y)-\mathbb E[s(Y)])+\epsilon\big(h(Y)-\mathbb E[h(Y)]\big).
\]
For each fixed $t\in\mathbb R$, the map $\epsilon\mapsto \kappa_\epsilon(t)=\mathbb E[e^{itZ_\epsilon}]$ is
differentiable, with
\[
\frac{d}{d\epsilon}\kappa_\epsilon(t)
=
\mathbb E\!\left[it\,e^{itZ_\epsilon}\big(h(Y)-\mathbb E[h(Y)]\big)\right].
\]
In particular,
\begin{align}\label{eq:kappa_derivative_at0}
\lim_{\epsilon\to 0}\frac{\kappa_\epsilon(t)-\kappa_0(t)}{\epsilon}
=\frac{d}{d\epsilon}\kappa_\epsilon(t)\Big|_{\epsilon=0}
=
\mathbb E\!\left[it\,e^{it(s(Y)-\mathbb E[s(Y)])}\big(h(Y)-\mathbb E[h(Y)]\big)\right].
\end{align}
Moreover, by Assumption~\eqref{eq:M1_dom}, for all $0<|\epsilon|<K$,
\[
\left|e^{-itm_\epsilon}\frac{\kappa_\epsilon(t)-\kappa_0(t)}{\epsilon}\right|
=\left|\frac{\kappa_\epsilon(t)-\kappa_0(t)}{\epsilon}\right|
\le
|t|\left|\mathbb E\!\left[e^{itZ_\epsilon}\big(h(Y)-\mathbb E[h(Y)]\big)\right]\right|
\le M_1(t)\,\|h\|.
\]
Since $M_1\in L^1(\mathbb R)$, dominated convergence applies to the right-hand side of
\eqref{eq:T1_eps}, giving for each $x\in\mathcal{X}$,
\begin{align}\label{eq:T1_limit}
\lim_{\epsilon\to 0}T_{1,\epsilon}(x)
=
\frac{1}{2\pi}\int_{\mathbb R}-e^{itx}\,e^{-it\mathbb E[s(Y)]}\,
\mathbb E\!\left[it\,e^{it(s(Y)-\mathbb E[s(Y)])}\big(h(Y)-\mathbb E[h(Y)]\big)\right]\,dt.
\end{align}
Since $\hat G_s$ is differentiable with derivative $\hat g_s$, we have
\[
\frac{d}{dx}\hat G_s(x-m_\epsilon)=\hat g_s(x-m_\epsilon),
\qquad
\frac{d}{dx}\hat G_s(x-\mathbb E[X_0])=\hat g_s(x-\mathbb E[X_0]).
\]
Therefore,
\[
T_{2,\epsilon}(x)
=
-\frac{1}{\epsilon}\Big(\hat g_s(x-m_\epsilon)-\hat g_s(x-\mathbb E[X_0])\Big).
\]
Because $\hat g_s$ is continuously differentiable, a first-order Taylor expansion gives
\[
\hat g_s(x-m_\epsilon)
=
\hat g_s(x-\mathbb E[X_0])-\hat g_s'(x-\mathbb E[X_0])\,(m_\epsilon-\mathbb E[X_0])
+o(|m_\epsilon-\mathbb E[X_0]|).
\]
Since $m_\epsilon-\mathbb E[X_0]=\epsilon\,\mathbb E[h(Y)]$, we obtain
\begin{align}\label{eq:T2_limit}
\lim_{\epsilon\to 0}T_{2,\epsilon}(x)
=
-\hat g_s'(x-\mathbb E[X_0])\,\mathbb E[h(Y)].
\end{align}
Finally, $\hat g_s(u)=g_s(u+\mathbb E[s(Y)])$ implies
$\hat g_s'(x-\mathbb E[X_0])=g_s'(x)$, hence the limit in \eqref{eq:T2_limit} equals $-g_s'(x)\mathbb E[h(Y)]$.
Combining \eqref{eq:split_two_terms}, \eqref{eq:T1_limit}, and \eqref{eq:T2_limit} yields, for each $x\in\mathcal{X}$,
\begin{align*}
\lim_{\epsilon\to 0}\frac{1}{\epsilon}\frac{d}{dx}\Big(\mathbb P(X_\epsilon\le x)-\mathbb P(X_0\le x)\Big)
=&\;\frac{1}{2\pi}\int_{\mathbb R}-e^{itx}\,e^{-it\mathbb E[s(Y)]}\,
\mathbb E\!\left[it\,e^{it(s(Y)-\mathbb E[s(Y)])}\big(h(Y)-\mathbb E[h(Y)]\big)\right]\,dt\\
&\;-\;g_s'(x)\,\mathbb E[h(Y)].
\end{align*}
This defines the candidate derivative $Pf_s(h)(x)$.

To show $Pf_s\in\mathcal{B}(\mathcal{S},\mathcal{Z}_0)$, use the inequality
\[
\sup_{x\in\mathcal{X}}\left|\int_{\mathbb R} e^{itx}F(t)\,dt\right|\le \int_{\mathbb R}|F(t)|\,dt
\qquad (F\in L^1(\mathbb R)).
\]
With
\[
F(t):=
\mathbb E\!\left[it\,e^{it(s(Y)-\mathbb E[s(Y)])}\big(h(Y)-\mathbb E[h(Y)]\big)\right],
\]
Assumption~\eqref{eq:M1_dom} (with $\epsilon=0$) gives $\int_{\mathbb R}|F(t)|\,dt\le \|M_1\|_{L^1}\|h\|$.
Hence,
\[
\left\|\frac{1}{2\pi}\int_{\mathbb R}-e^{itx}\,e^{-it\mathbb E[s(Y)]}F(t)\,dt\right\|_{L^1(\mathcal{X})}
\le \frac{|\mathcal{X}|}{2\pi}\|M_1\|_{L^1(\mathbb R)}\,\|h\|.
\]
Moreover, since $\mu$ is a probability measure, $|\mathbb E[h(Y)]|\le \|h\|=\|h\|$, so
\[
\|g_s'(\cdot)\mathbb E[h(Y)]\|_{L^1(\mathcal{X})}\le \|g_s'\|_{L^1(\mathcal{X})}\,\|h\|.
\]
Combining the two bounds yields $\|Pf_s(h)\|_{L^1(\mathcal{X})}\le C\|h\|$, proving $Pf_s\in\mathcal{B}(\mathcal{S},\mathcal{Z}_0)$.

Finally, the domination \eqref{eq:M1_dom} provides an $L^1(\mathcal{X})$ envelope independent of $\epsilon$ for the integrand
in \eqref{eq:T1_eps}, and $g_s'\in L^1(\mathcal{X})$ controls the shift term. Therefore the difference quotient
$\epsilon^{-1}(Pf(s+\epsilon h)-Pf(s))$ converges in $\mathcal{Z}=L^1(\mathcal{X})$ to $Pf_s(h)$, establishing
Fr\'echet differentiability at $s$ with derivative $Pf'(s)=Pf_s$.

\begin{comment}
  \subsection{Proof of Corollary~\ref{cor: PF1deriv}}
    \begin{align*}
    & Pf_s(h)(x)\\ 
    =&\lim_{\epsilon\rightarrow 0} \frac{1}{\epsilon}(Pf(s+\epsilon h)(x)-Pf(s)(x))\\
    =&\lim_{\epsilon\rightarrow 0}  \frac{1}{\epsilon}\frac{d}{dx} \left(\int_{(s+\epsilon h)^{-1}(A)} f(y)dy - \int_{s^{-1}(A)} f(y)dy\right)\\
    =&\lim_{\epsilon\rightarrow 0}  \frac{1}{\epsilon}\frac{d}{dx} \left(\int_{-\infty}^{(s+\epsilon h)^{-1}(x)} f(y)dy - \int_{-\infty}^{s^{-1}(x)}f(y)dy\right)\\
    =&\lim_{\epsilon\rightarrow 0}  \frac{1}{\epsilon}\frac{d}{dx} \left(\int_{s^{-1}(x)}^{(s+\epsilon h)^{-1}(x)} f(y)dy\right)\\
    =&\lim_{\epsilon\rightarrow 0}\left(f((s+\epsilon h)^{-1}(x))\frac{d}{dx}(s+\epsilon h)^{-1}(x)-f(s^{-1}(x))\frac{d}{dx}s^{-1}(x)\right)\\
    =&\frac{d}{dz}\frac{f((s+zh)^{-1}(x))}{(s+zh)'((s+zh)^{-1}(x))}\Big|_{z=0}.
\end{align*}
The penultimate expression is due to the Leibniz rule and the last expression is based on the derivative of an inverse function. Note that the proof only goes through because strictly increasing functions are invertible, and the summation of two strictly increasing functions is also strictly increasing.  
\end{comment}


\subsection{Proof of Theorem~\ref{thm: PF2deriv}}
\begin{proof}
Fix $s\in\mathcal{S}$ and $h,\Tilde h\in\mathcal{S}$. For $\epsilon\in\mathbb R$ define
\[
X_\epsilon := (s+\epsilon\Tilde h)(Y),\qquad m_\epsilon:=\mathbb E[X_\epsilon],\qquad
Z_\epsilon := X_\epsilon-m_\epsilon.
\]
Recall from Theorem~\ref{thm: PF1deriv} (with the corrected Fourier factors) that for each $\epsilon$,
\begin{align}\label{eq:Pf_first_var_eps_pf2}
Pf_{s+\epsilon\Tilde h}(h)(x)
=&\;\frac{1}{2\pi}\int_{\mathbb R}-e^{itx}\,e^{-it\,m_\epsilon}\,
\mathbb E\!\left[
it\,e^{itZ_\epsilon}\big(h(Y)-\mathbb E[h(Y)]\big)
\right]\,dt
-\;g_{s+\epsilon\Tilde h}'(x)\,\mathbb E[h(Y)].
\end{align}
Similarly,
\begin{align}\label{eq:Pf_first_var_0_pf2}
Pf_{s}(h)(x)
=&\;\frac{1}{2\pi}\int_{\mathbb R}-e^{itx}\,e^{-it\,\mathbb E[s(Y)]}\,
\mathbb E\!\left[
it\,e^{it(s(Y)-\mathbb E[s(Y)])}\big(h(Y)-\mathbb E[h(Y)]\big)
\right]\,dt
-\;g_{s}'(x)\,\mathbb E[h(Y)].
\end{align}

We consider the difference quotient
\begin{align}\label{eq:second_var_start_newletters}
\frac{1}{\epsilon}\Big(Pf_{s+\epsilon\Tilde h}(h)(x)-Pf_s(h)(x)\Big)
=\mathsf{A}_\epsilon(x)+\mathsf{B}_\epsilon(x),
\end{align}
where $\mathsf{A}_\epsilon$ collects the Fourier-integral terms and $\mathsf{B}_\epsilon$ collects the density-derivative terms.
Since $g_s$ is twice continuously differentiable and $g_s''\in L^1(\mathcal{X})$, we have
\[
\mathsf{B}_\epsilon(x)
=
\frac{-g'_{s+\epsilon\Tilde h}(x)\mathbb E[h(Y)]+g'_s(x)\mathbb E[h(Y)]}{\epsilon}.
\]
By the same shift argument used in the proof of Theorem~\ref{thm: PF1deriv}, the derivative of the shift contributes
\[
\partial_\epsilon g'_{s+\epsilon\Tilde h}(x)\big|_{\epsilon=0}
= -g_s''(x)\,\mathbb E[\Tilde h(Y)],
\]
and thus
\begin{align}\label{eq:B_limit}
\lim_{\epsilon\to 0}\mathsf{B}_\epsilon(x)
= g_s''(x)\,\mathbb E[h(Y)]\,\mathbb E[\Tilde h(Y)].
\end{align}
Define for each $\epsilon$ and $t$,
\[
\alpha_\epsilon(t):=e^{-it\,m_\epsilon},\qquad
\beta_\epsilon(t):=\mathbb E\!\left[it\,e^{itZ_\epsilon}\big(h(Y)-\mathbb E[h(Y)]\big)\right].
\]
Then the Fourier term in \eqref{eq:Pf_first_var_eps_pf2} equals
\[
\frac{1}{2\pi}\int_{\mathbb R}-e^{itx}\,\alpha_\epsilon(t)\,\beta_\epsilon(t)\,dt,
\]
and therefore
\begin{align}\label{eq:A_eps}
\mathsf{A}_\epsilon(x)
=\frac{1}{2\pi}\int_{\mathbb R}-e^{itx}\,
\frac{\alpha_\epsilon(t)\beta_\epsilon(t)-\alpha_0(t)\beta_0(t)}{\epsilon}\,dt.
\end{align}
For each fixed $t\in\mathbb R$, both $\epsilon\mapsto \alpha_\epsilon(t)$ and $\epsilon\mapsto \beta_\epsilon(t)$ are differentiable.
Since $m_\epsilon=\mathbb E[s(Y)]+\epsilon\,\mathbb E[\Tilde h(Y)]$, we have
\[
\alpha_\epsilon'(t)=\frac{d}{d\epsilon}e^{-it m_\epsilon}=-it\,\mathbb E[\Tilde h(Y)]\,e^{-it m_\epsilon}.
\]
Also, because
\[
Z_\epsilon=(s(Y)-\mathbb E[s(Y)])+\epsilon(\Tilde h(Y)-\mathbb E[\Tilde h(Y)]),
\]
we obtain
\[
\beta_\epsilon'(t)
=\mathbb E\!\left[it\cdot it\,e^{itZ_\epsilon}\big(\Tilde h(Y)-\mathbb E[\Tilde h(Y)]\big)\big(h(Y)-\mathbb E[h(Y)]\big)\right]
=
-\mathbb E\!\left[t^2 e^{itZ_\epsilon}\big(h-\mathbb E h\big)\big(\Tilde h-\mathbb E\Tilde h\big)\right].
\]
Thus, by the product rule,
\begin{align}\label{eq:alphabeta_prime}
\frac{d}{d\epsilon}\big(\alpha_\epsilon(t)\beta_\epsilon(t)\big)
&=\alpha_\epsilon'(t)\beta_\epsilon(t)+\alpha_\epsilon(t)\beta_\epsilon'(t)\\
&=\;e^{-it m_\epsilon}\,
\mathbb E\Bigg[
e^{itZ_\epsilon}\big(h-\mathbb E h\big)\Big(t^2\big(\Tilde h-\mathbb E\Tilde h\big)+it\,\mathbb E[\Tilde h]\Big)
\Bigg],
\nonumber
\end{align}
where we suppressed the argument $(Y)$ in the last expectation for readability.

Evaluating at $\epsilon=0$ gives
\begin{align}\label{eq:alphabeta_prime_0}
\lim_{\epsilon\to 0}\frac{\alpha_\epsilon(t)\beta_\epsilon(t)-\alpha_0(t)\beta_0(t)}{\epsilon}
=
e^{-it\,\mathbb E[s(Y)]}\,
\mathbb E\Bigg[
e^{it(s(Y)-\mathbb E[s(Y)])}\big(h-\mathbb E h\big)\Big(t^2\big(\Tilde h-\mathbb E\Tilde h\big)+it\,\mathbb E[\Tilde h]\Big)
\Bigg].
\end{align}
Using \eqref{eq:alphabeta_prime} and the triangle inequality, for $0<|\epsilon|<K$,
\begin{align*}
\left|\frac{\alpha_\epsilon(t)\beta_\epsilon(t)-\alpha_0(t)\beta_0(t)}{\epsilon}\right|
&\le
|t|^2\left|\mathbb E\!\left[e^{itZ_\epsilon}(h-\mathbb Eh)(\Tilde h-\mathbb E\Tilde h)\right]\right|
+|t|\,|\mathbb E[\Tilde h]|\left|\mathbb E\!\left[e^{itZ_\epsilon}(h-\mathbb Eh)\right]\right|.
\end{align*}
By Assumption~\eqref{eq:M2_dom}, the first term is bounded by $M_2(t)\|h\|\,\|\Tilde h\|$.
For the second term, the domination hypothesis from Theorem~\ref{thm: PF1deriv} yields
\[
|t|\left|\mathbb E\!\left[e^{itZ_\epsilon}(h-\mathbb Eh)\right]\right|
\le M_1(t)\,\|h\|,
\qquad
|\mathbb E[\Tilde h]|\le \|\Tilde h\|.
\]
Hence,
\[
\left|\frac{\alpha_\epsilon(t)\beta_\epsilon(t)-\alpha_0(t)\beta_0(t)}{\epsilon}\right|
\le (M_2(t)+M_1(t))\,\|h\|\,\|\Tilde h\|.
\]
Since $M_1,M_2\in L^1(\mathbb R)$, dominated convergence applies to \eqref{eq:A_eps}, and using \eqref{eq:alphabeta_prime_0} gives,
for each $x\in\mathcal{X}$,
\begin{align}\label{eq:A_limit}
\lim_{\epsilon\to 0}\mathsf{A}_\epsilon(x)
=
\frac{1}{2\pi}\int_{\mathbb R}-e^{itx}\,e^{-it\,\mathbb E[s(Y)]}\,
\mathbb E\Bigg[
e^{it(s(Y)-\mathbb E[s(Y)])}\big(h-\mathbb E h\big)\Big(t^2\big(\Tilde h-\mathbb E\Tilde h\big)+it\,\mathbb E[\Tilde h]\Big)
\Bigg]\,dt.
\end{align}
Combining \eqref{eq:second_var_start_newletters}, \eqref{eq:A_limit}, and \eqref{eq:B_limit} yields
\begin{align*}
&\lim_{\epsilon\to 0}\frac{Pf_{s+\epsilon\Tilde h}(h)(x)-Pf_s(h)(x)}{\epsilon}
\\
&\quad\quad =\;\frac{1}{2\pi}\int_{\mathbb R} e^{itx} e^{-it\,\mathbb E [s(Y)] }\mathbb E\Bigg[ e^{it(s(Y)-\mathbb E [s(Y)])}
\big(h(Y)-\mathbb E[h(Y)]\big)\bigg(t^2\big(\Tilde{h}(Y)-\mathbb E[\Tilde{h}(Y)]\big)+it\,\mathbb E[\Tilde{h}(Y)]\bigg)\Bigg]dt\\
&\quad\quad+\;g_s''(x)\,\mathbb E [h(Y)] \mathbb E [\Tilde{h}(Y)].
\end{align*}
This is exactly $Qf_s(h)(\Tilde h)(x)$.
Using the same Fourier bound as in the proof of Theorem~\ref{thm: PF1deriv} and the domination by
$(M_2+M_1)\|h\|\,\|\Tilde h\|$, we obtain
\[
\|\,Qf_s(h)(\Tilde h)\,\|_{L^1(\mathcal{X})}
\le
\frac{|\mathcal{X}|}{2\pi}\big(\|M_2\|_{L^1}+\|M_1\|_{L^1}\big)\,\|h\|\,\|\Tilde h\|
+\|g_s''\|_{L^1(\mathcal{X})}\,\|h\|\,\|\Tilde h\|,
\]
since $|\mathbb E[h(Y)]|\le \|h\|$ and $|\mathbb E[\Tilde h(Y)]|\le \|\Tilde h\|$.
Therefore $Qf_s\in \mathcal{B}(\mathcal{S},\mathcal{B}\big(\mathcal{S},\mathcal{Z}_0)\big)$, and the existence of the limit
in $\mathcal{Z}$ with an $L^1(\mathcal{X})$ dominating envelope yields the Fr\'echet differentiability claim.
\end{proof}


\section{Supplementary Material from Section~\ref{sec:SQP}}\label{proof:loc}
\subsection{Example from Section~\ref{sec:local-geometry}}
\begin{example}[Infinite-dimensional counterexample]\label{ex:infinite-counterexample}
In the infinite-dimensional setting, the injectivity of the first variation $Pf_{\bar{s}}(\cdot)$ does not imply that $Pf_{\bar{s}}(\cdot)$ is bounded below. The following counterexample shows that the range of the operator $Pf_{\bar{s}}(\cdot)$ is not closed. Specifically, we can construct a sequence $\{h_n\}$ in $\mathcal{S}$ such that $\|h_n\| = 1$ but $\|Pf_{\bar{s}}(h_n)\|_{\mathcal{Z}_0}\to 0$. Since the operator $Pf_{\bar{s}}(\cdot)$ is injective, it attains $Pf_{\bar{s}}(h) = 0_{\mathcal{Z}_0}$ if and only if $h=0$ $\mu$--a.e., which in turn shows that the range is not closed.

As a reminder of our settings:
\begin{enumerate}
  \item $Y$ is a real random variable with probability measure $\mu$ and density $f$ on a compact set $\mathcal{X}$;
  \item $\bar{s}\in\mathcal{S}$ is the transport map and $Z := \bar{s}(Y)$ has a $C^1$ density $g_{\bar{s}}$ (the pushforward density);
\end{enumerate}
As a special context we also suppose we are considering a target density $g_{\bar{s}}$ where
\begin{align*}
    g_{\bar{s}},g_{\bar{s}}' \in L^1(\mathcal{X}) \cap L^2(\mathcal{X}).
\end{align*}

Recall that the first variation of the PF operator at a feasible $\bar{s}$ is the bounded linear map $Pf_{\bar{s}}:\mathcal{S}\to \mathcal{Z}_0$ given by
\begin{equation}\label{eq:Ps-linearized}
Pf_{\bar{s}}(h)(x)
  = \frac{1}{2\pi}\int_{\mathbb{R}} -e^{itx} e^{-it\,\mathbb{E}[\bar{s}(Y)]}\,
    \phi_h(t)\,dt - g_{\bar{s}}'(x)\mathbb{E}[h(Y)],
\end{equation}
where
\begin{equation}\label{eq:phi-def-counter}
\phi_h(t)
 := it\,\mathbb{E}\big[(h(Y)-\mathbb{E}[h(Y)])\,
      e^{it(\bar{s}(Y)-\mathbb{E}[\bar{s}(Y)])}\big],\qquad t\in\mathbb{R}.
\end{equation}

Next, define for each integer $n$
\[
h_n(Y) := e^{inZ} = e^{in \bar{s}(Y)} = \cos nZ + i\sin nZ.
\]
Then $|h_n(Y)| = 1$ a.s., so
\[
\|h_n\|^2 = \mathbb{E}[|h_n(Y)|^2] = 1,
\qquad\Rightarrow\qquad
\|h_n\| = 1 \quad \forall n.
\]

Let $\varphi(u) := \mathbb{E}[e^{iuZ}]$ be the characteristic function of $Z$.
Since $g_{\bar{s}}\in L^1(\mathcal{X})$, the Riemann--Lebesgue lemma yields
$\varphi(u)\to 0$ as $|u|\to\infty$.
Using the density $g_{\bar{s}}$ of $Z$,
\[
\mathbb{E}[h_n(Y)]
 = \mathbb{E}[e^{inZ}]
 = \int_{\mathcal{X}} e^{inx} g_{\bar{s}}(x)dx
 = \varphi(n)\longrightarrow 0
\qquad\text{as }|n|\to\infty.
\]

Let us first consider the $\phi_{h_n}$ term. By \eqref{eq:phi-def-counter},
\[
\phi_{h_n}(t)
 = it\,\mathbb{E}\big[(h_n(Y)-\mathbb{E}[h_n(Y)])e^{it(Z-\mathbb{E}[\bar{s}(Y)])}\big].
\]
Using the definition of $Z$ and pulling out the constant term $e^{-it\mathbb{E}[\bar{s}(Y)]}$, we get
\begin{equation}\label{eq:phi-hn-char}
\phi_{h_n}(t)
 = it\,e^{-it\mathbb{E}[\bar{s}(Y)]}\Big(\varphi(n+t) - \varphi(n)\varphi(t)\Big),
\end{equation}
where $\varphi(u) = \mathbb{E}[e^{iuZ}]$ is the characteristic function.
Fix $t\in\mathbb{R}$. As $|n|\to\infty$, the Riemann--Lebesgue lemma gives
$\varphi(n+t)\to 0$ and $\varphi(n)\to 0$, while $\varphi(t)$ is fixed, so
\[
\varphi(n+t) - \varphi(n)\varphi(t) \longrightarrow 0.
\]
Thus
\begin{equation}\label{eq:phi-hn-ptwise}
\phi_{h_n}(t) \longrightarrow 0
\quad\text{for each fixed }t\in\mathbb{R}.
\end{equation}

We now show that $\phi_{h_n}\to 0$ in $L^2(\mathbb{R})$. From
\eqref{eq:phi-hn-char}, we have
\[
|\phi_{h_n}(t)|
 \le |t|\big(|\varphi(n+t)| + |\varphi(n)|\,|\varphi(t)|\big).
\]
Since $g_{\bar{s}},g_{\bar{s}}' \in L^2(\mathcal{X})$, by Plancherel’s Theorem \citep[Chapter 6, Theorem 2.4]{stein2003fourier} and the identity that
$\widehat{g_{\bar{s}}'}(u) = iu\widehat{g_{\bar{s}}}(u)$, we have
\[
\varphi = \widehat{g_{\bar{s}}} \in L^2(\mathbb{R}),
\qquad
t\varphi(t) \in L^2(\mathbb{R}).
\]
Using $|\varphi(n)|\le \|g_{\bar{s}}\|_{L^1}$ and the fact that translations are isometries
in $L^2$, we obtain a uniform upper bound
\[
|\phi_{h_n}(t)|
 \le C\Big(|t||\varphi(t+n)| + |t||\varphi(t)|\Big)
 \le (1+|t|)\,w(t),
\]
for some $w\in L^2(\mathbb{R})$ independent of $n$ and some constant $C>0$.
Combining this uniform bound with the pointwise convergence
\eqref{eq:phi-hn-ptwise} and the dominated convergence theorem in $L^2$ yields
\begin{equation}\label{eq:phi-hn-L2}
\|\phi_{h_n}\|_{L^2(\mathbb{R})} \longrightarrow 0
\qquad\text{as }|n|\to\infty.
\end{equation}

Now define
\[
\psi_n(t) := e^{-i\mathbb{E}[Z]}\,\phi_{h_n}(t).
\]
Since $|e^{-i\mathbb{E}[Z]}|=1$, we have
$\|\psi_n\|_{L^2_t} = \|\phi_{h_n}\|_{L^2_t}$. Let $\mathcal{F}^{-1}$
be the inverse Fourier transform.
By Plancherel's Theorem there exists $C_1>0$
such that
\[
\|\mathcal{F}^{-1}\psi_n\|_{L^2_x}
 \le C_1\|\psi_n\|_{L^2_t}
 = C_1\|\phi_{h_n}\|_{L^2_t}.
\]
But the first term in \eqref{eq:Ps-linearized} is exactly
$\mathcal{F}^{-1}\psi_n$, up to the fixed constant $-(2\pi)^{-1}$.
Thus
\[
\left\|\frac{1}{2\pi}\int_{\mathbb{R}} -e^{itx}e^{-it\,\mathbb{E}[\bar{s}(Y)]}\,
  \phi_{h_n}(t)\,dt \right\|_{L^2_x}
 \le C_2\|\phi_{h_n}\|_{L^2_t}
\]
for some constant $C_2>0$. By \eqref{eq:phi-hn-L2}, the right-hand side tends
to zero, so the Fourier part vanishes in $L^2_x$. Since $\mathcal{X}$ is compact, convergence in $L^2(\mathcal{X})$ implies convergence in $L^1(\mathcal{X})$, and therefore in $\mathcal{Z}_0$.

We now address the second term, $g_{\bar{s}}'(x)\mathbb{E}[h_n(Y)]$. Since $\mathbb{E}[h_n(Y)] = \varphi(n)\longrightarrow 0$,
we have
\[
\|g_{\bar{s}}'(\cdot)\,\mathbb{E}[h_n(Y)]\|_{\mathcal{Z}_0}
 = |\varphi(n)|\,\|g_{\bar{s}}'\|_{L^1(\mathcal{X})} \longrightarrow 0.
\]

Combining the two pieces we obtain
\[
\|Pf_{\bar{s}}(h_n)\|_{\mathcal{Z}_0} \longrightarrow 0,
\]
while $\|h_n\|=1$ for all $n$.
Consequently, $Pf_{\bar{s}}:\mathcal{S}\to \mathcal{Z}_0$ is
not bounded below:
\[
\inf_{\substack{h\in\mathcal{S} \\ \|h\|=1}} \|Pf_{\bar{s}}(h)\|_{\mathcal{Z}_0} = 0.
\]
\end{example}

\section{Additional Experiments}
\refstepcounter{experiment}\label{exp:g-to-g}
\subsection*{Experiment~\ref{exp:g-to-g}: Gaussian to Gaussian}
\paragraph{\it Linear Transport Map}
In our first experiment, we consider the source and target distributions to be both Gaussian. Specifically, the source distribution is taken to be $\mu =\mathcal{N}(-0.4, 1.2)$ and the target distribution is $\mathcal{\nu} = N(1.6, 0.55)$. The transport map is parameterized as a linear function $s(x) = \alpha x + \beta$ with $\alpha > 0$ enforced through an exponential reparameterization $\alpha = e^\gamma$. For numerical implementation, we restrict to a domain that covers at least $95\%$ of the probability mass of both the source and target Gaussians. This domain is then discretized into a uniform grid of evaluation points. On this grid, the pushforward condition is enforced by matching the transported density to the target density.

This experiment provides a controlled and interpretable benchmark: the true optimal affine map is known, the numerical constraints are transparent, and the finite–grid discretization makes the SQP procedure easy to evaluate step by step.

\paragraph{\it Quadratic Transport Map}
We also conducted the same experiment while parameterizing the transport map as a quadratic function of the form
\begin{align*}
s(x) = \alpha x + \beta + \gamma x^2.
\end{align*}
The problem setup, cost function, and discretization remain unchanged. In this case, the additional quadratic term provides extra flexibility, but the SQP solver consistently drives the coefficient $\gamma$ toward zero. This behavior is consistent with the theoretical fact that the optimal map between two one–dimensional Gaussian distributions under quadratic cost is affine. Hence, the experiment not only converges successfully but also validates that the method correctly recovers the linear structure of the true solution even when a richer parameterization is allowed.

\begin{figure}[ht]
    \centering
    \begin{subfigure}{0.48\textwidth}
        \centering
        \includegraphics[width=\linewidth]{graphs/mu0=-1.4, s0=2.2, mu1=0.6, s1=1.25.png}
        \caption{Linear map.}
        \label{fig:sub-linear}
    \end{subfigure}\hfill
    \begin{subfigure}{0.48\textwidth}
        \centering
        \includegraphics[width=\linewidth]{graphs/q mu0=-1.4, s0=2.2, mu1=0.6, s1=1.25.png}
        \caption{Quadratic map.}
        \label{fig:sub-quadratic}
    \end{subfigure}
    \caption{Comparison of Target Distribution: (a) linear transport map and (b) quadratic transport map.}
    \label{fig:linear-vs-quadratic}
\end{figure}
\begin{figure}[ht]
    \centering
    \begin{subfigure}{0.40\textwidth}
        \centering
        \includegraphics[width=\linewidth]{graphs/mu0=-1.4, s0=2.2, mu1=0.6, s1=1.25 benchmark.png}
        \caption{Linear map.}
        \label{fig:sub-linear}
    \end{subfigure}\hfill
    \begin{subfigure}{0.40\textwidth}
        \centering
        \includegraphics[width=\linewidth]{graphs/q mu0=-1.4, s0=2.2, mu1=0.6, s1=1.25 benchmark.png}
        \caption{Quadratic map.}
        \label{fig:sub-quadratic}
    \end{subfigure}
    \caption{Comparison of the Solved Transport Map: (a) linear transport map and (b) quadratic transport map.}
    \label{fig:linear-vs-quadratic}
\end{figure}


\refstepcounter{experiment}\label{exp:u-to-b}
\subsection*{Experiment~\ref{exp:u-to-b}: Uniform to Beta}
In this experiment, we transport from a Uniform$[0,1]$ with source density $f(x)=1$ to a Beta$(1,3)$ with target density $g(y)=3(1-y)^2$. We use a piecewise linear approximation of the transport map $s(x)$ on the interval $[0,1]$. As a reminder, we evaluate the constraint on a finite set of grid points: $\{d_j\}_{j=1}^{m_d} \subset \mathrm{supp}(g)$ where $m_d$ is the number of grid points. To improve resolution in regions where the target density $g$ varies rapidly,
we employ an adaptive discretization for the constraints based on quantile locations rather than
uniform grid points. Ideally, we want to evaluate the constraints more frequently at the place where the target probability mass is concentrated. Specifically, let $F_X$ and $F_Y$
denote the CDFs of the Uniform$[0,1]$ source and Beta$(1,3)$ target, we enforce the measure preserving constraints
at equal probability grids in the range of $F_Y$. We choose a set of uniformly spaced probability levels
\begin{align*}
v_j = p_{\mathrm{lo}} +
\frac{j-1}{m_d-1}\bigl(p_{\mathrm{hi}}-p_{\mathrm{lo}}\bigr),
\qquad j=1,\ldots,m_d,    
\end{align*}
with $0<p_{\mathrm{lo}}\ll 1$ and $p_{\mathrm{hi}}\approx 1$
to avoid numerical extremes near $y=0$ and $y=1$. 
The corresponding constraint grid on the target domain is then defined by
\begin{align*}
d_j := F_Y^{-1}(v_j).
\end{align*}
Because the source distribution is Uniform$[0,1]$, we maintain a uniform
grid $\{x_i\}_{i=1}^{m_x}\subset \mathrm{supp}(f)$ as the pieces in the piecewise approximation of the transport map, therefore each piece in the piecewise linear approximation is transporting the same amount of source mass. This quantile based construction concentrates constraint points in
regions where $g(y)$ has significant curvature (e.g., near $y=1$
when $(\alpha,\beta)=(1,3)$), ensuring uniform accuracy in probability
space and improved enforcement of the measure preserving condition.
The resulting SQP iterations remain stable and converge rapidly to the
analytically known optimal map $s^*$, validating the correctness and
robustness of our finite--dimensional formulation.


Because the analytic optimal map is known from quantile matching as 
\begin{align*}
s^*(x)=F^{-1}_{\mathrm{Beta}(1,3)}(x),
\end{align*}
this experiment provides a direct numerical benchmark. The learned map aligns closely with $s^*(x)$ across the entire domain, and the pushforward density $P_f(s)(y)$ matches the target $g(y)$ up to numerical tolerance.

\begin{figure}[ht]
    \centering
    \includegraphics[width=0.55\linewidth]{graphs/U2B_params13_24pieces_401_200ite1.png}
    \caption{Source $\mu(x)$ and target densities $\nu(y)$ for the Uniform$\to$Beta experiment, showing close agreement between the learned pushforward $P_f(s)(y)$ and the true target.}
    \label{fig:u2b-densities}
\end{figure}

\begin{figure}[ht]
    \centering
    \includegraphics[width=\linewidth]{graphs/U2B_params13_24pieces_401_200ite2.png}
    \caption{Computed transport map (SQP, piecewise linear) compared with the analytic quantile map for the Uniform$\to$Beta case.}
    \label{fig:u2b-map}
\end{figure}
\end{document} 